\documentclass[11pt]{article}

\usepackage[T1]{fontenc}
\usepackage[utf8]{inputenc}
\usepackage{lmodern}
\usepackage[margin=1in]{geometry}
\usepackage{amsmath,amssymb,amsthm,mathtools}
\usepackage{enumitem}
\usepackage{booktabs}
\usepackage{array}
\usepackage{longtable}
\usepackage{xcolor}
\usepackage{microtype}
\usepackage{etoolbox}
\usepackage{url}
\usepackage{tikz}
\usetikzlibrary{arrows.meta,calc,decorations.pathreplacing}
\usepackage[colorlinks=true,linkcolor=blue!50!black,citecolor=green!35!black,urlcolor=blue!60!black]{hyperref}
\usepackage[capitalize,nameinlink]{cleveref}

% ---------------------------------------------------------------------------
% Theorem environments
% ---------------------------------------------------------------------------
\theoremstyle{plain}
\newtheorem{theorem}{Theorem}[section]
\newtheorem{lemma}[theorem]{Lemma}
\newtheorem{corollary}[theorem]{Corollary}
\newtheorem{proposition}[theorem]{Proposition}
\newtheorem{conjecture}[theorem]{Conjecture}
\theoremstyle{definition}
\newtheorem{definition}[theorem]{Definition}
\newtheorem{remark}[theorem]{Remark}
\newtheorem{example}[theorem]{Example}

% The Lean proof is the proof; the text gives a faithful outline of it.
\newenvironment{outline}{\begin{proof}[Outline of the formal proof]}{\end{proof}}
% Results that are not formalized get complete proofs, labelled as such.
\newenvironment{informalproof}{\begin{proof}[Proof (not formalized)]}{\end{proof}}

% ---------------------------------------------------------------------------
% Lean references
%   \lean{name}           a declaration, relative to the namespace Transpose.Tape
%                         unless it starts with "Transpose." (fully qualified)
%   \formal{file}{names}  the formal counterpart of a numbered result: a file
%                         relative to lean/Transpose/ and a comma-separated list
%                         of declarations (same naming convention as \lean)
%   \notformalized        marks a numbered result that has no formal counterpart
% ---------------------------------------------------------------------------
% Lean names are typeset like paths, with line breaks allowed after "_", "." and "/".
\DeclareUrlCommand\leanname{\def\UrlFont{\ttfamily\upshape}}
\renewcommand{\UrlBreaks}{\do\_\do\.\do\/}
\newcommand{\lean}[1]{\leanname{#1}}
\newcommand{\leanfile}[1]{\leanname{#1}}
\newcommand{\leanlist}[1]{\begingroup\def\leansep{}%
  \renewcommand*{\do}[1]{\leansep\leanname{##1}\def\leansep{, }}\docsvlist{#1}\endgroup}
\newcommand{\formal}[2]{\par\noindent{\small\upshape\raggedright\textsf{Lean:} \leanlist{#2}\unskip
  \nobreak\hfil\penalty50\hskip1em\hbox{}\nobreak\hfill[\leanfile{#1}]\parfillskip=0pt\par}%
  \smallskip\noindent\ignorespaces}
\newcommand{\notformalized}{\par\noindent{\small\upshape\textsf{Not formalized;} a complete proof is given
  below.}\par\smallskip\noindent\ignorespaces}

% ---------------------------------------------------------------------------
% Notation
% ---------------------------------------------------------------------------
\newcommand{\N}{\mathbb{N}}
\newcommand{\Z}{\mathbb{Z}}
\newcommand{\Q}{\mathbb{Q}}
\newcommand{\R}{\mathbb{R}}
\newcommand{\F}{\mathbb{F}}
\newcommand{\T}{\mathsf{T}}
\newcommand{\clog}[1]{\lceil \log_2 #1 \rceil}
\newcommand{\abs}[1]{\lvert #1 \rvert}
\newcommand{\card}[1]{\lvert #1 \rvert}
\newcommand{\Haddr}{\mathcal{H}}
\newcommand{\eone}{e_1}
\newcommand{\etwo}{e_2}
\newcommand{\Mb}{M_{\mathrm{b}}}
\newcommand{\chicross}{\chi_{\mathrm{cross}}}
\newcommand{\src}{\mathrm{src}}
\newcommand{\snk}{\mathrm{snk}}
\newcommand{\eps}{\varepsilon}
\newcommand{\TCC}{\textsc{TCC}}
\newcommand{\dotminus}{\mathbin{\dot{-}}}% truncated subtraction, x \dotminus y = max(x-y, 0)
\DeclareMathOperator{\rank}{rank}
\DeclareMathOperator{\dist}{dist}
\DeclareMathOperator{\poly}{poly}


\title{Matrix Transposition and the Tape Coding Conjecture}
\author{Hans Brende \and Claude Opus 5.5 (Anthropic)}
\date{October 2026}

\begin{document}

\maketitle

\begin{abstract}
Can a deterministic multitape Turing machine transpose an $n\times n$ matrix with $\lceil\log_2n\rceil$-bit
entries in time linear in the input length $L=n^2\lceil\log_2n\rceil$? Transposition is the data-movement
bottleneck of FFT-based integer multiplication on tapes, and by a reduction of Harvey and van der Hoeven a
linear-time multiplier would give a linear-time transposer. The question is open. We study it against a
frozen formal statement in Lean~4, and every theorem in this paper except two explicitly marked barrier
results is machine-checked.

We prove that every transposer needs $\Omega(\log n)$ passes in which all heads move monotonically, whatever it
does with its data (Theorem~P); that the macroscopic work-head motion of every linear-time transposer satisfies
an isoperimetric inequality that forces it to fold back (Theorem~M*); and that both statements survive when the
machine is correct only on a $2^{-O(n\log n)}$ fraction of the inputs. We then isolate the remaining obstacle as
a statement about network coding: the \emph{Tape Coding Conjecture} asserts that complete exchange admits no
causal code on the block networks of linear-length tape diagrams, equivalently that coding beats multicommodity
flow there by a factor $o(\log n)$. We prove that the conjecture implies that no linear-time transposer exists,
give equivalent and weaker sufficient forms, prove that its plain form implies its robust form for several
classes of codes, and show why a proof must use the shape of tapes. Finally we show that XOR-frame recursion, the
only known technique that beats $O(L\log n)$, cannot reach linear time, and we describe what a construction would
need.

\end{abstract}

\begin{quote}\small
\textbf{Contributions.} Hans Brende chose the problem, fixed its formal statement and the verification standard, and
directed the research. Claude Opus 5.5, a large language model developed by Anthropic, developed the mathematics, wrote
the Lean formalization and wrote this paper under Hans Brende's direction.
\end{quote}

\tableofcontents

\section{Introduction}\label{sec:intro}

\subsection{The problem}

A matrix of $d\times d$ entries of $b=\lceil\log_2d\rceil$ bits, written row by row on the input tape of a
deterministic multitape Turing machine, occupies $L=d^2b$ cells. The machine must write the same matrix column
by column on its output tape. Can it do so in time $O(L)$?

On a random-access machine, or on a machine with a two-dimensional work tape, transposition is trivial, and a
Boolean circuit transposes with no gates at all. On one-dimensional tapes the problem is open. Recursive block
exchange gives time $O(L\log d)$. The best upper bound claimed is $O(L(\log d)^{1-\kappa})$ for a very small
$\kappa>0$, in an unrefereed preprint of OpenAI~\cite{OpenAI2026mult} whose transposition primitive mixes data by
XOR in sheared address frames. No superlinear lower bound is known; indeed none is known on machines with two or more work tapes for
off-line multiplication, sorting or the DFT either~\cite[\S1.4]{HvdH2025} (the superlinear lower bounds that
follow from~\cite{PPST1983} concern problems that are hard for nondeterministic linear time~\cite{Grandjean1988}).
The classical lower bounds for transposition hold for restricted models only: one work
tape~\cite{DMS1991,DietzfelbingerMaass1993}, data movement without coding~\cite{Stoss1973}, and reversal-bounded
conservative machines~\cite{Kirchherr1988}.

The question matters because of integer multiplication. Fast multiplication on Turing machines runs through
multidimensional discrete Fourier transforms~\cite{SchonhageStrassen1971,HvdH2021}, and on tapes such transforms
need transpositions to bring the data of each dimension together. Harvey and van der Hoeven proved a converse:
transposition reduces to multiplication, with no overhead for entries of $\lceil\log_2d\rceil$ bits~\cite[Cor.~6.9]{HvdH2025}. (Their own linear-time statement~\cite[Thm.~6.11]{HvdH2025}, in its simplest case, takes entries of $\lceil\log_2m\rceil$ bits for matrices of at most $m$ bits; up to linear time this is the same problem, \cref{prop:entry-width}.) So a linear-time multiplication machine would give a linear-time transposer, and a proof that transposition needs
superlinear time would be the first superlinear lower bound for off-line multiplication on multitape Turing
machines.

We write~(T) for the statement that a linear-time transposer exists. Its formal version,
\lean{Transpose.LinearTimeTransposable}, is fixed in a frozen Lean file (\cref{sec:model}), together with a checker
that accepts a proof of~(T) or of its negation.

\subsection{Results}

We do not decide~(T). We prove what we can about it, formally, and isolate one precise conjecture
whose proof would refute~(T).

\paragraph{Locality, counting and isoperimetry (\cref{sec:locality,sec:blocks,sec:traces}).} If a set of time steps
must output bits of the input that it never reads, those bits must enter through its entry data: the states at its
entry points and the work symbols it inherits (the Cut Lemma, \cref{lem:cut}). Cutting tapes into blocks of a row's
size and time into slots of the same length turns every run of linear length into a bounded-degree network on $O(d)$ slots, the
\emph{block network}, and the Cut Lemma into an isoperimetric inequality on it (\cref{cor:isoperimetry}). For machines
whose head motion depends on the data, inputs with the same block-level trace (a coarse record of head motion) behave like
an oblivious machine; for a linear-time machine some trace is shared by a $2^{-O(db)}$ fraction of the inputs (the popular trace, \cref{thm:popular}), and the
inequalities hold on it (\cref{lem:cut-trace,cor:isoperimetry-trace}).

\paragraph{Passes (\cref{sec:passes}).} \emph{Theorem~P.} An oblivious transposer whose run splits into $r$ passes, in
each of which every head moves monotonically, satisfies $(k+1)^r\ge1+kd/(2+\sigma_M/b)$, where $\sigma_M$ is the number
of bits of a state plus $k$ work symbols (\cref{thm:P}); \emph{Theorem~P$'$} gives $(k+1)^r\ge1+kd/(2+(\sigma_M+\gamma_T)/b)$ for every
machine halting within $T$ steps, where $\gamma_T=(k+4)\clog{(2T+1)}$ is $O(kb)$ when $T\le\poly(L)$
(\cref{thm:Pprime}). So, for $k\ge1$ and $T\le\poly(L)$, $\Omega(\log d)$ passes are necessary, whatever the
machine does with its data, and this remains true if the machine is correct only on a $2^{-O(db)}$ fraction of the inputs
(\cref{thm:P-on,thm:Pprime-on}).

\paragraph{Macroscopic motion (\cref{sec:mstar}).} \emph{Theorem~M*.} Rescale the work-head trajectories of a
linear-time transposer (for adaptive machines, along its popular trace classes) by $d$ in space (in blocks) and in
time (in slots). Along subsequences they converge to
$1$-Lipschitz paths, and every limit satisfies an isoperimetric inequality~$(\star)$: for every partition of macroscopic time into finitely many intervals, grouped into classes of measure at most $\theta$, the heads return to their positions across different classes with total multiplicity at least $1-2\theta$. Consequently the limits are not all monotone, and the heads
together waste at least one unit of range (\cref{thm:mstar,thm:mstar-adaptive,thm:excess}). The theorem holds for
non-uniform machines and for machines correct on a $2^{-o(L)}$ fraction of the inputs (\cref{thm:mstar-robust}), and
it is a statement about walks satisfying a cut inequality, not about machines (\cref{thm:mstar-walk}).

\paragraph{The Tape Coding Conjecture (\cref{sec:codes,sec:conjecture}).} The block-level results above use the machine only
through the flow of information along its block network. A run yields a \emph{causal network code} on the block
network of its time--space diagram that solves \emph{complete exchange}: every row must reach every column
(\cref{lem:class-code,thm:fast-codes}). On these networks almost every source--sink pair is at distance
$\Omega(\log d)$ while the total capacity is $O(L)$, so a solving code beats multicommodity flow by a factor
$\Omega(\log d)$ (\cref{thm:code-gap}). The \emph{Tape Coding Conjecture}
(\cref{conj:TCC}) asserts that such codes do not exist, even if they need to be correct only on a $2^{-O(d\log d)}$
fraction of the inputs; equivalently, coding beats routing on tape networks by $o(\log d)$ (\cref{thm:tcc-forms}). We
prove:
\begin{itemize}[leftmargin=2em]
\item the conjecture implies $\neg$(T); a proof of it yields the frozen checker's answer file directly
(\cref{thm:reduction}). The reduction is one-way: a refutation gives codes, not machines;
\item much less suffices: no codes at co-entropy $Cd$ for every $C$, for infinitely many $d$, on diagrams that write
their output at the end (\cref{thm:coent-linear,thm:io});
\item under the conjecture every machine fails at every large size, and no non-uniform machine with $2^{O(d\log d)}$
states is fast and correct on a $2^{-O(d\log d)}$ fraction of the inputs (\cref{thm:tcc-consequences});
\item every counterexample to the conjecture obeys Theorem~M* (\cref{cor:cex-star});
\item its plain form (codes correct on all inputs), which already refutes oblivious linear-time transposers
(\cref{thm:oblivious}), implies the robust form for linear and nearly linear codes, for affine and low-degree sets,
for codes with few error patterns, and for constant-size hints (\cref{sec:plainrobust}).
\end{itemize}

\paragraph{Barriers (\cref{sec:barriers}).} The same statement on general bounded-degree graphs would imply superlinear
circuit lower bounds for multiplication (\cref{prop:circuit-barrier}), and its bit-level analogue on tape networks
would separate oblivious from non-oblivious machines off-line (\cref{prop:rotation-barrier}). A proof must therefore use
the shape of tapes; and, since the plain form already excludes oblivious linear-time transposers
(\cref{thm:oblivious}), a proof would give a superlinear lower bound for oblivious transposition, which known
techniques do not provide.

\paragraph{XOR-frame recursions and constructions (\cref{sec:frames,sec:construction}).} In an abstract model of the recursion of the
OpenAI preprint (the correspondence is informal, \cref{sec:frames}), every run realizing a gadget makes at least as many frame changes as the gadget has wires (\cref{thm:frame-change}); by a translation from gadgets to runs, which we prove in full but have not formalized, every gadget has cost $s\ge W$ (\cref{prop:gadget-runs}), while linear time would need $s<W$; integer frames cost more. We record the conditions a construction must meet and the one architecture that we could not rule out; its open
cores are index-coding questions on bounded-degree expanders.

\subsection{How to read this paper: the formal proofs}

The proofs are the Lean proofs. Every numbered theorem, lemma, corollary, proposition and definition, except the four propositions marked as not formalized, names its Lean declarations and the file that contains them, in a line of the form
\begin{center}
{\small\textsf{Lean:} \texttt{declaration, \dots}\hfill[\texttt{file}]}
\end{center}
Declarations are in the namespace \texttt{Transpose.Tape} unless written with their full name, and files are relative to
\leanfile{lean/Transpose/}. The text after a statement, labelled \emph{Outline of the formal proof}, is a faithful and
straightforward guide to the formal proof, written for human readers; it is a convenience, not a proof. Where the
English statement simplifies (for instance by omitting harmless edge cases), the Lean statement is authoritative, and
the formal statements are often slightly stronger. The four results that are not formalized (\cref{prop:entry-width,prop:circuit-barrier,prop:rotation-barrier,prop:gadget-runs}) are marked as such and come with complete proofs. Results from the
literature are cited and not reproved.

The library depends only on Lean~4 and Mathlib (both v4.34.1). It contains no \texttt{sorry}, no unsafe or native code
evaluation, and no axioms beyond Lean's standard ones (\texttt{propext}, \texttt{Classical.choice}, \texttt{Quot.sound}).
\Cref{app:verify} explains how to build it and how to check these claims, and \cref{app:index} lists every declaration
cited in the paper.

\subsection{Related work}

\paragraph{Lower bounds on Turing machines.} Superlinear lower bounds are known for one-tape machines, starting with
Hennie~\cite{Hennie1965} and Maass~\cite{Maass1985}; for matrix transposition on one-tape off-line machines they are due
to Dietzfelbinger, Maass and Schnitger~\cite{DMS1991,DietzfelbingerMaass1993}. For multitape machines the strongest
general result is the separation of deterministic and nondeterministic linear time by Paul, Pippenger, Szemer\'edi and
Trotter~\cite{PPST1983}. Restricted models give more: Stoß proved $\Omega(L\log d)$ for permutations by data
movement~\cite{Stoss1973}, Kirchherr $\Omega(\log d)$ reversals for transposition on conservative
machines~\cite{Kirchherr1988}, and Pippenger and Fischer simulated multitape machines obliviously with a logarithmic
overhead~\cite{PippengerFischer1979}. Our Theorem~P is a reversal lower bound that allows arbitrary coding; lower bounds
for streams with few passes over external memory~\cite{GroheSchweikardt2005,GHS2009} need long or index-tagged items and
do not imply it.

\paragraph{External memory and network coding.} Aggarwal and Vitter determined the I/O complexity of sorting and
transposition for algorithms that move indivisible items~\cite{AggarwalVitter1988}. Whether coding can help is the
question behind the network coding conjecture of Li and Li~\cite{LiLi2004}, which would have implied lower bounds for
multiplication circuits~\cite{AFKL2019}, external-memory sorting~\cite{FHLS2019} and two-tape oblivious machines~\cite{AHJKL2006}.
Braverman and He~\cite{BravermanHe2026} and, concurrently, Zhang, Li and Li~\cite{ZhangLiLi2026} refuted it: in
undirected networks coding can beat multicommodity flow, by a factor $(\log n)^{\varepsilon}$ on networks of size $n$ after amplification~\cite{BGS2017}. Our conjecture is a version of the network coding
conjecture restricted to tape networks and to the densest demand, where a gain of order $\log d$ is at stake.

\paragraph{Multiplication.} Multiplication of $n$-bit integers in time $O(n\log n)$ on multitape machines is due to Harvey and van der
Hoeven~\cite{HvdH2021}, after Sch\"onhage and Strassen~\cite{SchonhageStrassen1971} and F\"urer~\cite{Furer2009}. The
OpenAI preprints claim multiplication below $n\log n$ and transposition below $L\log d$ on multitape machines~\cite{OpenAI2026mult} and exact complex linear circuits of size $o(n\log n)$ for the DFT of length $n$ for infinitely many lengths~\cite{OpenAI2026fourier};
we use them only as motivation, and \cref{sec:frames} shows the limits of their transposition technique.

\subsection{Organization}

\Cref{sec:model} fixes the model and the frozen statement. \Cref{sec:locality,sec:blocks} prove the cut lemmas,
\cref{sec:passes} Theorems~P and~P$'$, \cref{sec:traces} the extension to adaptive machines, and \cref{sec:mstar}
Theorem~M*. \Cref{sec:codes,sec:conjecture,sec:plainrobust} develop codes on tape diagrams, the Tape Coding Conjecture,
and the relation between its plain and robust forms. \Cref{sec:barriers} relates (T) to integer multiplication and explains why the conjecture is hard,
\cref{sec:frames} treats XOR-frame recursions, \cref{sec:construction} the construction side, and \cref{sec:open} lists
open problems. The appendices give the index of Lean declarations and instructions for building and verifying the
library.

\section{The model and the frozen statement}\label{sec:model}

This section fixes the machine model and the statement~(T) exactly as they are formalized in the
frozen file \leanfile{Transpose/Statement.lean}, and introduces the vocabulary used in the rest of
the paper: total runs, obliviousness, correctness on a set of inputs, and co-entropy.

\subsection{Machines}

A machine $M$ has a finite state set $Q=\{0,\dots,q\}$ with initial state $0$, a read-only
two-way input tape, $k\ge 0$ read/write work tapes, and a write-only one-way output tape.
\begin{itemize}[leftmargin=2em]
\item The input tape holds the input word $w\in\{0,1\}^*$ in cells $0,\dots,\abs{w}-1$; every
other cell reads as an end marker.
\item Each work tape is two-way infinite over the alphabet $\Gamma=\{0,\dots,g\}$, where $0$ is the
blank.
\item In each step the machine may append at most one bit to the output.
\end{itemize}
The transition function $\delta$ maps the current state, the input symbol under the input head and
the $k$ work symbols under the work heads either to ``halt'' or to a new state, an input-head move,
a symbol to write and a head move for each work tape, and optionally one output bit. Head moves
are $-1$, $0$ or $+1$. All heads start at cell $0$, the work tapes start blank, and the output
starts empty. A configuration records the state, the head positions, the work-tape contents and the
output written so far. The machine \emph{computes $\mathit{out}$ on $w$ within $t$ steps}
(\lean{Transpose.ComputesWithin}) if, for some $s\le t$, the $s$-th configuration exists, is
halting, and has output exactly $\mathit{out}$.

Every step writes the cell under each work head (possibly rewriting the old symbol). The machine
is deterministic, and its transition table is finite; in particular $q$, $g$ and $k$ do not depend
on the input.

\subsection{Matrices and their encoding}

For $d\ge1$ (the side length of the matrix, written $n$ in the Lean files) let
\[
  b(d)=\max\bigl(1,\lceil\log_2 d\rceil\bigr)
  \qquad\text{(formally } \lfloor\log_2(d-1)\rfloor+1\text{ with }\lfloor\log_2 0\rfloor:=0\text{, \lean{Transpose.entryBits})}.
\]
A $d\times d$ matrix $A$ with $b$-bit entries is encoded row-major
(\lean{Transpose.encode}): the $d^2$ entries $(0,0),(0,1),\dots,(0,d-1),(1,0),\dots$ are written
as consecutive fixed-width blocks of $b$ bits, bit $0$ first, with no delimiters. The input length
is $L=d^2b$, and $d$ is determined by $L$. Bit $t$ of entry $(i,j)$ sits at position
$\alpha=(di+j)b+t$, and in the encoding of the transpose $A^{\T}$ the same bit sits at position
\[
  \operatorname{tpos}(\alpha)=(dj+i)b+t .
\]

\subsection{The statement}

\begin{definition}[Statement (T)]\label{def:T}
\formal{Statement.lean}{Transpose.TransposesWithin, Transpose.LinearTimeTransposable}
A machine $M$ \emph{transposes within $T(\cdot)$} if for every $d\ge1$ and every $d\times d$
matrix $A$ with $b(d)$-bit entries, $M$ computes the encoding of $A^{\T}$ on the encoding of $A$
within $T(L)$ steps, where $L=d^2b(d)$. Statement~(T), \lean{Transpose.LinearTimeTransposable},
says that some machine $M$ and some constant $c$ satisfy this with $T(L)=c\,(L+1)$.
\end{definition}

The file \leanfile{Transpose/Statement.lean} depends on Lean core only and is frozen: its SHA-256
hash, together with that of the checker \leanfile{Verify.lean}, is recorded in
\leanfile{FROZEN.sha256}. A resolution of~(T) is, by definition, a file
\leanfile{Transpose/Answer.lean} containing \texttt{theorem Transpose.answer} whose type is exactly
\lean{Transpose.LinearTimeTransposable} or its negation; \leanfile{Verify.lean} re-checks with the kernel, from scratch, every declaration of its environment
that does not ship with the Lean toolchain (\leanfile{Answer.lean} itself and the Mathlib and library modules it imports, directly or indirectly), and allows no axioms beyond
\texttt{propext}, \texttt{Classical.choice} and \texttt{Quot.sound}; in particular it rejects
\texttt{sorry}, \texttt{native\_decide} and new axioms. This paper does \emph{not} resolve~(T), so no \leanfile{Answer.lean} exists and
\leanfile{Verify.lean} does not compile yet; \cref{sec:conjecture} shows the exact form a proof of
the Tape Coding Conjecture would take as such a file.

\subsection{Runs}\label{sec:runs}

Throughout, the semantics of the frozen model is accessed through the \emph{total run}
\lean{run}: $\mathrm{run}(M,w,\tau)$ is the configuration after $\tau$ steps, frozen once the
machine halts. We write $\delta_\tau(w)$ for the transition taken at step $\tau$
(\lean{transAt}; it is ``halt'' exactly when the configuration at time $\tau$ is halting), and
say that $M$ \emph{halts at $T$ on $w$} (\lean{HaltsAt}) if steps $0,\dots,T-1$ execute and the
configuration at time $T$ is halting.

\begin{lemma}[The total run agrees with the frozen semantics]\label{lem:bridge}
\formal{Tape/Statements/Core.lean}{stepN_init_eq_some_iff, computesWithin_iff, haltsAt_unique, outRev_reverse_eq, transposes_of_transposesWithin}
\begin{enumerate}[label=(\alph*)]
\item The $\tau$-step iteration of the frozen step function from the initial configuration
succeeds with $c$ if and only if no step before $\tau$ halts and $\mathrm{run}(M,w,\tau)=c$.
\item $M$ computes $\mathit{out}$ on $w$ within $t$ steps if and only if $M$ halts on $w$ at some
$T\le t$ and the output of $\mathrm{run}(M,w,T)$ is $\mathit{out}$. The halting time is unique.
\item The output after $\tau$ steps is the list of the bits appended at steps $0,\dots,\tau-1$, in
order.
\item If $M$ transposes within any $T(\cdot)$, then for every $d\ge1$ it transposes $d\times d$
matrices with $b(d)$-bit entries in the sense of \cref{def:obl} below.
\end{enumerate}
\end{lemma}

\begin{outline}
By induction on the number of steps, comparing one step of the frozen semantics with one step of
the total run. Part~(d) is a restatement of the definitions.
\end{outline}

\begin{lemma}[Cell contents]\label{lem:cells}
\formal{Tape/Statements/Core.lean}{work_eq_written, work_eq_blank}
If work head $t$ is at cell $p$ at time $\sigma<\tau$ and not at $p$ at any time strictly between
$\sigma$ and $\tau$, then at time $\tau$ cell $p$ holds the symbol written at step $\sigma$ (after halting, the symbol
written is the one already in the cell, \lean{writtenSym}). A cell
not visited by its head before time $\tau$ is blank at time $\tau$.
\end{lemma}

\begin{outline}
Induction on $\tau$: a step changes only the cells under the work heads.
\end{outline}

\subsection{Obliviousness, correctness and co-entropy}\label{sec:obl}

\begin{definition}[Oblivious sets, transposers]\label{def:obl}
\formal{Tape/Defs.lean}{ObliviousOn, matrixInputs, Transposes}
Let $W$ be a set of inputs and $T\in\N$. The machine $M$ is \emph{oblivious on $W$ with halting
time $T$} if every $w\in W$ halts at $T$, and any two inputs of $W$ give the same input-head
position, the same work-head positions and the same decision whether to append an output bit, at
every time $\tau$ (including $\tau\ge T$). We write $\mathrm{Mat}(d,b)$ for the set of encodings of
$d\times d$ matrices with $b$-bit entries, and say that $M$ \emph{transposes $d\times d$ matrices
with $b$-bit entries} if on each of them it computes the encoding of the transpose within some
number of steps (no time bound).
\end{definition}

Since time $T$ is included, obliviousness also fixes the head moves of the last step; every
theorem below that assumes obliviousness assumes this (slightly stronger) form. By
\lean{matrixInputs_eq}, $\mathrm{Mat}(d,b)$ is exactly the set of words of length $d^2b$, so
``oblivious on $\mathrm{Mat}(d,b)$'' means oblivious at input length $L=d^2b$.

Many results hold for machines that are correct, or fast, only on a subset of the inputs. We
measure the size of such a subset by its co-entropy.

\paragraph{Co-entropy.}\label{def:coent} A set $W\subseteq\mathrm{Mat}(d,b)$ has \emph{co-entropy at most $m$} if
$2^{d^2b}\le 2^m\,\abs{W}$, that is, if $W$ contains at least a $2^{-m}$ fraction of all inputs. In the formal
statements this is written out as the hypothesis $2^{d^2b}\le2^m\cdot\abs{W}$ (named \texttt{hco}), with
$\abs{W}$ the cardinality \texttt{W.ncard}.

A machine is \emph{correct on $W$} (at time $T$) if every $w\in W$ encodes some matrix $A$ and the
output of $\mathrm{run}(M,w,T)$ is the encoding of $A^{\T}$. In the formal statements this
hypothesis is written out explicitly (it is the hypothesis \texttt{hcorr} of, for example,
\lean{cut_bound_on}).

\begin{lemma}[Encoding facts]\label{lem:encoding}
\formal{Tape/Statements/Core.lean}{matrixInputs_eq, tpos_spec}
The encodings of $d\times d$ matrices with $b$-bit entries are exactly the words of length $d^2b$.
For $j<d$ and $t<b$, $\operatorname{tpos}((di+j)b+t)=(dj+i)b+t$.
\end{lemma}

\begin{outline}
Induction on the nested concatenation defining the encoding, and integer division.
\end{outline}

\subsection{Non-vacuity}

The hypotheses of the lower bounds below (obliviousness together with correctness) are
satisfiable, and the counting quantities they involve are non-trivial on concrete machines.

\begin{proposition}[Witness machines]\label{prop:witnesses}
\formal{Tape/Statements/Core.lean}{copyTM_obliviousOn, copyTM_transposes, tr2_obliviousOn, tr2_transposes, tr2_cutPairs, tr2k_obliviousOn, tr2k_transposes}
\begin{enumerate}[label=(\alph*)]
\item The copy machine (which copies its input to its output) is oblivious on the inputs of every
length $\ell$, with halting time $\ell$, and transposes $1\times1$ matrices with $b$-bit entries
for every $b$.
\item A $9$-state machine \lean{Witness.tr2} without work tapes is oblivious on $\mathrm{Mat}(2,1)$
with halting time $5$ and transposes $2\times2$ matrices with $1$-bit entries. For the all-zero
input and the step set $J=\{2,3,4\}$ the set of cut pairs (\cref{def:cut}) is $\{(0,1)\}$.
\item The same machine with one idle work tape, \lean{Witness.tr2k} (so $k=1$), is oblivious on
$\mathrm{Mat}(2,1)$ with halting time $5$ and transposes $2\times2$ matrices with $1$-bit entries.
\end{enumerate}
\end{proposition}

\begin{outline}
For the copy machine, by induction on time, for inputs of every length $\ell$: at time $\tau\le\ell$
the input head is at cell $\tau$ and the output is the first $\tau$ input bits, and at time $\ell$ the
machine reads the end marker and halts. For \lean{Witness.tr2} and \lean{Witness.tr2k}, by evaluating
(with \texttt{decide}) the runs on all $16$ inputs of length $4$: the input-head positions $0,1,2,3,3,3$ (and, for \lean{Witness.tr2k}, the work-head position $0$) and the output schedule are common, the machine halts at time $5$, and the output on
$x_1x_2x_3x_4$ is $x_1x_3x_2x_4$. The cut pairs follow from the input-head positions $2,3,3$ at steps $2,3,4$ of the
all-zero run and the output positions $1,2,3$ appended at these steps.
\end{outline}

\section{Locality and counting}\label{sec:locality}

The basic tool of the paper is a cut lemma: if a set $J$ of time steps must produce, in the
output, bits of the input that $J$ never reads, then those bits must enter $J$ through its
\emph{entry data}, the few states and work-tape symbols that $J$ inherits from the rest of the
run. This section states the cut lemma for an arbitrary set of steps; \cref{sec:blocks} bounds the
entry data of unions of time slots, and later sections apply it to passes, trace classes and
networks.

\subsection{Entry data and cut sets}

Fix a machine $M$, an input $w$ and a finite set $J$ of time steps.

\begin{definition}[Entry data]\label{def:entry}
\formal{Tape/Defs.lean}{entryPoints, entryCells, IsLastVisit}
An \emph{entry point} of $J$ is a step $\tau\in J$ with $\tau\ne0$ and $\tau-1\notin J$. An
\emph{entry cell} of $J$ is a pair $(t,\tau)$ of a work tape $t$ and a step $\tau\in J$ such that
the cell under head $t$ at time $\tau$ was last visited by that head at a step $\sigma\notin J$
(so $\sigma<\tau$). Cells never visited before $\tau$ are blank for every input and are not entry
cells. We write
\[
  \eone(J)=\#\{\text{entry points of }J\},\qquad \etwo(J)=\#\{\text{entry cells of }J\}.
\]
\end{definition}

Step $0$ is never an entry point: the state at time $0$ is the initial state for every input.

\begin{definition}[Cut set, sources, sinks, cut pairs]\label{def:cut}
\formal{Tape/Defs.lean}{cutSet, srcRows, snkCols, cutPairs, tpos}
Let $w_0$ be an input of length $n^2b$. Relative to the run of $M$ on $w_0$:
\begin{itemize}[leftmargin=2em]
\item the \emph{cut set} of $J$ is the set of input positions $d<n^2b$ that the input head visits
at no step of $J$ but whose transposed copy, output bit $\operatorname{tpos}(d)$, is appended at a
step of $J$;
\item the \emph{source rows} $\src(J)$ are the rows $i<n$ such that the input head lies in the
row's interval $[inb,(i+1)nb)$ at some step of $J$;
\item the \emph{sink columns} $\snk(J)$ are the columns $j<n$ all of whose output bits
$[jnb,(j+1)nb)$ are appended at steps of $J$;
\item the \emph{cut pairs} are $S_J=([n]\setminus\src(J))\times\snk(J)$.
\end{itemize}
\end{definition}

Every bit of an entry $(i,j)$ with $(i,j)\in S_J$ lies in the cut set, so
$b\,\abs{S_J}\le\abs{\mathrm{cut}(J)}$ (\lean{cutPairs_card_mul_le}).

\subsection{The Locality Lemma}

\begin{lemma}[Locality]\label{lem:locality}
\formal{Tape/Statements/Core.lean}{locality, locality_oblivious, locality_function}
Let $w,w'$ be inputs whose runs have the same work-head positions at all times $\tau<T$, and let
$J\subseteq[0,T)$. Suppose the two runs agree on
\begin{enumerate}[label=(\roman*)]
\item the state at every entry point of $J$,
\item the symbol under the work head at every entry cell of $J$ (entry cells taken in the run on
$w$), and
\item the input symbol read at every step of $J$.
\end{enumerate}
Then at every step $\tau\in J$ the two runs have the same state, the same symbols under the work
heads, and take the same transition. In particular, if $M$ is oblivious on a set $W$ with halting
time $T$, then for every $J\subseteq[0,T)$ there is a function $F$ such that for every $w\in W$
the transitions taken during $J$ are $F$ applied to the states of $w$ at the entry points, the
symbols of $w$ at the entry cells, and the input symbols read by $w$ during $J$.
\end{lemma}

The run need not be live: either run may halt during $J$, since halted runs are frozen.

\begin{outline}
Induction over the steps of $J$ in increasing order. At $\tau\in J$, the state is given by~(i)
if $\tau$ is an entry point, and otherwise is produced by the common transition at $\tau-1\in J$.
The symbol under work head $t$ is given by~(ii) if $(t,\tau)$ is an entry cell; is blank in both
runs if the cell was never visited (\cref{lem:cells}); and otherwise was written at its last visit
$\sigma\in J$, where by induction both runs took the same transition at the same head position.
With~(iii), the transition function receives identical arguments. The functional form follows
because, on an oblivious set, entry points, entry cells and input positions are the same for all
inputs.
\end{outline}

The same induction gives the forms used for machines whose heads depend on the data. In those
forms the head positions are not assumed equal; they are derived.

\begin{lemma}[Locality, data-dependent heads]\label{lem:locality-core}
\formal{Tape/Statements/Core.lean}{locality_core, locality_interval}
Let $J$ be a set of steps and $A$ a set of \emph{anchors} such that every non-anchor step of $J$
has its predecessor in $J$. Suppose that two runs agree, at every anchor in $J$, on the state, all
head positions and the symbols under the work heads; that the input symbols agree wherever the
input heads agree; and that whenever the work heads agree at a non-anchor step $\tau\in J$, each
cell they read has, at some time $a\le\tau$, the same content in both runs, with the last visits
to that cell during $[a,\tau)$ in both runs lying in $J$. Then the runs agree at every step of $J$
on the state, all head positions, the symbols under the work heads and the transition taken. In
particular, two runs that agree at time $a$ on the state and all head positions, on the contents
at time $a$ of every work cell visited during $[a,b)$, and on the input bits read during $[a,b)$,
agree on all of this at every time in $[a,b]$ (transitions for $\tau<b$).
\end{lemma}

\begin{outline}
The induction of \cref{lem:locality}, carrying the head positions along: at a non-anchor step the
positions agree because the previous transitions agree. The interval form is the case
$J=[a,b)$, $A=\{a\}$.
\end{outline}

\subsection{The Cut Lemma}

\begin{lemma}[Cut Lemma]\label{lem:cut}
\formal{Tape/Statements/Core.lean}{cut_bound, cut_bound_pairs, cut_bound_bits, card_le_entryData}
Let $M$ be oblivious on $\mathrm{Mat}(n,b)$ with halting time $T$ and transpose $n\times n$
matrices with $b$-bit entries. Let $w_0\in\mathrm{Mat}(n,b)$ and $J\subseteq[0,T)$. Then
\[
  2^{\abs{\mathrm{cut}(J)}}\;\le\;(q+1)^{\eone(J)}\,(g+1)^{\etwo(J)},
  \qquad\text{hence}\qquad
  2^{b\abs{S_J}}\le(q+1)^{\eone(J)}(g+1)^{\etwo(J)}
\]
and $b\abs{S_J}\le\eone(J)\clog{(q+1)}+\etwo(J)\clog{(g+1)}$.
\end{lemma}

The lemma has a form for machines that are oblivious and correct only on a subset of the inputs.
Nothing is assumed about inputs outside the subset.

\begin{lemma}[Cut Lemma on a set]\label{lem:cut-on}
\formal{Tape/Statements/Core.lean}{cut_bound_on, cut_bound_on_coentropy}
Let $M$ be oblivious on a set $W$ with halting time $T$ and correct on $W$ at time $T$, let
$w_0\in W$ and $J\subseteq[0,T)$. Then
\[
  2^{\abs{\mathrm{cut}(J)}}\,\abs{W}\;\le\;2^{n^2b}\,(q+1)^{\eone(J)}(g+1)^{\etwo(J)} .
\]
In particular, if $W$ has co-entropy at most $m$, then
$2^{\abs{\mathrm{cut}(J)}}\le 2^m(q+1)^{\eone(J)}(g+1)^{\etwo(J)}$.
\end{lemma}

\begin{outline}
Group the inputs of $W$ by their bits outside the cut set: there are at most
$2^{n^2b-\abs{\mathrm{cut}(J)}}$ groups. Within a group the inputs agree on every input bit read
during $J$, because the cut set is read at no step of $J$ and all inputs of $W$ share the input
trajectory. By \cref{lem:locality} the transitions during $J$, and hence the output bits appended
during $J$, are a function of the entry data. These output bits contain the transposed copy of
every cut-set bit, and by correctness they equal it; so distinct inputs of a group have distinct
entry data. Hence each group has at most $(q+1)^{\eone}(g+1)^{\etwo}$ members. The formal proof
packages this injection as an abstract counting core, shared with \cref{lem:cut-trace}; the case
$W=\mathrm{Mat}(n,b)$ gives \cref{lem:cut}.
\end{outline}

\subsection{Families of step sets}

Applying the Cut Lemma to several step sets and multiplying gives the form used for partitions.

\begin{corollary}[Cut Lemma for families]\label{cor:cut-family}
\formal{Tape/Statements/Core.lean}{cut_bound_sum, cut_bound_sum_bits, cut_bound_sum_on, sum_cutPairs_ge, uncovered_cols_le, uncovered_cols_le_len}
Under the hypotheses of \cref{lem:cut}, for every finite family $(J_u)_{u\in I}$ of subsets of
$[0,T)$ (not necessarily disjoint),
\[
  b\sum_{u}\abs{S_{J_u}}\;\le\;\clog{(q+1)}\sum_u\eone(J_u)+\clog{(g+1)}\sum_u\etwo(J_u),
\]
and under the hypotheses of \cref{lem:cut-on} with co-entropy $m$ the right side of the
exponential form gains a factor $2^{\abs{I}m}$. Moreover:
\begin{enumerate}[label=(\alph*)]
\item if every $J_u$ reads at most $\rho$ rows, then
$\bigl|\bigcup_u\snk(J_u)\bigr|\,(n-\rho)\le\sum_u\abs{S_{J_u}}$;
\item if the $J_u$ are the label classes of a labelling $\lambda$ of the steps $<T$ and $b\ge1$,
then the number of columns that are not sinks of a single class is at most the number of label
changes $\#\{1\le\tau<T:\lambda(\tau)\ne\lambda(\tau-1)\}$. This holds for any run that has
appended at least $n^2b$ output bits by time $T$, with no obliviousness, correctness or halting
assumption.
\end{enumerate}
\end{corollary}

\begin{outline}
The first part multiplies the bounds of \cref{lem:cut} and \cref{lem:cut-on}. For~(a), a sink
column of $J_u$ pairs with each of the at least $n-\rho$ rows not read by $J_u$. For~(b), the
output bits of a column are appended during an interval of time; if they are not all appended in
one label class, the label changes inside that interval, and the intervals of distinct columns
are disjoint.
\end{outline}

\section{Blocks, slots and the block network}\label{sec:blocks}

To apply the Cut Lemma to long stretches of a run we coarsen space and time. Cut each work tape
into \emph{blocks} of $r$ consecutive cells and time into \emph{slots} of $r$ consecutive steps.
In one slot a head moves at most $r$ cells, so it touches at most two blocks; the entry data of a
union of slots is then carried by a few blocks per slot boundary and by blocks that the head
leaves and later revisits. The result is a bounded-degree network on the slots, the
\emph{block network}, and the Cut Lemma becomes an isoperimetric inequality on it. In the
applications $r=\lceil nb/G\rceil$ with $G=\max(1,\clog{(g+1)})$, so that a block holds about one
matrix row.

\subsection{Blocks and residency}

Let $h:\N\to\Z$ be a head trajectory with unit steps (every head of the frozen model has unit
steps, \lean{unitSteps_inTraj}, \lean{unitSteps_workTraj}) and let $r\ge1$.

\begin{definition}[Blocks, slots, residency]\label{def:residency}
\formal{Tape/Defs.lean}{blockOf, curBlock, slotOf, numSlots, residentPair, pairRes, touchRes, AdmissibleRes, IsChord}
The block of a cell $p$ is $\lfloor p/r\rfloor$, and $\mathrm{cur}_\tau$ is the block of the head
at time $\tau$. Slot $y$ is the set of steps $[yr,(y+1)r)\cap[0,T)$; there are
$N=\lceil T/r\rceil$ slots. At time $\tau$ the \emph{resident pair} consists of $\mathrm{cur}_\tau$
and the block the head occupied just before its last block change at or before $\tau$ (if there
was one). The \emph{pair residency} $R_y$ is the union of the resident pairs over the steps of
slot $y$; the \emph{touched residency} is the set of blocks the head occupies during slot $y$. A
family $R_y$ of sets of blocks is an \emph{admissible residency with constant $c$} if the head's
block at every step $\tau<T$ is resident at the slot of $\tau$, and no slot has more than $c$
resident blocks. A \emph{chord} of $R$ for block $\beta$ from slot $x$ to slot $z$ is a pair
$x+1<z$ with $\beta\in R_x$, $\beta\in R_z$ and $\beta\notin R_y$ for $x<y<z$: the block stops
being resident after slot $x$ and becomes resident again at slot $z$.
\end{definition}

\begin{lemma}[Residency]\label{lem:residency}
\formal{Tape/Statements/Core.lean}{pairRes_subset, pairRes_admissible, touchRes_admissible, slot_block_step, exists_slot_block_eq, pairChange_spacing}
Let $h$ have unit steps and $r\ge1$.
\begin{itemize}[leftmargin=3em]
\item[(V1)] $R_y\subseteq\{\mathrm{cur}_{yr}-1,\mathrm{cur}_{yr},\mathrm{cur}_{yr}+1\}$. Pair residency
is admissible with $c=3$, and touched residency is admissible with $c=2$ and also lies within
$\pm1$ of $\mathrm{cur}_{yr}$.
\item[(V2)] $\abs{\mathrm{cur}_{(y+1)r}-\mathrm{cur}_{yr}}\le1$; consequently, for slots
$y_1\le y_2$, every block between $\mathrm{cur}_{y_1r}$ and $\mathrm{cur}_{y_2r}$ is the head's
block at the first step of some slot in $[y_1,y_2]$.
\item[(S)] Two changes of the resident pair are at least $r$ steps apart.
\end{itemize}
\end{lemma}

\begin{outline}
In $r$ steps a unit-step walk moves at most $r$ cells, hence crosses at most one block
boundary; this gives (V2) and its crossing form (a discrete intermediate value argument). For
(V1), induct over the steps of the slot: the head's block stays within $\pm1$ of
$\mathrm{cur}_{yr}$, and the other block of a resident pair is adjacent to the block the head
entered at its last block change. For (S), a change of the resident pair means that the head
enters a block outside the current pair, which requires crossing a whole block.
\end{outline}

\begin{lemma}[Uniqueness of visit ends and chord endpoints (P2)]\label{lem:P2}
\formal{Tape/Statements/Core.lean}{pairRes_begin_unique, pairRes_end_unique, pairRes_chord_source_unique, pairRes_chord_target_unique, touchRes_begin_unique, touchRes_end_unique, touchRes_chord_source_unique, touchRes_chord_target_unique}
For pair residency and for touched residency, with $h$ of unit steps and $r\ge1$:
at most one block becomes resident at each slot $y+1$; at most one block stops being resident
after each slot $y$ with $(y+1)r<T$; each slot is the source of at most one chord, and the target
of at most one chord.
\end{lemma}

\begin{outline}
By (V1), (V2) and the spacing of pair changes: two blocks entering at the same slot boundary would
require two pair changes within $r$ steps, or a jump of two blocks.
\end{outline}

\subsection{Charging and the Slot Cut Bound}

For a set $U$ of slots let $J=J(U)$ be the steps lying in slots of $U$. A \emph{memory boundary}
of $U$ is a slot $z\ge1$ in $U$ with $z-1\notin U$; we write $\mu(U)$ for their number. For an
admissible residency $R$ let $\chi(U)$ be the number of pairs $(z,\beta)$ such that some chord of
$R$ for $\beta$ enters $U$ at $z$, from a slot outside $U$ (\lean{chordsInto}), summed over the
work tapes.

\begin{lemma}[Charging]\label{lem:charging}
\formal{Tape/Statements/Core.lean}{entryTimes_charged, entryTimes_card_le, entryCells_slotSteps_card, entryPoints_slotSteps_card, charging_bound, srcRows_slotSteps_card_le}
Let $r\ge1$, let each work tape $t$ have an admissible residency $R_t$ with constant $c$, and let
$U$ be a set of slots. Then:
\begin{enumerate}[label=(\alph*)]
\item every external read of $J(U)$ on a tape, at a step $\tau$ on block $\beta=\mathrm{cur}_\tau$,
is charged to a slot $z\le\tau/r$ that comes after the cell's last write, and either $z$ is a
memory boundary of $U$ with $\beta\in R_z\cap R_{z-1}$, or a chord for $\beta$ enters $U$ at $z$;
\item the entry points of $J(U)$ are exactly the first steps of its memory boundaries:
$\eone(J(U))=\mu(U)$;
\item $\etwo(J(U))\le r\,\chi(U)+c\,k\,r\,\mu(U)$;
\item[(P1)] if $r\le nb$, every slot reads at most two input rows, so $J(U)$ reads at most
$2\abs{U}$ rows.
\end{enumerate}
\end{lemma}

\begin{outline}
For~(a), the slot of the cell's last write lies outside $U$ and $\beta$ is resident there. Let
$y^*$ be the last slot before $\tau$'s slot that lies outside $U$ and at which $\beta$ is
resident, and $z$ the first slot after $y^*$ at which $\beta$ is resident again; then $z\in U$,
and either $z=y^*+1$ is a memory boundary with $\beta\in R_z\cap R_{z-1}$, or a chord for $\beta$
enters $U$ at $z$. For~(c), charge each external read to its event $(z,\beta)$ together with the
offset of its cell in the block. Two reads with the same charge read the same cell; the later
read's last write is then no earlier than the earlier read, which does not precede slot $z$,
contradicting~(a). So the charge is injective, and there are at most $r(\chi(U)+c\,\mu(U))$ external reads per tape. Part~(b) is a
direct computation, and (P1) holds because a slot reads at most $r\le nb$ consecutive input
cells.
\end{outline}

\begin{corollary}[Slot Cut Bound]\label{cor:slot-cut}
\formal{Tape/Statements/Core.lean}{slot_cut_bound, slot_cut_bound_partition, uncovered_cols_le_slots, uncovered_cols_le_slots_len}
Let $M$ be oblivious on $\mathrm{Mat}(n,b)$ with halting time $T$ and transpose, let $r\ge1$, and
let $G'=\clog{(g+1)}$. For every set $U$ of slots, with pair residency,
\[
  b\,\abs{S_{J(U)}}\;\le\;\mu(U)\clog{(q+1)}+r\,G'\bigl(\chi(U)+3k\,\mu(U)\bigr).
\]
For a labelling $\lambda$ of the slots with $\Mb$ label changes, any admissible residency with
constant $c$, and $\chicross$ the number of chords whose two endpoints carry different labels
(summed over tapes),
\[
  b\sum_{\ell}\abs{S_{J_\ell}}\;\le\;\Mb\bigl(\clog{(q+1)}+c\,k\,r\,G'\bigr)+r\,G'\,\chicross,
\]
where $J_\ell$ is the set of steps in slots labelled $\ell$; and if $b\ge1$, at most $\Mb$
columns are not sinks of a single class (for any run that has appended $n^2b$ bits by time $T$).
\end{corollary}

\begin{outline}
\Cref{lem:cut} applied to $J(U)$ with \cref{lem:charging}(b),(c); for a labelling, sum over the
classes, observe that memory boundaries of the classes are label changes and chords entering a
class from another class are crossing chords, and apply \cref{cor:cut-family}(b).
\end{outline}

\subsection{Discrete isoperimetry}

Combining the Slot Cut Bound with (P1) gives the inequality behind all results on head motion.
Read it as follows: if time is cut into label classes of at most $m$ slots each, then about
$n^2b$ bits must cross between classes, and the only channels are the label changes (each worth
$O(k)$ blocks plus a state) and the chords between classes (one block each).

\begin{corollary}[Discrete Isoperimetry]\label{cor:isoperimetry}
\formal{Tape/Statements/Core.lean}{isoperimetry_res, isoperimetry, isoperimetry_touch, isoperimetry_res_on}
Let $M$ be oblivious on $\mathrm{Mat}(n,b)$ with halting time $T$ and transpose, let
$1\le r\le nb$, let each work tape have an admissible residency with constant $c$, and let
$\lambda$ be a labelling of the slots with $\Mb$ label changes and at most $m$ slots per label.
Then
\[
  2^{(n-\Mb)(n-2m)b}\;\le\;(q+1)^{\Mb}\,(g+1)^{r\,(c\,k\,\Mb+\chicross)} .
\]
Pair residency gives $c=3$ and touched residency $c=2$. If $M$ is oblivious and correct only on
a set $W$ of co-entropy at most $m_c$, the right side gains the factor $2^{(\Mb+1)m_c}$: each
label class pays $m_c$ bits.
\end{corollary}

\begin{outline}
By (P1) each class reads at most $2m$ rows, so by \cref{cor:cut-family}(a),(b) the cut pairs of
the classes number at least $(n-\Mb)(n-2m)$; the Slot Cut Bound bounds them from above. For the
robust form, apply \cref{lem:cut-on} to each class.
\end{outline}

When $m=o(n)$ and $\Mb=o(n)$, with $r\approx nb/G$, the inequality forces
$\chicross=\Omega(n)$: the work heads must carry information between distant times through
revisited blocks. \Cref{sec:mstar} turns this into a constraint on the macroscopic motion of the
heads.

\section{Passes}\label{sec:passes}

A \emph{pass} is a stretch of time during which every head moves monotonically, each in its own
direction. Transposition by recursive block exchange uses $O(\log n)$ passes (folklore), and the
external-memory intuition is that fewer passes cannot suffice. This section proves this for every
machine of the model whose running time is polynomial in $L$, whatever it does with its data. For
oblivious machines $r$ passes force $(k+1)^r\ge1+kn/(2+\sigma_M/b)$ (Theorem~P); for all machines
halting within $T$ steps they force $(k+1)^r\ge1+kn/(2+(\sigma_M+\gamma_T)/b)$ with
$\gamma_T=(k+4)\clog{(2T+1)}$, which is $O(kb)$ for $T\le\poly(L)$ (Theorem~P$'$). In both cases
$r\ge\log_{k+1}n-O(1)$ for $k\ge1$. Throughout,
\[
  \sigma_M=\clog{(q+1)}+k\,\clog{(g+1)}
\]
is the number of bits needed to write down a state and $k$ work symbols.

\subsection{Pass partitions and addresses}

\begin{definition}[Passes]\label{def:passes}
\formal{Tape/DefsP.lean}{MonoOn, AllMono, PassPartition, greedyBound, greedyPasses, addrCount}
A trajectory is \emph{monotone} on $[a,e)$ if all its steps in $[a,e)$ are non-decreasing, or all
are non-increasing. A \emph{pass partition} of the run of $M$ on $w$ into $r$ passes is a sequence
$0=p_0\le p_1\le\dots\le p_r=T$ such that on each $[p_i,p_{i+1})$ the input head and every work
head are monotone (empty passes are allowed). The \emph{greedy} partition extends each pass as far
as possible. We write
\[
  \Haddr=\Haddr(k,r)=\sum_{i<r}(k+1)^i ,\qquad\text{so that}\quad k\Haddr+1=(k+1)^r .
\]
\end{definition}

\begin{lemma}[Greedy passes]\label{lem:greedy}
\formal{Tape/Statements/Core.lean}{greedy_passPartition, greedyPasses_le, passPartition_iff, passPartition_indep, greedyPasses_indep, addrCount_mul}
The greedy passes form a pass partition, and a pass partition into $r$ passes exists if and only
if the greedy count is at most $r$. If $M$ is oblivious on $W$, the pass partitions of the runs on
the inputs of $W$ are the same for all inputs, and so is the greedy count. Moreover $k\Haddr+1=(k+1)^r$.
\end{lemma}

\begin{outline}
Minimality: the greedy partition stays ahead. By induction, the $i$-th greedy boundary is at least
the $i$-th boundary of any pass partition, since monotonicity on an interval is inherited by
subintervals. Conversely, a greedy partition with fewer passes is padded with empty passes. Independence
holds because pass partitions depend only on the head trajectories.
\end{outline}

\subsection{Where accessed cells come from}

\begin{lemma}[Fresh cells]\label{lem:fresh}
\formal{Tape/Statements/Core.lean}{pass_fresh}
Let $h$ be monotone on $[s,e)$ and $s\le\rho<a\le\tau\le e$. If $h(\tau)\ne h(a)$ then
$h(\rho)\ne h(\tau)$: during a monotone stretch, a cell occupied after time $a$, other than the
cell occupied at time $a$, was not occupied earlier in the stretch.
\end{lemma}

\begin{outline}
Monotonicity: $h(\rho)\le h(a)\le h(\tau)$ (or the reverse) with $h(a)\ne h(\tau)$.
\end{outline}

So, inside a pass, a work head that at time $\tau\ge a$ reads a cell other than the one it occupied at
time $a$ reads a cell last visited during $[a,\tau)$, a cell last visited in an \emph{earlier pass}, or a
blank cell. This is what makes the dependency structure of a pass schedule a forest of bounded size.

\subsection{The dependency forest}

Fix a run with a pass partition $p$ into $r$ passes and an output segment $[o_1,o_2)$.

\begin{definition}[Dependency forest]\label{def:forest}
\formal{Tape/DefsP.lean}{rootNode, sSet, childNode, forestNode, addrsBelow, addresses, forest, readSet, combSet}
\begin{itemize}[leftmargin=2em]
\item The \emph{root node} in pass $P_0$ is the set of steps of pass $P_0$ during which the output
length moves inside $[o_1,o_2)$.
\item For a node $J$, an earlier pass $P'$ and a work tape $t$, the \emph{S-set} consists of the
cells that head $t$ occupies during $J$, other than its cell at the first step of $J$, whose last
visit before that first step lies in pass $P'$. The \emph{child} of $J$ at $(P',t)$ is the set of
steps of pass $P'$ at which head $t$ lies between the least and the greatest cell of the S-set
(empty if the S-set is empty).
\item An \emph{address} is a pass $P_0<r$ followed by a list $(P_1,t_1),\dots,(P_j,t_j)$ with
$P_0>P_1>\dots>P_j$ and tapes $t_i$; the node at an address is obtained from the root in pass
$P_0$ by taking successive children. There are at most $\Haddr$ addresses
(\lean{card_addresses_le}). The \emph{forest} of the segment is the union of all its nodes.
\item $\mathrm{read}(J)$ is the set of input positions $<L$ read during $J$, and the \emph{comb}
of the segment is the set of input positions $d<n^2b$ with $\operatorname{tpos}(d)\in[o_1,o_2)$.
\end{itemize}
\end{definition}

Nodes are intervals of time, because heads are monotone within a pass and the output length is
monotone.

\begin{lemma}[Forest]\label{lem:forest}
\formal{Tape/Statements/Core.lean}{forest_subset, forest_entryPoints_subset, forest_entryCells_subset, forest_entryPoints_card_le, forest_entryCells_card_le, card_addresses_le, forest_info}
For every run with a pass partition into $r$ passes and every output segment, the forest consists
of steps $<T$; every entry point of the forest is the first step of a node; and every entry cell
of the forest is a read, at the first step of a node, of the cell under a work head. Hence
$\eone(\mathrm{forest})\le\Haddr$ and $\etwo(\mathrm{forest})\le k\Haddr$. Consequently, if $M$ is
oblivious on $\mathrm{Mat}(n,b)$ with halting time $T$ and transposes, then for the common run
\[
  2^{\abs{\mathrm{comb}\setminus\mathrm{read}(\mathrm{forest})}}\;\le\;
  \bigl((q+1)(g+1)^k\bigr)^{\#\mathrm{addresses}} .
\]
\end{lemma}

\begin{outline}
Entry points of a union of intervals are first steps of the intervals. For entry cells, take a
read by head $t$ at a step $\tau$ of a node $J$ that is not its first step $a$, of a cell whose last visit
before $\tau$ is at $\sigma$. If $\sigma\ge a$ then $\sigma\in J$, since $J$ is an interval. Otherwise the
cell is not the one occupied at time $a$ (that visit would come after $\sigma$), so by \cref{lem:fresh} it
was not visited earlier in the current pass: $\sigma$ lies in an earlier pass $P'$, the cell belongs to the
S-set at $(P',t)$, and $\sigma$ is a step of the child at $(P',t)$, which is part of the forest. So the read
is not external. For the last inequality, every output bit of the segment is
appended in some root node, so the comb bits that the forest does not read are in its cut set,
and the Cut Lemma on a set (\cref{lem:cut-on}) applies to the forest.
\end{outline}

\subsection{Disjointness and coverage}

\begin{lemma}[Disjointness]\label{lem:disjoint}
\formal{Tape/Statements/Core.lean}{forestNode_disjoint, read_chain}
At a fixed address, the nodes of two output segments, the first ending before the second starts,
are disjoint. If finitely many (at least one) segments have pairwise ordered output ranges (of any two,
one ends before the other starts), the input positions read by their nodes at a fixed address satisfy
$\sum_s\abs{\mathrm{read}(J_s)}+1\le L+\#\{s\}$.
\end{lemma}

\begin{outline}
By induction on the depth of the address, the nodes of an earlier and a later segment are
\emph{time-separated}, one lying entirely before the other: for the roots this is the monotonicity
of the output length. If two nodes of one pass are time-separated, then, the head being monotone in
that pass, every cell of the earlier node lies weakly on one side of the cell $c$ under the head at the
later node's start, while the S-set of the later node, which excludes $c$, lies strictly on the other side;
so the hulls of the two S-sets are disjoint, and monotonicity in the earlier pass makes the children
time-separated. For the second part, all nodes at a fixed address lie in one pass, in which the input head
is monotone; ordered in time, the read set of each node meets the union of the read sets of the earlier
nodes in at most one position (the input cell at its first step), so
$\sum_s\abs{\mathrm{read}(J_s)}\le\abs{\bigcup_s\mathrm{read}(J_s)}+\#\{s\}-1\le L+\#\{s\}-1$.
\end{outline}

Group the columns into $\lceil n/\gamma\rceil$ groups of $\gamma$ consecutive columns (the last may be
smaller); group $s$ owns the output segment of its columns, and its comb consists of at most $\gamma b$
consecutive input bits in each of the $n$ rows.

\begin{lemma}[Coverage]\label{lem:coverage}
\formal{Tape/Statements/Core.lean}{coverage_per_address, coverage_grouped}
For every run with a pass partition into $r$ passes and every $\gamma\ge1$, at a fixed address the comb
bits read by the nodes of the $\lceil n/\gamma\rceil$ groups number at most $(\gamma+1)nb$ in total.
Over all addresses, the comb bits read by the group forests number at most
$\#\mathrm{addresses}\cdot(\gamma+1)nb$.
\end{lemma}

\begin{outline}
At a fixed address the read sets of the groups form a time-ordered chain of intervals of input
positions (\cref{lem:disjoint}). An exact residue count bounds what such a chain can collect from
the combs: with $m=nb$, the positions $d$ of a set $R$ of consecutive positions with $d\bmod m$ in
a set $\mathrm{Res}$ number at most $\abs{\mathrm{Res}}(\abs{R}+m)/m$; summing over the chain with
the true group widths gives $(\gamma+1)nb$.
\end{outline}

\subsection{Theorem P}

\begin{theorem}[Theorem P$_\gamma$, grouped columns]\label{thm:Pgamma}
\formal{Tape/Statements/Core.lean}{theoremP_grouped, theoremP_grouped_bits}
Let $M$ be oblivious on $\mathrm{Mat}(n,b)$ with halting time $T$, transpose $n\times n$ matrices
with $b$-bit entries, and let the common run have a pass partition into $r$ passes. Then for every
$\gamma\ge1$
\[
  2^{n^2b}\le 2^{\Haddr(\gamma+1)nb}\bigl((q+1)(g+1)^k\bigr)^{\lceil n/\gamma\rceil\Haddr},
  \qquad
  n^2b\le\Haddr\bigl((\gamma+1)nb+\lceil n/\gamma\rceil\,\sigma_M\bigr).
\]
\end{theorem}

\begin{theorem}[Theorem P]\label{thm:P}
\formal{Tape/Statements/Core.lean}{theoremP, theoremP_bits, theoremP_bits_add_one, theoremP_div, theoremP_pow, theoremP_greedy}
Under the hypotheses of \cref{thm:Pgamma},
\[
  2^{n^2b}\le 2^{2nb\Haddr}\bigl((q+1)(g+1)^k\bigr)^{n\Haddr},\qquad
  nb\le\Haddr\,(2b+\sigma_M),\qquad
  k\,nb\le\bigl((k+1)^r-1\bigr)(2b+\sigma_M),
\]
also in the weaker form $n^2b\le\Haddr(n(2b+\sigma_M)+1)$, and the same holds with $r$ the greedy pass
count (the greedy passes form a pass partition,
\cref{lem:greedy}; \lean{theoremP_greedy} states the bits form). In particular, for $k\ge1$ and
$b=\lceil\log_2n\rceil$, a fixed machine needs $r\ge\log_{k+1}n-O(1)$ passes.
\end{theorem}

\begin{outline}
Theorem~P is the case $\gamma=1$ of Theorem~P$_\gamma$. For each group $s$, \cref{lem:forest}
bounds the comb bits of $s$ not read by the group's forest:
$2^{\abs{\mathrm{comb}_s\setminus U_s}}\le((q+1)(g+1)^k)^{\Haddr}$. \Cref{lem:coverage} bounds
the comb bits that are read: $\sum_s\abs{\mathrm{comb}_s\cap U_s}\le\Haddr(\gamma+1)nb$. The combs
partition the $n^2b$ input positions, so
$n^2b\le\sum_s\abs{\mathrm{comb}_s\setminus U_s}+\Haddr(\gamma+1)nb$. The per-group bounds are summed,
$\sum_s2^{\abs{\mathrm{comb}_s\setminus U_s}}\le\lceil n/\gamma\rceil((q+1)(g+1)^k)^{\Haddr}$, and the
inequality of arithmetic and geometric means (the final step shared with Theorem~P$'$) gives
$2^{\sum_s\abs{\mathrm{comb}_s\setminus U_s}}\le((q+1)(g+1)^k)^{\lceil n/\gamma\rceil\Haddr}$, hence the
claim. The other forms are arithmetic, using $k\Haddr+1=(k+1)^r$.
\end{outline}

Theorem~P allows arbitrary coding: the machine may combine bits in its finite control and on its
tapes in any way, dwell, and move its heads at arbitrary relative speeds. The closest previous lower
bound for transposition is Kirchherr's $\Omega(\log n)$ reversals for \emph{conservative} machines, which
only move bits~\cite{Kirchherr1988}; the stream lower bounds~\cite{GroheSchweikardt2005,GHS2009} need long or
index-tagged items (\cref{sec:intro}).

\subsection{Machines correct on part of the input}

\begin{theorem}[Theorem P on a set]\label{thm:P-on}
\formal{Tape/Statements/Core.lean}{theoremP_on, theoremP_on_bits, theoremP_on_div, theoremP_on_pow, theoremP_grouped_on, theoremP_grouped_on_bits, theoremP_on_robust_gamma, theoremP_on_robust}
Let $M$ be oblivious on a set $W$ with halting time $T$ and correct on $W$, let the common run
have a pass partition into $r$ passes, and let $W$ have co-entropy at most $m$. Then for every
$\gamma\ge1$
\[
  n^2b\le\lceil n/\gamma\rceil\,m+\Haddr\bigl((\gamma+1)nb+\lceil n/\gamma\rceil\,\sigma_M\bigr):
\]
each group of $\gamma$ columns pays $m$ bits. In particular:
\begin{enumerate}[label=(\alph*)]
\item at co-entropy $Cn$ ($C\in\N$): $nb\le Cn+\Haddr(2b+\sigma_M)$, so $\Haddr\ge n(b-C)/(2b+\sigma_M)$, and
$k\,nb\le kCn+((k+1)^r-1)(2b+\sigma_M)$;
\item at co-entropy $Cnb$ ($C\in\N$), for every $\gamma\ge1$:
$(\gamma-C)\,n^2b\le\gamma C\,nb+\Haddr\bigl(\gamma(\gamma+1)nb+(n+\gamma)\sigma_M\bigr)$; with
$\gamma=2C$, $Cnb\le 2C^2b+\Haddr\bigl(2C(2C+1)b+2\sigma_M\bigr)$, so
$\Haddr\gtrsim n/(2(2C+1))$ and, for $k\ge1$, still $r\ge\log_{k+1}n-O(\log C)$.
\end{enumerate}
\end{theorem}

\begin{outline}
As for Theorem~P$_\gamma$, with the co-entropy term of the Cut Lemma on a set (\cref{lem:cut-on}) kept in
each group's bound: each group's forest pays the co-entropy $m$ once. In the formal development this robust
form is the only direct proof; Theorem~P$_\gamma$ is its case $W=\mathrm{Mat}(n,b)$, $m=0$. The
corollaries are arithmetic.
\end{outline}

\subsection{Data-dependent heads: Theorem P\texorpdfstring{$'$}{'}}

For a machine whose head motion depends on the input, every input has its own run, pass partition
and dependency forests. A forest is then described by a \emph{structure code} per node: its first
and last steps, the positions of all heads and the output length at its first step
(\lean{nodeCode}); and by a \emph{side code}: the state and the symbols under the work heads at its
first step (\lean{sideCode}). The structure codes of nodes inside $[0,T)$ take at most
$(2T+1)^{k+4}$ values (\lean{structCount} is this bound; the count is part of the proof of
\cref{lem:forest-agree}).

\begin{lemma}[Forests of different inputs]\label{lem:forest-agree}
\formal{Tape/Statements/Adaptive.lean}{locality_anchored, forest_agree, forest_comb_agree, kraft_group}
Let $w,w'$ be inputs of the same length $L$, with pass partitions of $[0,T)$ and of $[0,T')$ into $r$
passes (each its own), and suppose their forests of a segment have equal structure and side codes at every
address, and that $w$ and $w'$ agree on every input position $<L$ read by both forests. Then the two
forests have the same nodes, and the runs agree on them (state, positions, output length, symbols under the
heads, transitions). If $M$ transposes $n\times n$ matrices with $b$-bit entries, $L=n^2b$, and $w,w'$ halt
at $T$ and $T'$ respectively, they agree on the segment's comb. Consequently, if $M$ transposes and every
$v\in\{0,1\}^{n^2b}$ halts at some $T_v\le T$ with a pass partition of $[0,T_v)$ into $r$ passes, then for
every segment
\[
  \sum_{v\in\{0,1\}^{n^2b}}2^{\abs{\mathrm{comb}\setminus\mathrm{read}_v(\mathrm{forest}_v)}}
  \;\le\;2^{n^2b}\Bigl((2T+1)^{k+4}(q+1)(g+1)^k\Bigr)^{\Haddr}.
\]
\end{lemma}

\begin{outline}
Nodes are intervals and the structure code records their first and last steps, so equal codes give equal
nodes at every address: the two forests are the same set of steps. The first part is then an anchored form
of the Locality Lemma (\lean{locality_anchored}, a corollary of \cref{lem:locality-core}) with the first
steps of the nodes as anchors: the structure and side codes give agreement at the anchors, and
\cref{lem:forest}, applied to each run, shows that every other step of the forest has its predecessor in
the forest and that every cell read there was last written inside the forest (or never); the input symbols
agree at commonly read positions because the inputs have the same length. The comb then agrees by
correctness. For
the Kraft-type inequality, fix the codes: within a class of inputs with equal codes, overwriting
the unread comb bits of $v$ by arbitrary values is injective, because two overwritten inputs that
coincide agree off the union of their free positions, hence on every position read by both
forests, hence (by the first part) on the comb, hence everywhere.
\end{outline}

\begin{theorem}[Theorem P$'$]\label{thm:Pprime}
\formal{Tape/Statements/Adaptive.lean}{theoremPprime_grouped, theoremPprime_grouped_bits, theoremPprime, theoremPprime_bits, theoremPprime_div, theoremPprime_pow, theoremPprime_greedy, theoremPprime_gamma, theoremPprime_ratio}
Let $M$ transpose $n\times n$ matrices with $b$-bit entries, and suppose every matrix input halts
within $T$ steps with a pass partition of its own run into $r$ passes. Let
$\gamma_T=(k+4)\clog{(2T+1)}$. Then for every $\gamma\ge1$
\[
  n^2b\le\Haddr\bigl((\gamma+1)nb+\lceil n/\gamma\rceil(\sigma_M+\gamma_T)\bigr),
  \qquad\text{in particular}\qquad nb\le\Haddr\,(2b+\sigma_M+\gamma_T),
\]
and $k\,nb\le((k+1)^r-1)(2b+\sigma_M+\gamma_T)$; these hold also in log-free form and with greedy pass
counts. The bound also holds with
$\gamma'=\clog r+2(k+2)\clog{(2T+3)}\ge\gamma_T$ in place of $\gamma_T$, in the forms
$n^2b\le\Haddr(n(2b+\sigma_M+\gamma')+1)$ and, for $b\ge1$,
$\Haddr\ge n/(2+(\sigma_M+\gamma'+1)/b)$.
\end{theorem}

\begin{outline}
By \cref{lem:coverage}, which holds run by run,
$\sum_s\abs{\mathrm{comb}_s\setminus U_s(v)}\ge n^2b-\Haddr(\gamma+1)nb$ for every input $v$. Sum
the inequality of \cref{lem:forest-agree} over the groups and average over the $2^{n^2b}$ inputs:
some $v_0$ has $\sum_s2^{m_s(v_0)}\le\lceil n/\gamma\rceil K$, where $m_s(v)$ is the number of
unread comb bits of group $s$ and $K=((2T+1)^{k+4}(q+1)(g+1)^k)^{\Haddr}$. The inequality of
arithmetic and geometric means gives $2^{\sum_sm_s(v_0)}\le K^{\lceil n/\gamma\rceil}$. The
remaining forms are arithmetic.
\end{outline}

If $T\le\poly(L)$ and $b=\lceil\log_2n\rceil$, then $\gamma_T=O(kb)$ and Theorem~P$'$ again gives
$r\ge\log_{k+1}n-O(1)$; this covers every machine allowed by~(T). Theorem~P$'$ also has the robust
form of \cref{thm:P-on}.

\begin{theorem}[Theorem P$'$ on a set, and against the frozen statement]\label{thm:Pprime-on}
\formal{Tape/Statements/Adaptive.lean}{forest_comb_agree_on, kraft_group_on, theoremPprime_grouped_on, theoremPprime_grouped_on_bits, theoremPprime_on, theoremPprime_on_bits, theoremPprime_on_div, theoremPprime_on_pow, theoremPprime_on_robust, theoremPprime_on_robust_gamma, theoremPprime_on_greedy, theoremPprime_linearTime}
Suppose that on every input of a set $W$ of co-entropy at most $m$, the machine $M$ has, at some
time $T'\le T$, the transposed output together with a pass partition of $[0,T')$ into $r$ passes
(it need not halt; other inputs are unconstrained). Then
$n^2b\le\lceil n/\gamma\rceil m+\Haddr((\gamma+1)nb+\lceil n/\gamma\rceil(\sigma_M+\gamma_T))$ for
every $\gamma\ge1$, and the inequalities of \cref{thm:P-on}(a),(b) hold with $\sigma_M+\gamma_T$ in place
of $\sigma_M$. Their asymptotic consequences change: $T\ge L$ forces $\gamma_T\ge2(k+4)\log_2n$, which is
not $o(b)$; if $\gamma_T\le c'b$ (for instance $T\le\poly(L)$), (b) gives
$\Haddr\ge(Cn-2C^2)b/(2C(2C+1)b+2(\sigma_M+\gamma_T))$ and, for $k\ge1$, $r\ge\log_{k+1}n-O(\log(C+c'))$. If $M$ transposes within $c(L+1)$ steps (\cref{def:T}) and every run on an
$n\times n$ input has greedy pass count at most $r$, then
$n^2b(n)\le\Haddr\,n\bigl(2b(n)+\sigma_M+(k+4)\clog{(2c(L+1)+1)}\bigr)$.
\end{theorem}

\begin{outline}
As for Theorem~P$'$, with the Kraft sum over $W$: the injectivity step uses correctness only on
the two inputs it compares, and the averaging step uses $\abs{W}\ge2^{n^2b-m}$, which charges $m$
once per group. The last statement instantiates $T=c(L+1)$.
\end{outline}

\subsection{Non-vacuity and non-triviality}

\begin{proposition}\label{prop:P-witness}
\formal{Tape/Statements/Core.lean}{theoremP_hyps_tr2, theoremP_hyps_tr2k}
The hypotheses of Theorem~P are satisfiable, with $k=0$ (\lean{Witness.tr2}, $n=2$, $b=1$) and
with $k=1$ (\lean{Witness.tr2k}).
\end{proposition}

\begin{outline}
Evaluation of the concrete machines on all inputs; the greedy passes give the pass partition
(\cref{lem:greedy}).
\end{outline}

\begin{proposition}\label{prop:Pprime-witness}
\formal{Tape/Statements/Adaptive.lean}{theoremPprime_hyps_of_oblivious, theoremPprime_hyps_tr2dd, theoremPprime_no_five_pass}
\begin{enumerate}[label=(\alph*)]
\item Every machine that is oblivious on $\mathrm{Mat}(n,b)$ with halting time $T$ and whose common run
has a pass partition into $r$ passes satisfies the run hypothesis of Theorem~P$'$ (every matrix input halts
within $T$ steps with a pass partition into $r$ passes). In particular the machines of Theorem~P satisfy all
hypotheses of Theorem~P$'$, with the same $T$ and $r$.
\item A machine \lean{Witness.tr2dd} transposes $2\times2$ matrices with $1$-bit entries, is
oblivious for no halting time (its running time is $5$ or $6$ depending on the input), and every
input halts within $6$ steps in a single pass.
\item No machine with one work tape, at most $2^8$ states and at most $4$ work symbols transposes
$1024\times1024$ matrices with $10$-bit entries so that every input halts within $2^{30}$ steps in
$5$ passes.
\end{enumerate}
\end{proposition}

\begin{outline}
(a) by \cref{lem:greedy}; (b) by evaluation; (c) by Theorem~P$'$ and arithmetic.
\end{outline}

\section{Adaptive machines: trace classes and the popular trace}\label{sec:traces}

The cut lemmas of \cref{sec:locality,sec:blocks} assume obliviousness: all inputs share one head
trajectory, so the entry data of a step set is a well-defined finite object. A general machine
moves its heads depending on the data. The way around this is to group the inputs by their
\emph{block-level trace}, a coarse record of the head motion that costs only $O(db)=o(L)$ bits when the running time is $O(L)$.
Some trace is shared by a $2^{-O(db)}$ fraction of all inputs (the \emph{popular trace}), and on
that set the machine behaves like an oblivious one at the level of blocks. The cut lemmas then hold
with an additive loss of $O(db)$ bits per label class.

\subsection{Block-level traces}

Let $r\ge1$ be the block size.

\begin{definition}[Block-level trace, trace class]\label{def:trace}
\formal{Tape/DefsTrace.lean}{blockTrace, SameTrace, TraceClass}
The \emph{block-level trace} of the run of $M$ on $w$ at slot $y$ consists of the pair residency of
every work tape at slot $y$, the set of input rows touched during slot $y$, and the output length
at the first step $yr$ of the slot. Two inputs have the \emph{same trace} (for a parameter $T$) if
their traces agree at all slots $y\le\lceil T/r\rceil$. A \emph{trace class} with representative
$w_0$ is a set $F$ of matrix inputs containing $w_0$, all of which halt after exactly $T$ steps and
have the trace of $w_0$.
\end{definition}

\begin{lemma}[Trace classes generalize oblivious sets]\label{lem:trace-obl}
\formal{Tape/Statements/Adaptive.lean}{traceClass_of_oblivious, traceClass_of_obliviousOn, traceClass_of_mem}
If $M$ is oblivious on $\mathrm{Mat}(d,b)$ with halting time $T$, all matrix inputs form one trace
class, of size $2^{d^2b}$. More generally, every nonempty set of matrix inputs on which $M$ is oblivious (with halting time $T$) is a trace class, with any of its members as representative; and any member of a trace class is a
representative.
\end{lemma}

\begin{outline}
The trace is a function of the head trajectories and the output schedule.
\end{outline}

\subsection{Block locality}

For a trace class the entry data of a union of slots must record, at each memory boundary, more than
the offsets of the heads within their blocks: two inputs with the same trace have the same resident blocks,
but a head need not stand in the same one of them in both runs (\cref{prop:offsets}). We record the full
head positions.

\begin{lemma}[Block Locality]\label{lem:block-locality}
\formal{Tape/Statements/Adaptive.lean}{block_locality}
Let $r\ge1$ and let $w,w'$ be inputs of length $d^2b$ with the same trace. Let $U$ be a set of
slots, and suppose that
\begin{enumerate}[label=(\roman*)]
\item at the first step $zr$ of every memory boundary $z$ of $U$ the two runs agree on the state,
the positions of all heads, and the contents of every cell of every block resident at $z$;
\item for every chord of the pair residency of $w$ entering $U$ at a slot $z$ for a block $\beta$,
the two runs agree at time $zr$ on the contents of the cells of $\beta$; and
\item $w$ and $w'$ agree on every input bit in a row touched during $U$.
\end{enumerate}
Then at every step of the slots of $U$ the two runs agree on the state, all head positions, the
output length, the symbols under the work heads and the transition taken.
\end{lemma}

No halting time and no correctness are assumed.

\begin{outline}
This is the entry form of the Locality Lemma (\lean{locality_entry}, a corollary of
\cref{lem:locality-core}), with step $0$ and the first steps of the memory boundaries as anchors.
For a read of a cell in block $\beta$ at a step $\tau$ of the slots of $U$, the charging argument
of \cref{lem:charging} provides the time $a$: either $a=0$, when $\beta$ is resident at no slot
outside $U$ before $\tau$'s slot, or $a=zr$, where $z$ is the first slot after the last such slot
at which $\beta$ is resident, a memory boundary or the target of a chord entering $U$. Every visit
to the cell during $[a,\tau)$, in either run, lies in $U$, because the head's block is always
resident and the residencies of the two runs coincide.
\end{outline}

Recording only the offsets of the work heads within their blocks does not suffice:

\begin{proposition}[Offsets are not enough]\label{prop:offsets}
\formal{Tape/Statements/Adaptive.lean}{block_locality_offsets_counterexample}
There are a machine $M$ that transposes $1\times1$ matrices with $1$-bit entries and two inputs
$w,w'$, both halting at time $9$ with the same trace for $r=4$, such that slot $2$ is a memory
boundary of $U=\{2\}$, the runs agree at time $8$ on the state, the input-head position, the
work-head offsets modulo $4$ and the contents of the resident blocks, no chord enters slot $2$, and
no input row is touched during slot $2$; yet the work heads are in different blocks at time $8$
and the transitions at step $8$ differ.
\end{proposition}

\begin{outline}
An explicit machine, evaluated on both inputs.
\end{outline}

\subsection{The Cut Lemma on trace classes}

Write $G'=\clog{(g+1)}$, and let $D$ be such that $2^{d^2b}\le 2^D\abs{F}$ (the co-entropy of the
class).

\begin{lemma}[Cut Lemma on a trace class]\label{lem:cut-trace}
\formal{Tape/Statements/Adaptive.lean}{cut_bound_trace, cut_bound_trace_bits, cut_bound_trace_on, cut_bound_trace_bits_on}
Let $r\ge1$, let $F$ be a trace class with representative $w_0$, and let $M$ be correct on $F$ (for
instance, let $M$ transpose). For every set $U$ of slots, with $\mu=\mu(U)$ and $\chi=\chi(U)$ taken
for the pair residency of the representative,
\[
  2^{b\abs{S_{J(U)}}}\,\abs{F}\;\le\;2^{d^2b}\Bigl((q+1)(2T+1)^{k+1}\Bigr)^{\mu}(g+1)^{r(3k\mu+\chi)},
\]
\[
  b\,\abs{S_{J(U)}}\;\le\;D+\mu\bigl(\clog{(q+1)}+(k+1)\clog{(2T+1)}\bigr)+r\,G'\,(3k\mu+\chi).
\]
\end{lemma}

\begin{outline}
Inject the class into pairs (entry data, input bits outside a row-based cut set) through the counting core
of \cref{lem:cut-on}: since members read different positions but the same rows, the cut set is replaced by
the bits of the rows not touched during $J(U)$ whose transposed copies are appended during $J(U)$; it
contains the $b\abs{S_{J(U)}}$ bits of the cut pairs. The entry data consist of the state and the head
positions (which lie in $[-T,T]$) at each memory boundary, together with the contents of the resident blocks
there and of the blocks entering along chords. The rows touched and the output positions appended during
$U$ are class invariants, since the rows of each slot and the output lengths at slot starts are part of the
trace and all members halt at $T$. Injectivity follows from \cref{lem:block-locality} and correctness on
$F$.
\end{outline}

\begin{corollary}[Discrete Isoperimetry on a trace class]\label{cor:isoperimetry-trace}
\formal{Tape/Statements/Adaptive.lean}{isoperimetry_trace, isoperimetry_trace_on}
Under the hypotheses of \cref{lem:cut-trace}, with $1\le r\le db$ and a labelling of the slots with
$\Mb$ label changes and at most $m$ slots per label,
\[
  (d\dotminus\Mb)(d\dotminus2m)\,b\;\le\;(\Mb+1)D+\Mb\bigl(\clog{(q+1)}+(k+1)\clog{(2T+1)}\bigr)
  +r\,G'\,(3k\Mb+\chicross),
\]
where $x\dotminus y=\max(x-y,0)$ and $\chicross$ is taken for the representative's pair residency.
\end{corollary}

\begin{outline}
Sum the bits form of \cref{lem:cut-trace} over the at most $\Mb+1$ label classes and bound the left
side as in \cref{cor:isoperimetry}.
\end{outline}

\subsection{The popular trace}

\begin{definition}[Slot length and horizon]\label{def:slotlen}
\formal{Tape/DefsM.lean}{cellBits, slotLen, horizon}
For a machine $M$ let $G=\max(1,\clog{(g+1)})$ be the number of bits per work cell. The slot
length at size $d$ is $r=\lceil db(d)/G\rceil$, so that a block of $r$ cells holds about one matrix row.
For a time constant $c$ the \emph{horizon} is $c_0=cG+1$; if $T\le cL$ there are at most $c_0d$
slots (\lean{numSlots_le_horizon}).
\end{definition}

\begin{theorem}[Popular trace]\label{thm:popular}
\formal{Tape/Statements/Adaptive.lean}{exists_popular_traceClass, traceBits_le}
Let $1\le r\le db$ and suppose every matrix input halts within $T_{\max}$ steps. Then for some
$T\le T_{\max}$ there is a trace class $F$ with $2^{d^2b}\le\mathcal{N}\cdot\abs{F}$, where
$\mathcal{N}=\mathcal{N}(k,d,\lceil T_{\max}/r\rceil,T_{\max})$ is an explicit bound on the number
of pairs (halting time, trace) (\lean{traceCount}). If $T_{\max}\le cL$ and $r$ is the slot length,
then $\clog{\mathcal{N}}\le C\,d\,b(d)$ for an explicit constant $C$ depending only on $M$ and $c$
(\lean{traceConst}). So the popular class has co-entropy $D=O(db)=o(L)$.
\end{theorem}

\begin{outline}
Pigeonhole over the pairs (halting time, trace): block indices lie in $[-N,N]$, each residency has
at most three blocks, each slot touches at most two rows, and output lengths are at most
$T_{\max}$. The bound on $\clog{\mathcal{N}}$ is arithmetic from $T\le cd^2b$ and $b\le d\le2^b$.
\end{outline}

The trace used here records absolute block indices, rows and output lengths per slot, which costs
$O(d\log d)$ bits for running time $O(L)$. For machines that write their whole output at the end, a coarser
trace that records only block and row moves and residency patterns, $O(1)$ bits per slot, suffices, and the popular class then has co-entropy $O(d)$; \cref{sec:conjecture} uses this to show
that the coding conjecture is needed only at co-entropy $O(d)$.

\section{Theorem M*: the macroscopic motion of the work heads}\label{sec:mstar}

Discrete Isoperimetry (\cref{cor:isoperimetry,cor:isoperimetry-trace}) says that, for every
partition of time into small classes with few boundaries, $\Omega(n)$ revisited blocks must link
different classes. This section passes to the limit $n\to\infty$. Rescaled, the head trajectories of
a linear-time transposer converge (along subsequences) to $1$-Lipschitz paths, and the counting
inequality becomes an integral inequality~$(\star)$ for these paths: in every partition of
macroscopic time into pieces of measure at most $\theta$, the heads must return to the same
positions in different pieces with total multiplicity at least $1-2\theta$. In particular the
heads cannot all sweep monotonically: some head makes a macroscopic U-turn, and together the heads
waste at least one unit of range.

The argument uses the cut inequality only, not the machine; we state it for arbitrary walks
(Theorem~M* for walks) and obtain the machine forms (oblivious, trace classes, adaptive, robust,
non-uniform) and, in \cref{sec:codes}, the form for network codes as corollaries.

\subsection{Rescaled walks and their limits}

Fix $G\ge1$. At size $n$ the block size is $r=\lceil nb(n)/G\rceil$ (\lean{TSDiagram.blockSize}), so a
block holds a row. For a walk $h:\N\to\Z$ the \emph{rescaled walk} is
\[
  P^{(n)}(u)=\frac{1}{n}\,\Bigl\lfloor \frac{h(\lfloor un\rfloor r)}{r}\Bigr\rfloor ,
\]
the block at the first step of slot $\lfloor un\rfloor$, in units of $n$ blocks
(\lean{blockTraj}). For a machine, $P^{(n)}_t$ is the rescaled trajectory of work head $t$ on the
all-zero input (\lean{rescaledTraj}), or on an input $w$ (\lean{rescaledTrajOf}). If $T\le cL$,
macroscopic time runs over $[0,c_0]$ with $c_0=cG+1$.

\begin{lemma}[Compactness and regularity]\label{lem:walk-basics}
\formal{Tape/Statements/Coding.lean}{blockTraj_limit_lipschitz, exists_subseq_blockTraj, numSlots_le_of, frozen_V2, pairRes_congr}
Let $n_i\to\infty$ and let $h_i$ be walks with unit steps and block sizes $r_i\ge1$.
\begin{enumerate}[label=(\alph*)]
\item Every uniform limit on $[0,c_0]$ of the rescaled walks is $1$-Lipschitz.
\item If the walks start at $0$, then for any $k$ walks per size some subsequence of the rescaled
walks converges uniformly on $[0,c_0]$.
\item If $T\le cn^2b$ and a block holds a row ($nb\le rG$), there are at most $(cG+1)n$ slots.
\item[(V2$'$)] For a walk frozen at time $T$, every block strictly between its blocks at the first
steps of two slots $y_1\le y_2$ is pair-resident at some slot of $[y_1,y_2]$; and the pair residency
depends only on the walk before $T$.
\end{enumerate}
\end{lemma}

\begin{outline}
(a) A unit-step walk moves at most $r$ cells per slot, that is at most one block per $1/n$ of
macroscopic time, up to rounding that vanishes in the limit. (b) Arzelà--Ascoli in the form of a
diagonal argument at rational times (a countable product of compact intervals is compact and
metrizable) together with asymptotic equicontinuity; the classical statement does not apply
directly because the rescaled walks are step functions. (c) Arithmetic. (V2$'$) is
\cref{lem:residency}(V2) for frozen walks.
\end{outline}

\begin{lemma}[Machine instances]\label{lem:machine-walks}
\formal{Tape/Statements/Core.lean}{pairRes_V2, numSlots_le_horizon, rescaledTraj_limit_lipschitz, exists_subseq_tendstoUniformlyOn}
For every machine $M$, time constant $c$ and $n_i\to\infty$: some subsequence of the rescaled
work-head trajectories converges uniformly on $[0,c_0]$, every such limit is $1$-Lipschitz, and if
$T\le cL$ there are at most $c_0n$ slots. If $M$ halts at $T$ on $w$, its work heads satisfy
(V2$'$).
\end{lemma}

\begin{outline}
Instances of \cref{lem:walk-basics}.
\end{outline}

\subsection{Levels, labellings and the inequality \texorpdfstring{$(\star)$}{(*)}}

\begin{definition}[Levels and crossings]\label{def:levels}
\formal{Tape/DefsM.lean}{levelSet, ConsecIn, crossPairs, levelCross, levelExcess, FinUnionIcc, PiecewiseConstOn}
Let $P:\R\to\R$ and $c_0>0$. The \emph{level set} of $p$ is $\Lambda(p)=P^{-1}(p)\cap[0,c_0]$. An
\emph{excised set} $S$ is a finite union of closed intervals, and a \emph{labelling}
$\beta:\R\to\N$ is piecewise constant off $S$: constant on every interval of $[0,c_0]\setminus S$
that avoids a fixed finite set. For a level $p$ with $\Lambda(p)$ finite, $\kappa(p)$ is the number
of consecutive pairs $s<s'$ of $\Lambda(p)$ with $s,s'\notin S$ and $\beta(s)\ne\beta(s')$; for
$\Lambda(p)$ infinite, $\kappa(p)=0$. The \emph{excess} of the level is
$\#\Lambda(p)-1$ (truncated at $0$, and $\infty$ for an infinite level set).
\end{definition}

So $\kappa(p)$ counts the returns of the path to height $p$ that cross from one label class to
another. For $k$ paths $P_t$ on $[0,c_0]$, a number $\theta\ge0$, an excised set $S$ and a
labelling $\beta$ whose classes off $S$ have measure at most $\theta$, the inequality is
\[
  (\star)\qquad
  1\;\le\;\sum_{t<k}\int\kappa_t(p)\,dp\;+\;2\theta\;+\;2k\,\bigl|S\cap[0,c_0]\bigr| .
\]

\begin{definition}[The cut inequality]\label{def:cutineq}
\formal{Tape/DefsM.lean}{CutIneq}
Walks $h_t$ ($t<k$) of length $T$ satisfy the \emph{cut inequality} at size $n$ with block size
$r$, $G$ bits per cell, slack $D$ and boundary cost $s$, if for every labelling of the
$N=\lceil T/r\rceil$ slots with at most $m$ slots per label and $\Mb$ label changes,
\[
  (n\dotminus\Mb)(n\dotminus2m)\,b(n)\;\le\;(\Mb+1)D+\Mb\,s+rG\,\bigl(3k\Mb+\chicross\bigr),
\]
where $\chicross$ counts the chords of the pair residencies whose endpoints carry different labels, and
$x\dotminus y=\max(x-y,0)$.
\end{definition}

For the work heads of a machine on a trace class it follows from \cref{cor:isoperimetry-trace} (since
$\clog{(g+1)}\le G$), with $D$ the co-entropy of the class and $s=\clog{(q+1)}+(k+1)\clog{(2T+1)}$; the formal
proofs obtain it through the class code and the cut-set bound for codes (\cref{lem:class-code},
\cref{cor:isoperimetry-code}).

\subsection{Theorem M* for walks}

\begin{theorem}[Theorem M* for walks]\label{thm:mstar-walk}
\formal{Tape/Statements/Coding.lean}{theoremMStar_walk, theoremMStar_level_walk}
Let $n_i\to\infty$ and, eventually, let $k$ walks $h_{i,t}$ have unit steps before $T_i\le cL_i$
and satisfy the cut inequality at size $n_i$ with block size $\lceil n_ib(n_i)/G\rceil$, slack
$D_i$ and boundary cost $s_i$, where $(D_i+s_i)/L_i\to0$. Let $c_0=cG+1$ and consider the rescaled
walks frozen at $T_i$.
\begin{enumerate}[label=(\alph*)]
\item If the walks start at $0$, some subsequence converges uniformly on $[0,c_0]$.
\item Every subsequential uniform limit $(P_t)_{t<k}$ consists of $1$-Lipschitz paths and
satisfies $(\star)$ for every $\theta$, $S$ and $\beta$. Moreover
\[
  \sum_t\int(\#\Lambda_t(p)-1)\,dp\;\ge1,\qquad
  1+\sum_t\bigl|P_t([0,c_0])\bigr|\;\le\;k\,c_0,
\]
and the paths are not all monotone.
\end{enumerate}
\end{theorem}

\begin{outline}
The discrete part works at one size. From $S$ and $\beta$ build a slot labelling: the label of a
slot off $S$ is determined by $\beta$, and slots in $S$ are cut into short pieces. Its label changes
and class sizes are bounded by counting good and bad slots. The \emph{chord split}
(\lean{crossChords_split}) bounds the crossing chords by the label changes times $(\lceil\eta n\rceil+1)$, plus
twice the slots in $S$, plus the \emph{long} crossing chords, those spanning at least $\eta n$
slots. Long chords are counted by levels (\lean{levelChords_limsup_le}): for a good level $p$, the
long crossing chords at block $\lfloor pn\rfloor$ inject, for large $n$, into the pairs counted by
$\kappa(p)$, by uniform closeness, (V1) and (V2$'$). Reverse Fatou then gives
$\limsup_n\#\{\text{long crossing chords}\}/n\le\int\kappa(p)\,dp$, with an explicit integrable
majorant and no measurability assumption on $\kappa$. If $(\star)$ failed, the cut inequality would
fail at a single large size for explicit $\eta$ and $\delta$; the arithmetic is isolated in one
lemma. The remaining claims are the labelling-free consequences of \cref{thm:excess} below.
\end{outline}

The interval labellings and excised sets avoid the area formula and Borel measurability of first
returns: only finitely many breakpoints occur.

\subsection{Labelling-free consequences}

\begin{theorem}[Excess and range]\label{thm:excess}
\formal{Tape/Statements/Core.lean}{levelCross_le_levelExcess, one_le_excess_of_star, excess_add_range_le, range_add_one_le_of_star, lintegral_levelExcess_eq_zero_of_monotone, not_all_monotone_of_star}
For every path, excised set and labelling, $\kappa(p)\le\#\Lambda(p)-1$. Let $P_t$ ($t<k$) satisfy
$(\star)$ for every $\theta\in(0,\tfrac12)$, $S$ and $\beta$. Then:
\begin{enumerate}[label=(\alph*)]
\item the total excess multiplicity is at least one:
$\sum_t\int(\#\Lambda_t(p)-1)\,dp\ge1$;
\item if the $P_t$ are $1$-Lipschitz on $[0,c_0]$, then
$\int(\#\Lambda(p)-1)\,dp+\abs{P([0,c_0])}\le c_0$ for each $t$, and hence
$1+\sum_t\abs{P_t([0,c_0])}\le kc_0$;
\item a monotone path has total excess $\int(\#\Lambda(p)-1)\,dp=0$, so the $P_t$ are not all monotone.
\end{enumerate}
\end{theorem}

\begin{outline}
For~(a), take $S=\emptyset$ and the labelling $\max(0,\lfloor mu\rfloor)$ with $\theta=1/m$, and let
$m\to\infty$. Part~(b) is the easy half of Banach's indicatrix~\cite{Banach1925} for $1$-Lipschitz
paths, proved by dyadic counting with the range split off. For~(c), a level with two preimages
under a monotone path is a value on an interval of constancy, and there are countably many such
values.
\end{outline}

\begin{lemma}[Analysis tools]\label{lem:analysis}
\formal{Tape/Statements/Core.lean}{countable_localExtr_values, levelSet_encard_majorant, lintegral_encard_levelSet_le, ae_finite_levelSet, levelChords_limsup_le, crossChords_split}
For every $P:\R\to\R$ the values at local extrema form a countable set. For $P$ that is
$K$-Lipschitz on $[0,c_0]$, $\#\Lambda(p)$ has a measurable majorant with integral at most $Kc_0$;
hence $\int\#\Lambda(p)\,dp\le Kc_0$ and almost every level set is finite. Moreover:
\begin{enumerate}[label=(\alph*)]
\item (Chord split.) Let a residency $R$ satisfy (P2), let $\eta\in\R$, and let a slot labelling
$\lambda$ satisfy $\beta(y/n)\ne\beta(y'/n)$ whenever $y/n,y'/n\notin S$ and $\lambda(y)\ne\lambda(y')$.
Then the chords of $R$ with target $<N$ whose ends carry different labels number at most
$\Mb(\lceil\eta n\rceil+1)+2\,\#\{y<N:y/n\in S\}+\#\mathcal X$, where $\mathcal X$ is the set of
long chords $x\to z$: $z<N$, $z-x\ge\eta n$, $x/n,z/n\notin S$, $\beta(x/n)\ne\beta(z/n)$.
\item (Long chords by levels.) Let $\eta>0$, let $P$ be $1$-Lipschitz on $[0,c_0]$, and for large $i$ let
$W_i:\N\to\Z$, residencies $R_i$ and $N_i\le c_0n_i$ satisfy $\abs{W_i(y)-\sigma}\le D$ for
$\sigma\in R_i(y)$ and (V2$'$), with $u\mapsto W_i(\lfloor un_i\rfloor)/n_i\to P$ uniformly on $[0,c_0]$.
Then $\limsup_i\#\mathcal X_i/n_i\le\int\kappa(p)\,dp$ for every excised set $S$ and labelling $\beta$.
\end{enumerate}
\end{lemma}

\begin{outline}
Each value at a local extremum is the minimum or maximum of $P$ over a rational interval. The
majorant is $\liminf_mN_m(p)$, where $N_m(p)$ counts the dyadic subintervals $I$ of order $m$ such that $p$
lies within $K\abs{I}/2$ of the value of $P$ at the centre of $I$ (an interval of values containing $P(I)$).
Each $N_m$ is measurable with integral $Kc_0$; once $c_02^{-m}$ is below the gaps of a finite subset of
$\Lambda(p)$, $N_m(p)$ is at least its size, so $\liminf_mN_m\ge\#\Lambda$ at every level, and Fatou's lemma
bounds the integral. For~(a), a crossing chord shorter than $\eta n$ has its source among the
$\lceil\eta n\rceil$ slots before a label change; a chord touching $S$ is charged to its source or target slot
in $S$; by (P2) both charges are injective, and the remaining crossing chords are long. Part~(b) is the level
count in the outline of \cref{thm:mstar-walk}.
\end{outline}

\subsection{Machine forms}

\begin{theorem}[Theorem M*, oblivious machines]\label{thm:mstar}
\formal{Tape/Statements/Core.lean}{theoremMStar_level, theoremMStar_level_sub, theoremMStar_level_subseq, theoremMStar_excess, no_stationary_limits}
Let $n_i\to\infty$ and let $M$ be, eventually, oblivious on $\mathrm{Mat}(n_i,b(n_i))$ with halting
time $T_i\le cL_i$ and transpose. If the rescaled work-head trajectories converge uniformly on
$[0,c_0]$ to $P_t$, then the $P_t$ satisfy $(\star)$ for every $\theta$, $S$, $\beta$ (equivalently,
$\sum_t\int\kappa_t\ge1-2\theta-2k\abs{S\cap[0,c_0]}$ for $\theta\ge0$), and the conclusions of
\cref{thm:excess}. A convergent subsequence always exists. In particular the limits cannot all be
identically $0$.
\end{theorem}

\begin{theorem}[Theorem M* for trace classes and adaptive machines]\label{thm:mstar-adaptive}
\formal{Tape/Statements/Adaptive.lean}{rescaledTrajOf_close, tendstoUniformlyOn_of_mem, mstar_discrete_class_on, theoremMStar_level_class, theoremMStar_level_class_on, theoremMStar_class_nonuniform, theoremMStar_level_adaptive, no_stationary_limits_class_on}
Let $n_i\to\infty$.
\begin{enumerate}[label=(\alph*)]
\item Members of one trace class have rescaled trajectories within $4/n$ of each other, so limits
along trace classes do not depend on the representatives.
\item Along trace classes $F_i$ with halting times $T_i\le cL_i$ and co-entropies $D_i=o(L_i)$ on which $M$
is correct, the rescaled trajectories of the representatives satisfy the conclusion of \cref{thm:mstar-walk}
(the labelling-free parts by \cref{thm:excess}); the underlying finite-$n$ inequality holds at every size. The same holds for non-uniform machines
$M_i$ with $k$ tapes, $G$ bits per cell and $o(L_i)$ state bits.
\item If every matrix input halts within $cL_i$ steps and $M$ transposes, then the popular classes
(\cref{thm:popular}) have $D_i\le Cn_ib(n_i)$, with $C=C(M,c)$ the constant of \cref{thm:popular}
(\lean{traceConst}), a convergent subsequence of the representatives'
trajectories exists, and every subsequential limit is $1$-Lipschitz and satisfies $(\star)$.
\end{enumerate}
\end{theorem}

\begin{theorem}[Robust and non-uniform M*]\label{thm:mstar-robust}
\formal{Tape/Statements/Adaptive.lean}{theoremMStar_level_robust, theoremMStar_level_nonuniform, theoremMStar_excess_robust, theoremMStar_excess_nonuniform}
Let $n_i\to\infty$ and suppose that, eventually, $M$ halts within $cL_i$ steps with the
transposed output on every input of a set $W_i$ of co-entropy $m_i$, with $m_i/L_i\to0$; inputs
outside $W_i$ are unconstrained. Then the popular trace classes inside $W_i$ have co-entropy at most
$m_i+Cn_ib(n_i)$, convergent subsequences of the representatives' trajectories exist, and every
subsequential limit is $1$-Lipschitz, satisfies $(\star)$, has total excess at least $1$ and total
range at most $kc_0-1$, and its paths are not all monotone; here $C=C(M,c)$ is the constant of
\cref{thm:popular}. The same holds for machines $M_i$ depending on $i$, with $k$ work tapes, $G$ bits per cell
and $o(L_i)$ state bits (then $C$ depends only on $k$, $G$ and $c$).
\end{theorem}

\begin{outline}[Outline of the formal proofs of \cref{thm:mstar,thm:mstar-adaptive,thm:mstar-robust}]
All machine forms reduce to Theorem~M* for code-carrying diagrams (\cref{thm:mstar-diagram}),
hence to Theorem~M* for walks. For a trace class, the class code on the representative's diagram
(\cref{lem:class-code}) carries the state and head positions as
$\clog{(q+1)}+(k+1)\clog{(2T+1)}=o(L)$ extra bits per memory edge, and the cut-set bound for codes
gives the cut inequality. The oblivious machine is the case of one class with $D=0$
(\cref{lem:trace-obl}); the adaptive form is the robust form with $W_i$ all inputs; the robust form
extracts the popular class inside $W_i$ (\cref{thm:popular}, in its form within a set of inputs,
\lean{exists_popular_traceClass_on}). Representative independence holds
because the work-head blocks of two members of a class differ by at most $2$ before halting and $4$
after.
\end{outline}

\begin{corollary}[Theorem M* under (T)]\label{cor:mstar-T}
\formal{Tape/Statements/Adaptive.lean}{adaptive_hyps_of_linearTime, linearTimeTransposable_popular_limits}
If (T) holds, then some machine $M$ and constant $c$ have the following property along every
$n_i\to\infty$: there are trace classes of co-entropy $o(L_i)$ (the popular traces), a convergent
subsequence of their representatives' rescaled trajectories, and every subsequential limit is
$1$-Lipschitz and satisfies $(\star)$ on $[0,(2c)G+1]$.
\end{corollary}

\begin{outline}
A machine meeting \cref{def:T} with constant $c$ halts within $2cL$ steps on every matrix input of every
size $n\ge1$ and transposes (\cref{lem:bridge}); apply \cref{thm:mstar-adaptive}(c).
\end{outline}

So a linear-time transposer must move its work heads in a macroscopically folding way. The
condition is far from sufficient: it constrains only the macroscopic shape of the head motion. It does
exclude simple zig-zags: if every head follows the same triangle wave $P\ge0$ of slope $\pm1$, the labelling
$\beta(u)=\lfloor mP(u)\rfloor$ with $m$ large has classes of measure at most $\tfrac14$ and keeps every level
set inside one class, so $\kappa_t\equiv0$ and $(\star)$ fails at $\theta=\tfrac14$, $S=\emptyset$. Its value
is as a test that any candidate construction must pass, and as the limit object behind the coding results of
\cref{sec:codes}.

\begin{proposition}[$(\star)$ is a real constraint]\label{prop:star-fails}
\formal{Tape/Statements/Core.lean}{theoremMStar_conclusion_can_fail}
For the constant path $P=0$, $S=\emptyset$, $\beta(u)=\max(0,\lfloor4u\rfloor)$ and $\theta=\tfrac14$, the
hypotheses on $S$, $\beta$ and $\theta$ hold and $(\star)$ fails, for every $k$ and $c_0$.
\end{proposition}

\begin{outline}
$\kappa\le\#\Lambda-1$, and the excess of a constant path vanishes almost everywhere.
\end{outline}

\section{Codes on time--space diagrams}\label{sec:codes}

The block-level lower bounds of the previous sections (the cut lemmas for unions of slots, Discrete
Isoperimetry and Theorem~M*) use the transition function only through block locality: what they use is
that information flows along the block network of the run, with bounded capacity per edge. This section
makes that structure explicit. A \emph{time--space diagram} records only head motions and an output
schedule; its block network is a directed acyclic graph with bounded degree; and a \emph{causal network
code} on it is an arbitrary assignment of functions to edges. A machine that is correct on a trace class
yields such a code (\cref{lem:class-code}), and conversely these block-level results hold for codes
(\cref{lem:cutset,cor:isoperimetry-code,thm:mstar-diagram}); the cell-level results (the Locality Lemma,
Theorems~P and~P$'$) are not claimed for codes. \Cref{sec:conjecture} then asks whether such codes exist
at all.

\subsection{Diagrams and block networks}

\begin{definition}[Time--space diagram]\label{def:diagram}
\formal{Tape/DefsCE.lean}{TSDiagram, TSDiagram.Valid, runDiagram}
A $k$-tape \emph{time--space diagram} consists of a length $T$, walks $h_t:\N\to\Z$ for $t<k$ (the
work heads), a walk $h_{\mathrm{in}}$ (the input head) and an output schedule $\mathrm{out}:\N\to\N$
($\mathrm{out}(\tau)$ is the number of output bits before step $\tau$). It is \emph{valid} for output
length $L$ if all walks start at $0$ and have unit steps, and the output schedule starts at $0$,
grows by at most one per step, and equals $L$ from time $T$ on. The diagram of a run of $M$ on $w_0$
up to time $T$ consists of its head trajectories and output lengths.
\end{definition}

\begin{definition}[Block network]\label{def:blocknet}
\formal{Tape/DefsCE.lean}{TSDiagram.blockSize, TSDiagram.nodes, TSDiagram.memEdges, TSDiagram.chordEdges, TSDiagram.edges, TSDiagram.rowsAt, TSDiagram.outNode, TSDiagram.visible}
Let $w=b(d)$, $L=d^2w$, $G\ge1$ and $r=\lceil dw/G\rceil$. The block network of a diagram has as
nodes the slots $y<\lceil T/r\rceil$, in their causal order; \emph{memory edges} $y\to y+1$; and a
\emph{chord edge} $x\to z$ whenever some work tape has a chord from $x$ to $z$ of its pair
residency. The rows \emph{attached} at node $y$ (its sources) are the input rows that the input walk
touches during slot $y$; the \emph{sink} of output position $o$ is the slot during which bit $o$ is
appended.
\end{definition}

Every node has in-degree at most $k+1$: one memory edge, and at most one chord per tape by (P2)
(\cref{lem:P2}). There are at most $kN$ chord edges (\lean{card_chordEdges_le}).

\begin{definition}[Causal codes, complete exchange, capacities]\label{def:code}
\formal{Tape/DefsCE.lean}{NetCode, NetCode.inVals, NetCode.edgeVal, NetCode.outVal, TSDiagram.Solves, TSDiagram.CapProfile, TSDiagram.blockCap, TSDiagram.profileCap}
A \emph{causal code} assigns to every edge $e=(u,v)$ a function $\mathrm{enc}_e$ computing a natural
number from the values on the edges into $u$ and the input bits visible at $u$ (the bits of the rows
attached at $u$), and to every output position $o$ a function $\mathrm{dec}_o$ computing a bit from
the same data at the sink of $o$. The code is evaluated along the spine. It \emph{solves complete
exchange} on a set $F$ of inputs within capacities $\mathrm{cap}$ if for every input in $F$ every edge
value is below $2^{\mathrm{cap}(e)}$ and every output bit is the transposed input bit. The capacities
are \emph{within the block profile} if a memory edge $y\to y+1$ carries at most the blocks resident at
slot $y$ (summed over the tapes, at most $3k$) and a chord edge carries one block per tape having that
chord, a block being $rG$ bits. The \emph{profile capacities} are the largest such capacities; with
$s$ extra bits on every memory edge we speak of the block capacities with $s$ extra bits.
\end{definition}

Being a solving code is monotone in the capacities (\lean{solves_mono}), so ruling out codes at the
profile capacities rules them out within every capacity in the profile.

The code may compute arbitrary functions. The only restrictions are causality, the capacities, and
the locality of the inputs: a row is visible only at the nodes where the input walk touches it.

\subsection{The cut-set bound}

\begin{lemma}[Cut-set bound for codes]\label{lem:cutset}
\formal{Tape/Statements/Coding.lean}{cutset_bound_bitwise, cutset_bits_bitwise, cutset_bound, card_le_of_solves}
Let a causal code solve complete exchange on $F$ within $\mathrm{cap}$, and let $U$ be a set of
nodes. Let $B_U$ be the set of input bits whose row is attached at no node of $U$ and whose
transposed position has its sink in $U$. Then
\[
  2^{\abs{B_U}}\,\abs{F}\;\le\;2^{L}\cdot2^{\sum_{e\text{ entering }U}\mathrm{cap}(e)} ,
\]
so if $F$ has co-entropy at most $D$, then $\abs{B_U}\le D+\sum_{e\text{ entering }U}\mathrm{cap}(e)$.
In particular $2^{w\abs{S_U}}\abs{F}\le2^L2^{\sum\mathrm{cap}}$, where $S_U$ is the set of pairs of a row
not attached in $U$ and a column all of whose positions have their sinks in $U$. Taking $U$ to be all
nodes: rows attached nowhere cost co-entropy, $2^{u\,dw}\abs{F}\le2^L$ for $u$ such rows.
\end{lemma}

\begin{outline}
Group the inputs of $F$ by their bits outside $B_U$. Within a group, the values computed inside $U$
are a function of the values on the edges entering $U$ and of the visible bits in $U$, which do not
depend on $B_U$; the outputs in $U$ determine $B_U$. So each group injects into the values on the
entering edges.
\end{outline}

\begin{corollary}[Discrete Isoperimetry for codes]\label{cor:isoperimetry-code}
\formal{Tape/Statements/Coding.lean}{isoperimetry_code, sum_cross_blockCap_le, cutIneq_of_solves}
Let a causal code solve complete exchange on a set of co-entropy at most $D$, on the block network
of a valid diagram with $r\ge1$. For every labelling of the nodes with $\Mb$ label changes and at
most $m$ nodes per label,
\[
  (d\dotminus\Mb)(d\dotminus2m)\,w\;\le\;(\Mb+1)D+\sum_{e:\;\text{labels of the ends differ}}\mathrm{cap}(e),
\]
where $x\dotminus y=\max(x-y,0)$.
At the block capacities with $s$ extra bits, the right-hand sum is at most
$\Mb s+rG(3k\Mb+\chicross)$, so the walks of the diagram satisfy the cut inequality
(\cref{def:cutineq}) with slack $D$ and boundary cost $s$.
\end{corollary}

\begin{outline}
Sum the bits form of \cref{lem:cutset} over the label classes, bound the left side as in
\cref{cor:isoperimetry}, and observe that an edge between different labels is a memory edge at a
label change or a crossing chord.
\end{outline}

\subsection{Distances, flows and the code gap}

\begin{definition}[Distances and routings]\label{def:routing}
\formal{Tape/DefsCE.lean}{ball, netDist, TSDiagram.rowDist, IsWalk, Routing, TSDiagram.RoutesAt}
$\mathrm{dist}$ is the undirected hop distance in the block network. For an input bit $\alpha$,
$\mathrm{rowDist}(\alpha)$ is the distance from the nearest node at which its row is attached to the sink
of its transposed position ($0$ if the row is attached nowhere). A \emph{routing at rate $\lambda$}
within $\mathrm{cap}$ is a fractional multicommodity flow, one commodity per input bit whose row is
attached somewhere, sending $\lambda$ units of each commodity along walks from a node attaching its
row to its sink, with total traversals of every edge at most its capacity.
\end{definition}

\begin{lemma}[Weak duality and volume]\label{lem:volume}
\formal{Tape/Statements/Coding.lean}{weak_duality, rowDist_sum_ge, card_chordEdges_le, sum_cap_le}
A routing at rate $\lambda$ satisfies $\lambda\sum_\alpha\mathrm{rowDist}(\alpha)\le\sum_e\mathrm{cap}(e)$. For a
valid diagram with $1\le r\le dw$ and every $j$,
\[
  j\,L\;\le\;\sum_\alpha\mathrm{rowDist}(\alpha)+j\bigl(4N(2k+3)^{j-1}w+u\,dw\bigr),
\]
where $N$ is the number of nodes and $u$ the number of rows attached nowhere. There are at most $kN$
chord edges, and capacities within the profile sum to at most $4kNrG$.
\end{lemma}

\begin{outline}
Weak duality: each unit of commodity $\alpha$ traverses at least $\mathrm{rowDist}(\alpha)$ edges. For the
distance bound, the ball of radius $j-1$ around a node attaching a row has at most $(2k+3)^{j-1}$ nodes
(degree at most $2k+2$), each node attaches at most two rows, and a sink receives at most $2w$ bits of a
row, so few bits are within distance $j-1$ of their sinks.
\end{outline}

\begin{theorem}[The code gap]\label{thm:code-gap}
\formal{Tape/Statements/Coding.lean}{volume_gap, flow_rate_lt, code_gap}
Fix $k$, $G\ge1$, $c$ and $C$, and let $E=4\clog{(2k+3)}(4k+1)G(c+1)$. For all large $d$, on every valid
$k$-tape diagram of length at most $cL$ with capacities within the profile and at most $C$ rows
attached nowhere,
\[
  w\sum_e\mathrm{cap}(e)\;<\;E\sum_\alpha\mathrm{rowDist}(\alpha),
\]
so every routing has rate $<E/w$. In particular, if a causal code solves complete exchange on a set of
co-entropy at most $Cdw$ within such capacities, the inequality holds: the code (rate~$1$) beats
multicommodity flow by a factor larger than $w/E=\Omega(\log d)$.
\end{theorem}

\begin{outline}
Take $j=\lfloor(w-1)/(2\clog{(2k+3)})\rfloor$, so that $(2k+3)^{2j}\le d-1$. By \cref{lem:volume}, all
but $o(L)$ bits are at distance at least $j=\Omega(w)$ from their sinks, while
$\sum\mathrm{cap}\le4kG(c+1)L$. Weak duality bounds the rate. For a solving code, co-entropy $Cdw$
leaves at most $C$ rows attached nowhere (\cref{lem:cutset}).
\end{outline}

So a code on a block network of a linear-time diagram, if it exists, beats routing by a factor
$\Omega(\log d)$. The connection goes back to Adler, Harvey, Jain, Kleinberg and
Lehman~\cite[Thm.~6.2]{AHJKL2006}, who turn an oblivious algorithm for transposition into a network code on
its computation graph; here \cref{lem:class-code,thm:fast-codes} turn every linear-time transposer into such a
code on its own block network (one way only: a code need not come from a machine). The Tape Coding
Conjecture says that no such code exists.

\subsection{Machines give codes}

\begin{lemma}[Class codes]\label{lem:class-code}
\formal{Tape/Statements/Coding.lean}{traceClass_code}
Let $d\ge1$, let $F$ be a trace class with representative $w_0$ (slot length $r=\lceil dw/G\rceil$,
$G=\max(1,\clog{(g+1)})$) on which $M$ is correct. Then the diagram of the run on $w_0$ up to $T$ is
valid, and its block network carries a causal code that solves complete exchange on $F$, where a
memory edge $y\to y+1$ carries the state, the $k+1$ head positions and the contents of the resident
blocks, at most $\clog{(q+1)}+(k+1)\clog{(2T+1)}+r\clog{(g+1)}\sum_t\abs{R_t(y)}$ bits, a chord edge
carries $r\clog{(g+1)}$ bits per tape having that chord, and no other pair carries anything.
\end{lemma}

\begin{outline}
This is the Block Locality Lemma in functional form. A memory edge $y\to y+1$ carries the state, the head
positions and the contents of the resident blocks at time $(y+1)r$, and a chord edge the contents of its
blocks at the time of its target slot. The code at a node chooses a member of $F$ consistent with the values
it receives and its visible input bits, and sends that member's payloads; the decoder outputs that member's
bit. By \cref{lem:block-locality} for the single slot, all consistent members have the same outgoing payloads
and output bits, and an induction over the slots shows that on every input of $F$ each edge carries that
input's own payload.
\end{outline}

\begin{theorem}[Fast machines give codes]\label{thm:fast-codes}
\formal{Tape/Statements/Coding.lean}{exists_popular_traceClass_on, codes_of_fastOn_bits, solves_profileCap_addTapes, codes_of_fastOn}
Let $M$ be correct within $cL$ steps on every input of a set $F_0$ of co-entropy at most $C_0dw$.
\begin{enumerate}[label=(\alph*)]
\item If $s\ge\clog{(q+1)}+(k+1)\clog{(2cL+1)}$, then there are a valid $k$-tape diagram of length at most
$cL$ (in the proof, the diagram of $M$'s run on the representative of a popular class) and a causal code on
it at the block capacities with $s$ extra bits ($G=\max(1,\clog{(g+1)})$, $r=\lceil dw/G\rceil$), solving
complete exchange on a set of co-entropy at most $(C_0+C_M)dw$, where $C_M$ depends only on $M$ and $c$.
\item Extra bits become tapes: a code at the block capacities with $s\le a\,rG$ extra bits is a code at
the profile capacities of the diagram with $a$ additional stationary tapes.
\item Hence, if $2\clog{(q+1)}\le dw$ and $d$ is large enough, there is such a code at the profile
capacities on a diagram with $k+1$ tapes.
\end{enumerate}
\end{theorem}

\begin{outline}
(a) The popular trace inside $F_0$ (\cref{thm:popular}, which holds within any set of inputs) is a
class of co-entropy at most $(C_0+C_M)dw$; apply \cref{lem:class-code}. (b) A stationary tape changes
no node, edge, source, sink or chord, and its single resident block carries the extra bits on every
memory edge. (c) Combine (a) and (b) with $a=1$.
\end{outline}

\subsection{Theorem M* for code-carrying diagrams}

Theorem~M* holds for codes: its proof uses the machine only through the cut inequality.

\begin{theorem}[Theorem M* for code-carrying diagrams]\label{thm:mstar-diagram}
\formal{Tape/Statements/Coding.lean}{theoremMStar_level_diagram, theoremMStar_diagram, theoremMStar_diagram_profile}
Let $d_i\to\infty$, and let valid $k$-tape diagrams $D_i$ of length at most $cL_i$ carry causal codes
solving complete exchange on sets of co-entropy $m_i$, at the block capacities with $s_i$ extra bits,
where $(m_i+s_i)/L_i\to0$. Then the rescaled work walks of the $D_i$ have a uniformly convergent
subsequence on $[0,cG+1]$, and every subsequential limit is $1$-Lipschitz and satisfies $(\star)$. At
the profile capacities ($s_i=0$) every limit also has total excess at least $1$, total range at most
$k(cG+1)-1$, and its paths are not all monotone.
\end{theorem}

\begin{outline}
By \cref{cor:isoperimetry-code} the walks satisfy the cut inequality with slack $m_i$ and boundary cost
$s_i$; apply Theorem~M* for walks (\cref{thm:mstar-walk}).
\end{outline}

\begin{corollary}[Counterexamples to the conjecture obey $(\star)$]\label{cor:cex-star}
\formal{Tape/Statements/Coding.lean}{mstar_of_robustNoCode_counterexamples, mstar_of_not_robustNoCode}
Every family of counterexamples to the Tape Coding Conjecture (\cref{conj:TCC}): valid diagrams of
length at most $cL$ with codes at the profile capacities on sets of co-entropy $Cdw$, along
$d_i\to\infty$, has rescaled work walks whose subsequential limits are $1$-Lipschitz, satisfy $(\star)$,
have total excess at least $1$ and total range at most $k(cG+1)-1$, and are not all monotone. If the
conjecture fails, such a family exists, with a convergent subsequence.
\end{corollary}

\begin{outline}
\Cref{thm:mstar-diagram} with $m_i=Cd_iw_i=o(L_i)$ and $s_i=0$.
\end{outline}

\section{The Tape Coding Conjecture}\label{sec:conjecture}

\subsection{Statement}

\begin{conjecture}[Tape Coding Conjecture]\label{conj:TCC}
\formal{Tape/DefsCE.lean}{RobustNoCode_TM}
For all $k$, $G\ge1$, $c$ and $C$, the following holds for all sufficiently large $n$ (with
$w=b(n)$, $L=n^2w$ and $r=\lceil nw/G\rceil$). No causal code solves complete exchange at the profile
capacities on the block network of a valid $k$-tape time--space diagram of length at most $cL$, for a
set of inputs of co-entropy at most $Cnw$.
\end{conjecture}

In words: on the block network of any linear-length tape diagram, complete exchange cannot be
realized by any coding scheme, even if the scheme only has to work for a $2^{-O(n\log n)}$ fraction
of the inputs. By \cref{thm:code-gap}, a solving code would beat multicommodity flow by a factor
$\Omega(\log n)$. The conjecture can therefore be read as ``on tape networks, coding beats routing by
$o(\log n)$'', and this reading is one of its equivalent forms (\cref{thm:tcc-forms}). The internal
name of the conjecture in the Lean library is \lean{RobustNoCode_TM} (``no code for \emph{complete
exchange} on \emph{tape-machine} networks, robust form'').

The \emph{plain} form, \lean{NoCode_TM}, asks the same only for codes that are correct on all inputs.
It refutes oblivious linear-time transposers (\cref{thm:oblivious}). Whether plain implies robust is
the subject of \cref{sec:plainrobust}.

\subsection{The reduction}

\begin{theorem}[The Tape Coding Conjecture refutes (T)]\label{thm:reduction}
\formal{Tape/Statements/Coding.lean}{not_linearTimeTransposable_of_robustNoCode, not_linearTimeTransposable_of_robustCE_TM, not_linearTimeTransposable_of_robustCE_TM_dist, codes_of_linearTimeTransposable, not_robustCE_TM_of_linearTimeTransposable}
If the Tape Coding Conjecture holds (in any of its robust forms of \cref{thm:tcc-forms}), then
$\neg\,$\lean{Transpose.LinearTimeTransposable}. In contrapositive form, and more strongly, a linear-time
transposer yields constants $k$, $G\ge1$, $c$, $C$ such that for \emph{every} large $n$ there is a valid
$k$-tape diagram of length at most $cL$ whose block network carries a causal code, within capacities in the
block profile, solving complete exchange on a set of co-entropy at most $Cnw$ (a refutation of the
conjecture needs such codes only for infinitely many $n$).
\end{theorem}

\begin{outline}
The formal proof goes through \cref{thm:coent-linear}. The conjecture rules out codes at co-entropy
$Cn\le Cnw$ for every $C$. A machine meeting \cref{def:T} with constant $c$ is correct within $2cL$ steps on every matrix input of every size $n\ge1$; after output buffering, the popular class of its coarse trace has co-entropy $O(n)$, and its
class code (\cref{lem:class-code}), with the state and the head positions carried by extra stationary
tapes (\cref{thm:fast-codes}(b)), solves complete exchange at the profile capacities at every large
size. This contradicts the conjecture at a single large $n$. The flow and distance forms are
equivalent (\cref{thm:tcc-forms}), and the second statement follows because $Cn\le Cnw$.
\end{outline}

Concretely, a proof \texttt{h} of \lean{RobustNoCode_TM} resolves~(T) by the file
\begin{center}
\texttt{theorem Transpose.answer : $\neg$ Transpose.LinearTimeTransposable :=}\\
\texttt{\ \ Transpose.Tape.not\_linearTimeTransposable\_of\_robustNoCode h}
\end{center}
which \leanfile{Verify.lean} (in \texttt{lean/}) accepts, once \leanfile{Transpose/Answer.lean} also imports
\leanfile{Transpose.Tape.Statements.Coding}. The reduction is one-way. A refutation of the conjecture is a
choice of constants $k$, $G$, $c$, $C$ with codes for infinitely many $n$; nothing makes these codes uniform,
finite-state, or local to slots, as a machine's code would be, and no converse of \cref{thm:reduction} is
proved. So a refutation would remove this coding obstruction but would not by itself decide (T).

\begin{theorem}[Consequences for every machine]\label{thm:tcc-consequences}
\formal{Tape/Statements/Coding.lean}{linearTime_fails_of_robustNoCode, no_fast_transposer_of_robustNoCode, no_fast_transposer_of_robustNoCode_bits, linearTime_fails_of_robustNoCode_logBits, not_linearTimeTransposable_of_robustNoCode_logBits}
Assume the Tape Coding Conjecture.
\begin{enumerate}[label=(\alph*)]
\item For every machine $M$ and constant $c$, for all large $n$ some $n\times n$ matrix is not
transposed within $c(L+1)$ steps: every machine fails at \emph{every} large size.
\item (Non-uniform, average case.) For all $k$, $g$, $c$, $C_0$, $C_1$ and all large $n$, no machine
with $k$ work tapes, $g+1$ work symbols and at most $2^{C_1nw}$ states (its transition table may
depend on $n$) halts within $cL$ steps with the transposed output on every input of a set of
co-entropy at most $C_0nw$.
\item (Extra bits.) For every $s$, if no code exists at the block capacities with $s(k,n)$ extra bits on
every memory edge (\lean{RobustNoCode_TM_bits}), then for all $k$, $g$, $c$, $C_0$ and all large $n$ no
machine with $k$ work tapes, $g+1$ work symbols and $\clog{(q+1)}+(k+1)\clog{(2cL+1)}\le s(k,n)$ halts
within $cL$ steps with the transposed output on every input of a set of co-entropy at most $C_0nw$.
For $s(k,n)=(3k+4)\clog n$ this hypothesis also implies (a) and $\neg$(T).
\end{enumerate}
\end{theorem}

\begin{outline}
Corollaries of the reduction at co-entropy $O(n)$ (\cref{thm:coent-linear}), whose single-size form
applies to each machine separately. The form with extra bits per memory edge is direct: a run's
payload on a memory edge is exactly its resident blocks, its state and its head positions.
\end{outline}

\subsection{Equivalent forms}

Two further forms make the connection with network coding explicit. The \emph{flow form}
\lean{RobustCE_TM} says: for every $\varepsilon>0$ and all large $n$, whenever a causal code solves
complete exchange within capacities $\mathrm{cap}$ in the block profile on a set of co-entropy at most
$Cnw$, the demand is routable at rate $1/(\varepsilon w)$ within $\mathrm{cap}$; that is, coding beats
multicommodity flow by a factor of at most $\varepsilon\log n$. The \emph{distance form}
\lean{RobustCE_TM_dist} replaces the flow rate by its volume bound:
$\sum_d\mathrm{rowDist}(d)\le\varepsilon w\sum_e\mathrm{cap}(e)$.

\begin{theorem}[Equivalent forms]\label{thm:tcc-forms}
\formal{Tape/Statements/Coding.lean}{robust_tfae, robustCE_TM_iff_robustNoCode, robustCE_TM_dist_iff_robustNoCode, robustNoCode_iff_attached, robustCE_TM_att_iff, robustCE_TM_att_iff_robustCE_TM, robustCE_TM_gamma_iff, robustCE_TM_gamma_per_iff, robustNoCode_of_gainBound, plain_tfae, ce_TM_iff_noCode, ce_TM_dist_iff_noCode, noCode_of_robustNoCode, ce_TM_of_robustCE_TM, ce_TM_dist_of_robustCE_TM_dist}
The following are equivalent: the Tape Coding Conjecture; its flow form; its distance form; the
no-code and flow forms restricted to diagrams whose input walk attaches every row; the no-code form
with the uniform profile of $3k$ blocks on every memory edge; and the no-code form with
$(3k+4)\clog n$ extra bits on every memory edge. The $\gamma$-form ``for some $\gamma(n)=o(\log n)$ with
$\gamma>0$ (one $\gamma$ for all constants, or one per choice of constants), every solving code's demand is
routable at rate $1/\gamma(n)$'' is equivalent as well. The explicit gain bound ``every code solving complete
exchange at the profile capacities has $E\sum_d\mathrm{rowDist}(d)\le w\sum_e\mathrm{cap}(e)$'', with $E$ as in
\cref{thm:code-gap}, implies the conjecture (the converse holds vacuously but is not formalized). The plain
forms (flow, distance, no code, for codes correct on all inputs) are equivalent to each other, and are
implied by the robust ones.
\end{theorem}

\begin{outline}
A solving code at the profile capacities makes the distance form fail by \cref{thm:code-gap}, and
weak duality (\cref{lem:volume}) turns any routing at rate $1/(\varepsilon w)$ into the distance form.
So each form is equivalent to ``no code''. The restriction to attached rows loses nothing: append a
final input sweep, which keeps every original node, edge and sink and the residency of every complete
slot, and let the code mask each node's input to its original rows (the time constant becomes $2c+3$).
The profile forms are \cref{thm:profile}.
\end{outline}

\begin{theorem}[The capacity profile is immaterial]\label{thm:profile}
\formal{Tape/Statements/Coding.lean}{robustNoCode_with_iff, robustNoCode_with_k_iff, exists_shift_iff_unbounded, robustNoCode_with_of_robustNoCode, robustNoCode_bits_iff, robustNoCode_logBits_iff, robustNoCode_iff_forall_cap, noCode_iff_forall_cap}
For every unbounded memory profile $m$ (equivalently, $3k\le m(k+a)$ for some $a$, for every $k$; for
instance $m(k)=3k$, $3k+1$, or one block per tape), the no-code statement with $m(k)$ blocks on every
memory edge is equivalent to the Tape Coding Conjecture. So is the no-code statement with $s(k,n)$
extra bits on every memory edge, for every $s$ with $s(k,n)\le C(k)\,nw$; for instance
$s=(3k+4)\clog n$. The conjecture stated at the profile capacities is equivalent to ruling out codes
within every capacity in the profile, and the same holds for the plain form.
\end{theorem}

\begin{outline}
Extra bits and extra blocks are traded for stationary tapes (\cref{thm:fast-codes}(b)): a stationary
tape changes no node, edge, source, sink or chord, and its resident block can carry any payload of at
most one block. Conversely the per-node profile is at most any uniform profile with enough tapes.
\end{outline}

\subsection{Weaker sufficient conditions}

The conjecture asks for no codes at co-entropy $Cnw$, at all large $n$. Much less suffices to
refute~(T).

\begin{theorem}[Co-entropy $O(n)$ suffices]\label{thm:coent-linear}
\formal{Tape/Statements/Coding.lean}{bufTM_run, codes_coentLinear_of_linearTimeTransposable, not_linearTimeTransposable_of_coentLinear, linearTime_fails_of_coentLinear, no_fast_transposer_of_coentLinear, not_linearTimeTransposable_of_coentLinear_end, not_linearTimeTransposable_of_coentLinear_end_io, coentLinear_of_robustNoCode, noCode_of_coentLinear}
Write $\mathrm{NoCode}(m)$ for the statement ``for all $k$, $G\ge1$, $c$ and all large $n$, no causal
code solves complete exchange at the profile capacities on a valid $k$-tape diagram of length at most
$cL$, for a set of co-entropy at most $m(n)$'' (\lean{NoCode_TM_coent}).
\begin{enumerate}[label=(\alph*)]
\item (T) yields codes at co-entropy $Cn$ (not $Cnw$), on diagrams that append their output in their
last $L$ steps.
\item Hence, if $\mathrm{NoCode}(Cn)$ holds for every $C$, then $\neg$(T); every machine fails at every
large size; and the non-uniform average-case statement of \cref{thm:tcc-consequences}(b) holds with
co-entropy $C_0n$. It suffices to rule out codes on diagrams with the output at the end, and for
infinitely many $n$.
\item This hypothesis lies between the plain and the robust form: the Tape Coding Conjecture implies
it, and it implies $\mathrm{NoCode}(0)$, the plain form.
\end{enumerate}
\end{theorem}

\begin{outline}
First make the machine write its output only at the end (\lean{Impl.bufTM}): it writes the output on an
extra tape, rewinds, and copies it out, halting at $T+2L+2$ with the same output. Then the output length
at every slot is fixed by the halting time, and a \emph{coarse trace} suffices: block and row moves in
$\{-1,0,1\}$ together with the resident-block and touched-row patterns within $\pm1$, at most
$24^{k+1}$ values per slot by (V1) and (V2) of \cref{lem:residency} (for the input head: a slot is at most
one row long), which determine the block-level trace. The popular coarse
class therefore has co-entropy $O(n)$, and the class code applies unchanged.
\end{outline}

\begin{theorem}[One superlinear level, infinitely many sizes]\label{thm:io}
\formal{Tape/Statements/Coding.lean}{not_linearTimeTransposable_of_coent, not_linearTimeTransposable_of_coent_nw, not_linearTimeTransposable_of_coentLinear_io, not_linearTimeTransposable_of_coent_io, coent_io_of_coent, coentLinear_io_of_robustNoCode_io, coentLinear_io_iff_exists_coent_io, coent_io_nw_of_robustNoCode_io, not_linearTimeTransposable_of_robustNoCode_io, not_linearTimeTransposable_of_robustCE_TM_io, robustNoCode_io_of_robustNoCode, robust_io_tfae, no_fast_transposer_of_coent, traceClass_coent_lt}
If $f(n)\ge Cn$ eventually for every $C$ (for instance $f(n)=nw$), then $\mathrm{NoCode}(f)$ refutes
(T); so does $\mathrm{NoCode}(f)$ for infinitely many $n$. The infinitely-often forms of the conjecture
(flow, distance, no code) are equivalent to each other, formally weaker than the conjecture, and still
refute (T). For infinitely many $n$, the single superlinear level and the linear family are equivalent:
$\mathrm{NoCode}(Cn)$ holds for infinitely many $n$, for every $C$, if and only if $\mathrm{NoCode}(f)$
holds for infinitely many $n$ for some $f$ with $f(n)\ge Cn$ eventually for every $C$. Under
$\mathrm{NoCode}(m)$, for all large $n$: no machine with $k$ work tapes, $g+1$ symbols and $2^{O(nw)}$
states is correct within $cL$ steps on a set of co-entropy $m_0$ with $m_0+Kn\le m(n)$, for a constant $K$
depending only on $k$, $g$ and $c$; and every trace class on which such a machine is correct within $cL$
steps has co-entropy larger than $m(n)$.
\end{theorem}

\begin{outline}
(T) gives codes at every large $n$, so a single failing size per choice of constants contradicts it.
For a single level, $Cn\le f(n)$ from some point on. Of the equivalence of the single-level and the
linear infinitely-often families, one direction is monotonicity in the level and the other a diagonal
argument.
\end{outline}

\begin{theorem}[Further equivalent and sufficient forms]\label{thm:rungs}
\formal{Tape/Statements/Coding.lean}{coentLinear_of_robustCE_TM_lin, robustCE_TM_lin_iff, not_linearTimeTransposable_of_robustCE_TM_lin, hintCodes_linear_of_linearTimeTransposable, hintLinear_iff_coentLinear, not_linearTimeTransposable_of_noCode_hintLinear}
The flow form at co-entropy $Cn$ is equivalent to $\mathrm{NoCode}(Cn)$ for every $C$, and refutes (T).
(T) gives codes indexed by an $O(n)$-bit oracle hint (\cref{sec:plainrobust}), and ruling out codes with
a $Cn$-bit hint for every $C$ is equivalent to $\mathrm{NoCode}(Cn)$ for every $C$, so it refutes (T)
as well.
\end{theorem}

\begin{outline}
As in \cref{thm:tcc-forms}, at the linear rung; hints and co-entropy are interchanged by
\cref{thm:hint}.
\end{outline}

\subsection{Oblivious machines}

\begin{theorem}[The plain conjecture refutes oblivious (T)]\label{thm:oblivious}
\formal{Tape/Statements/Coding.lean}{codes_of_obliviousLinearTimeTransposable, not_obliviousLinearTimeTransposable_of_noCode, not_obliviousLinearTimeTransposable_of_ce_TM, not_obliviousLinearTimeTransposable_of_ce_TM_dist, not_obliviousLinearTimeTransposable_of_noCode_io, noCode_io_of_robustNoCode_io}
Let \lean{ObliviousLinearTimeTransposable} be (T) for a machine that is oblivious at every size. An
oblivious linear-time transposer gives codes for independent uniform sources (all inputs). Hence the
plain form of the conjecture, in any of its forms, refutes oblivious (T); so does the plain no-code
form for infinitely many $n$. The robust infinitely-often form implies the plain infinitely-often form.
\end{theorem}

\begin{outline}
For an oblivious machine all inputs form one trace class with $D=0$ (\cref{lem:trace-obl}), so the
class code works on all inputs, with the state and the head positions on one extra stationary tape
(\cref{thm:fast-codes}(b)).
\end{outline}

\section{Plain versus robust}\label{sec:plainrobust}

The plain conjecture rules out codes that are correct on \emph{all} inputs; the robust one rules out
codes correct on a $2^{-O(nw)}$ fraction of the inputs. Robust implies plain trivially. The plain
form already refutes oblivious linear-time transposers (\cref{thm:oblivious}); if plain implied
robust, it would refute~(T). This section isolates the exact gap and proves the implication for
several natural classes of codes and sets. The general case is open.

\subsection{The gap is an oracle hint}

A \emph{hint code} with $2^s$ hint values consists of a function $h$ of the whole input and a family
of causal codes $(C_v)_v$ such that on every input $x$ the code $C_{h(x)}$ solves complete exchange
for $x$. It models an oracle that sees the whole input and tells every node $s$ bits; the hint is
uncharged (no edge carries it and no node computes it). Write $\mathrm{NoCode}^{\mathrm{hint}}(s)$
for ``no hint code with $2^{s(n)}$ values exists'' (\lean{NoCode_TM_hintBits}), with the same
quantifiers as the conjecture.

\begin{theorem}[Robust is plain with an oracle hint]\label{thm:hint}
\formal{Tape/Statements/Coding.lean}{robustNoCode_iff_hint, not_linearTimeTransposable_of_noCode_hint, hintCodes_of_linearTimeTransposable, noCode_of_noCode_hint, solves_shiftCode, exists_shift_cover, hint_of_code, code_of_hint, robustNoCode_iff_coent, noCode_iff_coent_zero, coent_mono, hintBits_of_coent, coent_of_hintBits, noCode_bits_iff}
\begin{enumerate}[label=(\alph*)]
\item On a fixed diagram and at fixed capacities: a hint code with $2^s$ values gives a code on a set
of co-entropy $s$; a code on a set of co-entropy $m$ gives a hint code with $2^m(L+1)$ values.
\item The Tape Coding Conjecture is equivalent to: for every $C$, no hint code with an $Cnw$-bit
hint exists. With a constant hint this gives the plain form. (T) yields, for every large $n$, a single
valid diagram with codes indexed by an $O(nw)$-bit hint solving complete exchange on every input.
\item The co-entropy ladder: $\mathrm{NoCode}(0)$ is the plain form, $\mathrm{NoCode}(Cnw)$ for every
$C$ is the robust form, the rungs are monotone, and hints and co-entropy agree up to $3w+1$ bits.
\item A side channel whose contents the code computes itself is free: for $s(k,n)\le C(k)\,nw$, the
plain form with $s$ extra bits on every memory edge is equivalent to the plain form.
\end{enumerate}
\end{theorem}

\begin{outline}
(a) In one direction take the most popular hint value. In the other, a code for $F$ is a code for
every shift $F\oplus z$ (each node XORs its visible rows with $z$, and each output bit is XORed back),
and $2^m(L+1)$ shifts of $F$ cover every input, by counting with Bernoulli's inequality. (b) and (c)
follow from (a). (d) Extra memory bits become stationary tapes, as in \cref{thm:profile}.
\end{outline}

So the entire gap between plain and robust is whether an $O(nw)$-bit oracle broadcast can be
replaced by one that the network computes itself. For the application to~(T), co-entropy $O(n)$
suffices (\cref{thm:coent-linear}), so only hints of $O(n)$ bits matter; this is still $\omega(1)$
bits. The first rung of the ladder is already open: does ``no code on all inputs'' imply ``no code
on half of the inputs''? This is a one-shot question of zero error versus constant error for this
class of networks.

\subsection{Partial results}

Throughout this subsection we assume the plain conjecture \lean{NoCode_TM} and conclude, for all
constants and all large $n$, that certain codes do not exist at the profile capacities. The
constructions prepend or append computations to the diagram (an input sweep, a selection phase,
copies of the work walks), so they change $k$ and $c$ by constant factors; this is harmless because
the plain conjecture quantifies over all constants.

\begin{theorem}[Affine sets and additive hints]\label{thm:pr-affine}
\formal{Tape/Statements/PlainRobust.lean}{PlainRobust.noCode_affine_of_noCode, PlainRobust.noCode_of_affine_subset, PlainRobust.noCode_additive_of_noCode, PlainRobust.noCode_additiveCover_of_noCode, PlainRobust.noCode_affine_exc_of_noCode}
Under the plain conjecture, no code exists on a set $F$ that contains a nonempty affine set
$\{x:\sum_ig_i(\text{row }i\text{ of }x)=b\}$, where the $g_i$ are $\F_2$-linear maps to $\F_2^s$ with
$s\le Cnw$; more generally, on a set $F$ such that, for an additive hint $\sum_ig_i(\text{row }i)$ of
$s\le Cnw$ bits and $K=O(1)$ shift tables, every input is moved into $F$ by one of the shifts; and on
such an affine set from which a set $E$ with $2^s\cdot2\abs{E}^2<2^L$ is removed.
\end{theorem}

\begin{outline}
Prepend an input sweep: the input head goes from $0$ to $L-1$ and back while the work heads are
parked, so every row is read in the prefix and the original diagram follows unchanged. The prefix
computes the additive hint along its spine, and the original part runs the code shifted by the shift
table entry of the hint, which it reads from the low $s$ bits of its memory edges. This is the
\emph{syndrome trick} of external-memory coding in the language of tape diagrams. The version with
$K$ shifts adds a selection among the $K$ shifted codes (\cref{thm:pr-select}), and the version with
exceptions adds voting (\cref{thm:pr-vote}).
\end{outline}

\begin{theorem}[Codes failing on few inputs]\label{thm:pr-vote}
\formal{Tape/Statements/PlainRobust.lean}{PlainRobust.noCode_cosmall_of_noCode}
Under the plain conjecture, no code is correct on all inputs outside a set $E$ with
$2\abs{E}^2<2^L$.
\end{theorem}

\begin{outline}
Triple every work walk; the block network is unchanged and the profile capacities triple. Run three
shifted copies of the code, with shifts $0,z_1,z_2$ such that $z_i\oplus z_j\notin E\oplus E$, and
output the majority of the three decoded bits: on every input at most one copy errs.
\end{outline}

\begin{theorem}[Selection among candidates]\label{thm:pr-select}
\formal{Tape/Statements/PlainRobust.lean}{PlainRobust.noCode_hintConst_of_noCode, PlainRobust.noCode_iff_hintConst, PlainRobust.noCode_netCover_of_noCode, PlainRobust.noCode_netHint_of_noCode}
Under the plain conjecture, no hint code with a constant number of hint bits exists; equivalently, the
plain conjecture is unchanged by a constant oracle hint. More generally, no code exists on a set $F$
for which a \emph{network-computable} hint (computed by a causal code within the profile capacities of
a linear-length prefix diagram) together with $O(1)$ shift tables moves every input into $F$; and no
family of codes indexed by a network-computable hint of $O(nw)$ bits solves complete exchange on every
input.
\end{theorem}

\begin{outline}
Run the $K$ candidate codes side by side (clamped and packed). Each node compares the candidates'
would-be output bits at its own output positions and, for every pair of candidates without a recorded
difference, records one position where they differ, with the first candidate's bit; this state
travels along the spine. An appended input sweep fetches the true bit at every recorded position, and
a fresh copy of the diagram runs a candidate that never lost a comparison. A correct candidate never
loses, and every incorrect one loses against it, so the winner is correct. This is discrimination,
not verification: telling two candidates apart takes one fetched bit. For network-computable hints,
concatenate the prefix computation with the diagram and carry the hint along the spine.
\end{outline}

\begin{theorem}[Linear and nearly linear codes]\label{thm:pr-linear}
\formal{Tape/Statements/PlainRobust.lean}{PlainRobust.noCode_linear_of_noCode, PlainRobust.noCode_linearExcept_of_noCode, PlainRobust.noCode_linearPlus_of_noCode, PlainRobust.noCode_linearPlusImage_of_noCode}
Under the plain conjecture, no $\F_2$-linear code solves complete exchange on a set of co-entropy at
most $Cnw$. The same holds for codes that are linear outside a set of edges carrying $O(nw)$ bits, for
codes that are linear plus nonlinear corrections of $O(nw)$ bits in total, and for codes that are
linear plus corrections taking at most $2^{O(nw)}$ joint values.
\end{theorem}

\begin{outline}
For a linear code the set of inputs on which it is correct is closed under XOR, hence a subspace; if it
contains a set of co-entropy $Cnw$ it has codimension at most $Cnw$, so it is the kernel of an additive
syndrome map, and \cref{thm:pr-affine} applies. For the nearly linear codes, the nonlinear values form
an internally computed hint of $O(nw)$ bits: replacing them by constants gives an affine code, a
$2^{-O(nw)}$ fraction of the good inputs share their nonlinear values by pigeonhole, and XOR-ing these
inputs with one of them gives inputs on which the linear part is correct.
\end{outline}

\begin{theorem}[Bounded-degree constraints and decompositions]\label{thm:pr-degree}
\formal{Tape/Statements/PlainRobust.lean}{PlainRobust.noCode_lowDegree_of_noCode, PlainRobust.noCode_quadForm_of_noCode, PlainRobust.noCode_decomp_of_noCode, PlainRobust.noCode_crossQuad_of_noCode}
Under the plain conjecture, no code exists on a set containing the zero set of a Boolean function of
degree at most $D$ (for constant $D$, and the zero set nonempty), in particular the zero set of a
quadratic form. No code exists on a set that, for a partition of the input positions into $K=O(1)$
parts, contains every input supported on a single part; for instance the common zero set of any
family of quadratic forms each of whose monomials joins two different parts.
\end{theorem}

\begin{outline}
The zero set of a function of degree at most $D$ is covered by $2^D$ translates of itself: if a
derivative of order $D'\le D$ is nonzero at $x$, some translate of $x$ by a sum of a sub-list of the
directions lies in the zero set, by a parity argument; so selection among $2^D$ shifted codes applies.
For a decomposition, run $K$ copies of the code, copy $j$ on the input masked outside part $j$, which
lies in $F$, and XOR the outputs: transposition is linear.
\end{outline}

\begin{theorem}[Errors]\label{thm:pr-errors}
\formal{Tape/Statements/PlainRobust.lean}{PlainRobust.noCode_fewErrors_of_noCode, PlainRobust.noCode_sparseErrors_of_noCode, PlainRobust.noCode_sparseEntryErrors_of_noCode, PlainRobust.noCode_fewErrorsCover_of_noCode}
Under the plain conjecture, every code that stays within the profile capacities on every input has
more than $2^{Hnw}$ distinct error vectors $C(x)^{\T}\oplus x$, for every constant $H$. In particular,
no such code gets all but at most $An$ output bits right on every input (transposition with $O(n)$
wrong bits is as hard as exact transposition); the same holds for $O(n)$ wrong entries, and for a
choice among $O(1)$ such codes per input.
\end{theorem}

\begin{outline}
Syndrome correction. Fix a check matrix $\Phi$ with $t=2Hnw+1$ rows that is injective on the possible
error vectors (a random one is). Run the code; every output bit equal to $1$ adds the corresponding
column of $\Phi$ on the spine; an appended input sweep adds the column of every input bit equal to $1$,
so the spine carries the syndrome $\Phi e(x)$; a second run of the code XORs each output bit with the
error bit decoded from the syndrome. The $t$ extra bits are paid for by stationary tapes.
\end{outline}

\begin{theorem}[Column covers]\label{thm:pr-columns}
\formal{Tape/Statements/PlainRobust.lean}{PlainRobust.noCode_colCover_of_noCode, PlainRobust.noCode_rowMapCover_of_noCode, PlainRobust.noCode_colAggregate_of_noCode}
Under the plain conjecture, no code exists on a set $F$ for which $K=O(1)$ locally computable maps
$\Phi_\kappa$ into $F$ exist, each keeping the input columns of its class $\kappa$ (the classes
partitioning the columns); in particular for maps acting on rows. For even $n$, no code exists on a set
containing $\{x:\sum_jv(\text{column }j\text{ of }x)=0\}$ for a fixed $v$ with values in $\F_2^m$.
\end{theorem}

\begin{outline}
Run $K$ copies, copy $\kappa$ on $\Phi_\kappa(x)$ (each node computes the copy's visible input from its
own), and read column $j$ from the copy of its class. For the aggregate constraint, the inputs
$[x_{\mathrm{left}}\mid x_{\mathrm{left}}]$ and $[x_{\mathrm{right}}\mid x_{\mathrm{right}}]$ satisfy it,
since each term appears twice.
\end{outline}

\subsection{Why the general case is hard}

We record, without formal proofs, why these methods stop short of the general implication.
\begin{itemize}[leftmargin=2em]
\item \emph{A general hint-removal theorem must use the code, not only the size of $F$.} A counting
argument over (automaton, shift table) pairs shows that a random set of density $2^{-m}$ has no
streaming covering hint of $O(nw)$ bits.
\item \emph{A hint is cheap in information but not in flow.} An $s$-bit broadcast tells at most $s$
bits about all outputs together, but delivering it to $\Theta(n)$ nodes costs $\Theta(L)$ bits of
capacity in aggregate; edge-removal style arguments do not apply as they stand. The first rung of the
ladder is a one-shot instance of zero error versus constant error, related to the edge-removal problem
and to the continuity of network capacity under dependent sources~\cite{LangbergEffros2021}.
\item \emph{The machine-side view.} Plain would refute (T) if (T) implied oblivious (T), or merely a
linear-time transposer whose block-level trace does not depend on the input. Codes of trace classes are
self-detecting: at every slot the simulation knows whether the run still follows the popular trace. A
hint-removal argument restricted to self-detecting codes would suffice; none is known.
\end{itemize}

\section{Barriers}\label{sec:barriers}

This section explains why the Tape Coding Conjecture is hard to prove and why a proof must use the
shape of tape networks. Its results are not formalized; each gets a complete proof.

\subsection{Multiplication}

By Harvey and van der Hoeven~\cite[Cor.~6.9]{HvdH2025}, a multiplication machine running in time
$M(m)$ yields a transposition machine for $n\times n$ matrices with $b$-bit entries running in time
$O(\ell_0\,\overline M(m))$, where $m\ge n^2b$, $\ell_0=\max(1,\lg^*n-\lg^*b)$ and
$\overline M(m)=m\max_{k\le m}M(k)/k$. For $b=\lceil\log_2n\rceil$ we have $\ell_0=1$, so a
linear-time multiplication machine gives a linear-time transposer. Hence $\neg$(T), and in particular
the Tape Coding Conjecture, implies that integer multiplication cannot be done in linear time on
multitape Turing machines. No superlinear lower bound for off-line multiplication on multitape
machines is known; the classical $\Omega(n\log n)$ bound of Paterson, Fischer and Meyer is for
on-line multiplication~\cite{PFM1974}.

\subsection{General graphs: the circuit barrier}

We use the following elementary model. Let $H$ be an undirected graph with source vertices
$a_0,\dots,a_{n-1}$ and sink vertices $b_0,\dots,b_{n-1}$, every edge of unit capacity. The
\emph{complete bipartite demand} asks to deliver, for every pair $(i,j)$, a message from $a_i$ to $b_j$.
\begin{itemize}[leftmargin=2em]
\item A \emph{routing at rate $\lambda$} is a fractional multicommodity flow that sends $\lambda$ units
from $a_i$ to $b_j$ for every pair, with total flow at most $1$ on every edge.
\item A \emph{causal code with $N$ rounds} fixes an acyclic orientation of $H$ and, in each round, sends
one bit along every edge in its orientation, each bit being a function (which may depend on the round) of the
bits that entered the tail in the same round and, at a source $a_i$, of $a_i$'s messages. Each pair $(i,j)$ is assigned a
round in which $b_j$ must output the message bit $m_{ij}$, a function of the bits entering $b_j$ in
that round; the messages $m_{ij}$ are arbitrary bits, and the code must work for all of them. Its
\emph{rate} is $1/N$ per pair and round.
\end{itemize}
In \cref{prop:rotation-barrier} we use three extensions of this model. A source may be attached at a set of
vertices: at each of them the code may use its messages, also in the output functions, a routing may start at any
of them, and $\dist_H(a_i,b_j)$ is the distance from the nearest. The functions may depend on the round. And edges
may carry symbols of $\kappa$ bits per round, against routings with capacity $\kappa$ per edge; then
\eqref{eq:weak-duality-graph} below reads $\lambda\sum_{i,j}\dist_H(a_i,b_j)\le\kappa\abs{E(H)}$.
The ratio of the rate of a causal code to the largest routing rate is its \emph{coding gain}. A routing
needs one unit of capacity per edge traversed, so for every routing,
\begin{equation}\label{eq:weak-duality-graph}
  \lambda\sum_{i,j}\dist_H(a_i,b_j)\;\le\;\abs{E(H)} .
\end{equation}

\begin{proposition}[Shift circuits are coding gains]\label{prop:circuit-barrier}
\notformalized
Let $C$ be a Boolean circuit with $s$ gates of fan-in at most two (arbitrary gate functions) whose
inputs are $x\in\{0,1\}^n$ and a shift input, such that for every $k<n$ some value of the shift input
makes the outputs equal the cyclic shift $x_{i-k\bmod n}$ at output $i$. Then there is a graph $H$ of
maximum degree at most $4$ with $O(s+n)$ vertices and a causal code with $n$ rounds for the complete
bipartite demand on $H$ whose coding gain is at least $c\,n\log n/(s+n)$, for an absolute constant
$c>0$ and all large $n$.
\end{proposition}

For example, a circuit multiplying two $n$-bit integers, followed by $n$ OR gates, is such a circuit:
for $y=2^k$ the $2n$-bit product $p=x\,2^k$ has $p_j$ and $p_{j+n}$ never both equal to $1$, so the
cyclic shift is $p_j\vee p_{j+n}$.

\begin{informalproof}
\emph{The graph.} Take the circuit as a directed acyclic graph on its gates and on the inputs that feed some gate (at most
$n+2s$ vertices; parallel edges allowed). Replace
every vertex of fan-out $f>2$ by a binary tree of $f-1$ copy vertices, so that in-degrees and
out-degrees are at most $2$; this adds at most $2s+n$ vertices. Let $H$ be the underlying undirected
graph, with $a_i$ the input vertex of $x_i$ and $b_j$ the vertex of output $j$. Its maximum degree is
at most $4$ and it has $O(s+n)$ vertices.

\emph{One code per shift.} Fix $k<n$ and fix the shift input to a value realizing the shift by $k$.
Every vertex then computes a fixed function of its in-edges, so the circuit, read in topological order
with one bit per edge, is a causal code on $H$ (with the circuit's orientation) that delivers $x_i$ from
$a_i$ to $b_{i+k\bmod n}$.

\emph{Time-sharing.} Run $n$ rounds; in round $k$ use the code for shift $k$ with the source $a_i$ sending
the message $m_{i,i+k\bmod n}$ as $x_i$. Every pair $(i,j)$ is served in round $j-i\bmod n$, and every edge
carries one bit per round. This is a causal code with $n$ rounds, of rate $1/n$.

\emph{The flow bound.} In a graph of maximum degree $4$ at most $4^{t+1}$ vertices lie within distance $t$
of a vertex. With $t=\lfloor\tfrac12\log_4n\rfloor-1$, at most $\sqrt n$ sinks lie within distance $t$ of
$a_i$, so
$\sum_{i,j}\dist(a_i,b_j)\ge n(n-\sqrt n)(t+1)\ge\tfrac15n^2\log_4n$ for large $n$. By
\eqref{eq:weak-duality-graph}, every routing has rate $\lambda\le5\abs{E(H)}/(n^2\log_4n)=O((s+n)/(n^2\log n))$.
The coding gain is therefore at least $(1/n)/\lambda=\Omega(n\log n/(s+n))$.
\end{informalproof}

So if, on every bounded-degree graph, complete bipartite exchange had coding gain at most $\gamma(n)$,
then every circuit computing all cyclic shifts, and every multiplication circuit, would have size
$\Omega(n\log n/\gamma(n))$. With $\gamma=o(\log n)$ this is a superlinear circuit lower bound for an
explicit function, which is far beyond current techniques: the best lower bounds for explicit functions
are $3.1n-o(n)$ over the full binary basis~\cite{LiYang2022}. A proof of the Tape Coding Conjecture must
therefore use the structure of tape networks: unless it also yields a superlinear circuit lower bound for an
explicit function, it cannot be a statement about bounded-degree networks in general. The single-shift version of this argument, conditional on the network coding conjecture, is due
to Afshani, Freksen, Kamma and Larsen~\cite{AFKL2019}, and the step from a small circuit to a coding gain
is also observed by Braverman and He~\cite{BravermanHe2026}.

Tape networks are not simple either. The chords of one work tape form two non-crossing families along the
time spine, excursions to the right and to the left, so a $k$-tape network has a book embedding with $2k$ pages,
while bounded-degree expanders have $3$-page book embeddings~\cite{DSW2016}. We expect, but do not prove, that after
spreading each vertex over a constant number of slots every page can be realized by the excursions of one work head,
so that tape networks with a constant number of work tapes contain subdivided bounded-degree expanders, on which cut
bounds would not exclude a code. This is the heuristic reason why we do not expect cut arguments alone to prove the
conjecture.

\subsection{Tape networks: an oblivious separation}

\emph{Programmable rotation} takes $x\in\{0,1\}^L$ and $\rho<L$ (in binary, $\clog L$ bits appended to $x$)
and outputs $x$ rotated by $\rho$. A non-oblivious machine computes it in time $O(L)$: copy $x$ to a work tape,
then count $\rho$ down in binary while moving a head $\rho$ cells (amortized $O(1)$ per decrement), and output the
work tape cyclically from that head. Oblivious machines can rotate in time $O(L\log L)$, by a barrel shifter or
by the oblivious simulation of Pippenger and Fischer~\cite{PippengerFischer1979}.

The \emph{bit-level network} of an oblivious machine $M$ at input length $L+\clog L$ has a vertex for every
step $\tau<T$; a spine edge $\tau\to\tau+1$; and, for every work tape and every pair of consecutive visits
$\sigma<\tau$ of its head to one cell, an edge $\sigma\to\tau$. The source of input bit $x_i$ is attached
at the steps where the input head reads cell $i$, and the sink of output bit $o$ is the step that appends
it. All of this is the same for every input, because $M$ is oblivious.

\begin{proposition}[Rotation codes on oblivious tape networks]\label{prop:rotation-barrier}
\notformalized
Let $M$ be an oblivious machine with $k$ work tapes that computes programmable rotation within $T(L)\le
L\log_2L$ steps. Then its bit-level network carries a causal code with $L$ rounds for the complete
bipartite demand between the input bits and the output bits whose coding gain over routing, with every edge
of capacity $\kappa=\max(\clog{(q+1)},\clog{(g+1)})$, is at least $c\,L\log L/T(L)$ for a constant $c>0$
depending only on $k$, $q$ and $g$, and all large $L$.
\end{proposition}

\begin{informalproof}
Fix $\rho<L$. With the bits of $\rho$ fixed, the run of $M$ is a causal code on its bit-level network: the spine
edge out of step $\tau$ carries the state, and the edge from one visit of a cell to the next carries the
symbol written at the first visit; every step computes its outgoing values from its incoming ones and the
input bit it reads. It delivers $x_i$ to output position $i-\rho\bmod L$. Time-sharing over the $L$ values of $\rho$, as in
\cref{prop:circuit-barrier}, serves every pair (input bit, output bit) once in $L$ rounds.

For the flow bound, every vertex has degree at most $D=2k+2$. Let $\mu_i$ be the number of steps at which
the input head reads cell $i$; then $\sum_i\mu_i\le T$. Within distance $t$ of the steps attaching $x_i$
lie at most $\mu_iD^{t+1}$ steps, hence at most $\mu_iD^{t+1}$ sinks. With
$t=\lfloor\tfrac12\log_DL\rfloor-1$, at most $T\sqrt L\le L^{3/2}\log_2L=o(L^2)$ pairs are within distance
$t$, so $\sum_{i,o}\dist\ge\tfrac12L^2(t+1)$ for large $L$. The network has at most $(k+1)T$ edges, each of
the stated capacity $\kappa$, so a routing has rate $\lambda\le(k+1)T\kappa/(\tfrac12L^2(t+1))=O(T/(L^2\log L))$,
while the code has rate $1/L$. The gain is $\Omega(L\log L/T)$.
\end{informalproof}

So a bit-level analogue of the Tape Coding Conjecture, asserting that the complete bipartite demand between
input and output bits on bit-level networks of oblivious machines running in linear time, at the capacities of
\cref{prop:rotation-barrier}, has coding gain $o(\log L)$, would imply that oblivious machines need time
$\omega(L)$ for programmable rotation, while non-oblivious machines need only $O(L)$. No superlinear
\emph{off-line} separation of oblivious and non-oblivious multitape machines is known: the optimality
claimed by Pippenger and Fischer is for on-line simulation~\cite{PippengerFischer1979}, and the known
superlinear lower bounds specific to oblivious machines are on-line or for one head~\cite{Vitanyi1984}. The block
version used for transposition and the bit version used here are not known to imply each other.
Independently of this calibration, by \cref{thm:oblivious} a proof of the conjecture would show that oblivious
machines cannot transpose in linear time, a lower bound that current techniques do not give. By contrast, Theorem~P and Theorem~M* are unconditional, and
the network coding conjecture, which would have given such bounds directly~\cite{AHJKL2006}, is
false~\cite{BravermanHe2026,ZhangLiLi2026}.

\section{Limits of XOR-frame recursions}\label{sec:frames}

The only technique claimed to beat $O(L\log n)$ for transposition on tapes is an unrefereed
construction of OpenAI~\cite{OpenAI2026mult}, which claims $O(L(\log n)^{1-\kappa})$ for a tiny
$\kappa>0$. Its core is an address-interchange primitive that mixes data by XOR in sheared address
frames, applied recursively. This section abstracts the recursion and shows that it cannot give linear
time: every gadget it could use makes at least as many frame changes as the trivial one, so its cost is at
least $W$. The results are about an
abstract model, so they rule out a class of constructions, not linear-time transposition.

\subsection{The framework}

Fix $m\ge2$. Two address chunks $H$ and $D$ are to be interchanged; write each as a vector of $m$
fields. For a rational $m\times m$ matrix $M$ the \emph{frame} $\Phi_M$ is the address map
$(H,D)\mapsto(H+MD,D)$; frames compose additively, $\Phi_A\Phi_B=\Phi_{A+B}$. Changing the frame of a
stream from $A$ to $B$ costs $\rank(B-A)$ recursive interchanges of narrower chunks.

\begin{definition}[Gadget]\label{def:gadget}
\formal{Frames/Gadget.lean}{Transpose.ShearFrame.Gadget, Transpose.ShearFrame.Gadget.cost}
A \emph{gadget} has $W$ wires; wire $w$ starts in a source frame $S_w$, passes through the frames of its
gates in time order, and ends in a sink frame $T_w$. Gates are XOR operations between wires in a common
frame. The scalar network permutes the wires by a permutation $\rho$ (G1), and
$T_{\rho(w)}=S_w+I$ (G2): every wire ends sheared by the identity. The \emph{cost} $s$ of the gadget is
the sum over wires of the ranks of consecutive frame differences along the wire, from source to sink.
Formally, \lean{Transpose.ShearFrame.Gadget} records $\rho$, the source frames and the frame paths subject to
(G2); the gates and (G1) enter only through the run model of \cref{def:frame-runs}.
\end{definition}

The recursion built on a gadget distributes a stream over the $W$ wires and runs the gadget, each frame
change being a recursive interchange of narrower chunks. Its cost per unit of data satisfies
$F_k\le(s/W)F_{k-1}+O(1)$ at width $m^k$, and every implementation in the template also satisfies
$F_k\ge(s/W)F_{k-1}+c_0$. A gadget with $s<W$ would therefore give linear-time transposition, a gadget
with $s/W=m^\tau$, $0<\tau<1$, would give $O(L(\log n)^\tau)$, and $s\ge W$ excludes linear time within the
template; $s=W$ gives $\Theta(L\log\log n)$. (This reading of the template is informal: it rests on the
preprint's construction~\cite{OpenAI2026mult}, which we use only to motivate the model.)

\subsection{Cost bounds}

\begin{theorem}[Cycle bound and projection bound]\label{thm:cycle}
\formal{Frames/Gadget.lean}{Transpose.ShearFrame.Gadget.invariant_cost_ge, Transpose.ShearFrame.Gadget.cost_ge_of_invariant_family, Transpose.ShearFrame.Gadget.cost_ge_m}
\formal{Frames/GadgetProjection.lean}{Transpose.ShearFrame.PGadget.two_mul_cost_ge}
Every nonempty $\rho$-invariant set of wires has total cost at least $m$, so $s\ge jm$ for $j$ disjoint
such sets, and $s\ge m$ if $W\ge1$. If every wire passes through at least one gate and every gate frame is
idempotent (a projection, as in the preprint), then $2s\ge Wm$.
\end{theorem}

\begin{outline}
Rank is subadditive, so the cost of a wire is at least $\rank(T_w-S_w)$, and over a cycle $C$ of $\rho$,
$\sum_{w\in C}(T_w-S_w)=\abs{C}I$ has rank $m$. For idempotent frames, $\rank(P_V-P_U)\ge\abs{\dim V-\dim
U}$ and $\rank(P_A+P_B)\ge\max(\dim A,\dim B)$; pairing the source term of $w$ with the sink term of
$\rho(w)$ and optimizing gives $2s\ge Wm$.
\end{outline}

The general lower bound uses a run model in which frames are arbitrary and only frame changes are counted.

\begin{definition}[Runs]\label{def:frame-runs}
\formal{Frames/GadgetRuns.lean}{Transpose.ShearFrame.Cfg, Transpose.ShearFrame.Event, Transpose.ShearFrame.moves}
A configuration assigns to each of $W$ wires a frame (from any set $F$) and a vector in $\F_2^W$. An
event is either a gate $i\leftarrow i+j$, allowed only for $i\ne j$ at a common frame, or a move of one
wire to a new frame. A run starts with wire $u$ at frame $S_u$ holding the unit vector $e_u$.
\end{definition}

\begin{theorem}[At least $W$ frame changes]\label{thm:frame-change}
\formal{Frames/GadgetRuns.lean}{Transpose.ShearFrame.moves_ge_W, Transpose.ShearFrame.cost_ge_W, Transpose.ShearFrame.moves_ge_W_matrix}
Let $\mathrm{shift}:F\to F$ have no fixed point. Every valid run that ends with wire $\rho(u)$ at frame
$\mathrm{shift}(S_u)$ holding $e_u$, for every $u$, contains at least $W$ moves. Hence every gadget has cost $s\ge W$ (for matrix frames, $\mathrm{shift}(S)=S+I$): reading its gates as
gate events and its frame changes as moves, (G1) and (G2) say that the run ends with wire $\rho(u)$ at
$S_u+I$ holding $e_u$, and a genuine frame change costs the rank of a nonzero matrix, at least $1$ (this
translation is not formalized; the structure of \cref{def:gadget}, which omits the gates and (G1), only
satisfies $s\ge m$). So the framework cannot give linear time.
\end{theorem}

\begin{outline}
Let $\sigma_f$ be the span of the vectors on the wires at frame $f$, and $N_f$ the span of the $e_u$ with
$\mathrm{shift}(S_u)=f$. The potential $\sum_f\dim(\sigma_f\cap N_f)$ is unchanged by gates (they act
within one frame by transvections), grows by at most one per move, is $0$ initially (no unit vector starts
at its target, since $\mathrm{shift}$ has no fixed point), and is $W$ at the end.
\end{outline}

\begin{theorem}[Mixing across frames does not help]\label{thm:group-algebra}
\formal{Frames/GroupAlgebraRuns.lean}{Transpose.ShearFrame.GroupAlgebra.changes_ge_W, Transpose.ShearFrame.GroupAlgebra.cost_ge_W}
In the group-algebra model, where each wire holds a formal sum of sheared items (a function from an
additive group of frames to $\F_2^W$), gates $i\leftarrow i+j$ are allowed between any wires, and a change
shifts the content of one wire, every run from the configuration in which wire $u$ holds exactly item $u$ at frame $S_u$ to one in which
wire $\rho(u)$ holds item $u$ at frame $S_u+I$ and nothing else ($I\ne0$) makes at least $W$ changes.
\end{theorem}

\begin{outline}
The span of the vectors $w\mapsto(\text{content of }w\text{ at frame }I+S_u)_u$ has dimension $0$ at the
start and $W$ at the end; gates do not increase it and each change increases it by at most one.
\end{outline}

\subsection{Coarsening and integer frames}

\begin{theorem}[Coarsening]\label{thm:coarsening}
\formal{Frames/GadgetCoarse.lean}{Transpose.ShearFrame.vmoves_ge_W, Transpose.ShearFrame.visible_ge_W_of_coarse, Transpose.ShearFrame.det_ne_zero_of_moves_le, Transpose.ShearFrame.rank_sum_of_moves_le}
For every map $\pi:F\to F'$ with $\pi(\mathrm{tgt}(u))\ne\pi(S_u)$ for all $u$, every valid run as in
\cref{thm:frame-change} (with targets $\mathrm{tgt}(u)$) makes at least $W$ moves that change the
$\pi$-class of the moved wire. For rational matrix frames with targets $S_u+I$, a run with at most $W$
moves changes every frame by a matrix of nonzero determinant, and the total rank of its jumps is at least $Wm$ (for a gadget whose wires each change frame once this
means cost at least the baseline $Wm$, through the same unformalized translation as in
\cref{thm:frame-change}).
\end{theorem}

\begin{outline}
The potential of \cref{thm:frame-change} needs only a target frame per item and survives coarsening: a move
that does not change the $\pi$-class of the wire does not change the coarsened run. For single jumps: if a jump $J$ of a run with at most $W$ moves were singular, choose $\ell\ne0$ with
$\ell^{\T}J=0$; the coarsening $M\mapsto\ell^{\T}M$ separates every $S_u$ from $S_u+I$ (they differ by
$\ell^{\T}\ne0$) but hides $J$, so fewer than $W$ moves would be visible. So every jump has rank $m$, and
there are at least $W$ jumps by \cref{thm:frame-change}.
\end{outline}

\begin{theorem}[Mod-$2$ bounds for integer frames]\label{thm:mod2}
\formal{Frames/GadgetMod2.lean}{Transpose.ShearFrame.two_mul_le_rank_sum}
For integer $2\times2$ frames, the total rank of the frame jumps of every valid run that ends with wire
$\rho(u)$ at frame $S_u+I$ holding $e_u$ (as in \cref{thm:frame-change}) is at least $2W$:
gadgets with $2\times2$ integer frames never beat the baseline.
\end{theorem}

\begin{theorem}[Mod-$2$ bounds, every $m$]\label{thm:mod2m}
\formal{Frames/GadgetOneSided.lean}{Transpose.ShearFrame.one_sided_bound}
\formal{Frames/GadgetTwoSided.lean}{Transpose.ShearFrame.two_sided_bound}
For integer $m\times m$ frames, the total rank $s$ (over $\Q$) of the jumps of every valid run that ends with
wire $\rho(u)$ at frame $S_u+I$ holding $e_u$ satisfies
$2(2^m-1)W\le2^ms$ and $2(2^m-1)W\le(2^{m-1}+1)s$. In particular $s\ge2.8W$ for $m=3$ and
$s\ge(4-o(1))W$ as $m\to\infty$.
\end{theorem}

\begin{outline}[Outline of the formal proofs of \cref{thm:mod2,thm:mod2m}]
Average \cref{thm:coarsening} over a family of coarsenings: $M\mapsto\ell^{\T}(M\bmod2)$ for nonzero
$\ell\in\F_2^m$ (one-sided), or $M\mapsto\ell^{\T}(M\bmod2)v$ for pairs with $\ell^{\T}v=1$ (two-sided). Each
coarsening separates every source from its target, so it sees at least $W$ jumps. A jump of rank $r$ over
$\F_2$ is seen by $2^m-2^{m-r}$ of the one-sided coarsenings, and by a bounded number of the two-sided ones,
computed from the left kernel and the fixed vectors of the jump; and the rank over $\F_2$ is at most the
rank over $\Q$ (\leanfile{Frames/RankMod.lean}). For $m=2$ a finite check over $\F_2^{2\times2}$ gives the
constant.
\end{outline}

\subsection{Tight runs}

A valid run that takes every item $u$ from its source frame $S_u$ to a target frame $\mathrm{tgt}(u)\ne S_u$
with at most $W$ frame-changing moves is \emph{tight}; by \cref{thm:coarsening} it makes exactly $W$ such moves.
Tight runs are rigid.

\begin{theorem}[Structure of tight runs]\label{thm:tight}
\formal{Frames/GadgetCargo.lean}{Transpose.ShearFrame.tight_potentials, Transpose.ShearFrame.tight_sigma_split, Transpose.ShearFrame.cargo_rule, Transpose.ShearFrame.cargo_rule_bool}
In a tight run, after any prefix with $t$ frame-changing moves the ``arrived'' potential is $t$ and the ``not yet
departed'' potential is $W-t$, and at every frame the span of the vectors there splits into its arrived and its
not-departed parts. If moreover each item has only two frames (every frame other than $S_u$ is
$\mathrm{tgt}(u)$, as after a two-class coarsening), then every frame-changing move carries a vector lying in the
span of the items whose target is its destination and in the span of the items born at its origin (the cargo
rule). (The formal statements assume finitely many frames; a run uses only finitely many, but this transfer
is not formalized.)
\end{theorem}

\begin{theorem}[Item itineraries]\label{thm:itineraries}
\formal{Frames/GadgetItinerary.lean}{Transpose.ShearFrame.chain_exists_visible, Transpose.ShearFrame.itinerary_condition, Transpose.ShearFrame.chain_itinerary, Transpose.ShearFrame.itinerary_condition_matrix, Transpose.ShearFrame.chain_itinerary_matrix}
In every valid run as in \cref{thm:coarsening}, every item has a chain of frame-changing hops (moves whose cargo
contains the item) from its source to its target. Let $\pi$ be a two-class coarsening with
$\pi(\mathrm{tgt}(u))\ne\pi(S_u)$ for all $u$, and suppose the run makes exactly $W$ moves that change the
$\pi$-class. Then every such move carrying an item starts on the item's birth side ($\pi$ of its origin equals
$\pi(S_u)$), and every item has a chain from its source to its target all of whose $\pi$-visible hops are
admissible in this sense. For frames in $\F_2^{m\times m}$ with targets $S_u+I$ and a coarsening by a functional
$\Lambda$ with $\operatorname{tr}\Lambda=1$ (the class of $M$ is $\langle\Lambda,M\rangle$), admissibility of a hop of $s$
from $f_0$ with jump $J$ and $\langle\Lambda,J\rangle=1$ reads $\langle\Lambda,S_s+f_0\rangle=0$.
\end{theorem}

\begin{theorem}[Every carrying move lies on an admissible chain]\label{thm:carrying}
\formal{Frames/GadgetItineraryC.lean}{Transpose.ShearFrame.carrying_move_on_chain, Transpose.ShearFrame.carrying_move_on_admissible_chain, Transpose.ShearFrame.visible_move_on_admissible_chain}
Every frame-changing move whose cargo contains an item lies on a chain of visible hops of that item from its
source to its target; under tightness, every visible move lies on an admissible chain of some item.
\end{theorem}

\begin{theorem}[Mixing lemma]\label{thm:mixing}
\formal{Frames/GadgetMixing.lean}{Transpose.ShearFrame.mixed_move_forces_extra_hop, Transpose.ShearFrame.mixing_lemma, Transpose.ShearFrame.mixing_lemma_matrix}
Weight frames by a function $d$ with $d(f,f)=0$, $d(f,g)=0\Rightarrow f=g$ and the triangle inequality (for
matrices, $d(A,B)=\rank(B-A)$), and let $d(S_u,\mathrm{tgt}(u))\ge m$ for every item $u$. Suppose a
valid run as in \cref{thm:frame-change} contains a frame-changing move whose cargo contains two distinct
items, and every other frame-changing move carries a single item. Then some item of that cargo has
$d$-weighted hop count at least $m+1$: a mixed move forces an extra hop.
\end{theorem}

\begin{outline}[Outline of the formal proofs of \cref{thm:tight,thm:itineraries,thm:carrying,thm:mixing}]
The vectors of a configuration reached by a valid run are linearly independent, since gates are
transvections; so the spans at distinct frames are independent and their dimensions add up to $W$. The
arrived and not-departed potentials each change by at most one per visible move and sum to at most $W$, so in
a tight run both move in lockstep, which forces the split and the cargo rule. Chains are found by following an
item backwards (and forwards) through the moves whose cargo contains it, using that an item is present at a
frame only if some move carried it there. Admissibility (\cref{thm:itineraries}) is the cargo rule of the coarsened two-class run. For the mixing
lemma, suppose every item of the mixed cargo $x$, moved into frame $g$, has hop count at most $m$. An item
present at a frame after (before) a segment was (will be) present at a frame within $d$-distance of it
bounded by the segment's hop count; with the triangle inequality, no item of the support of $x$ is present at
$g$ before the mixed move, and no later frame-changing move into $g$ carries one. As the later moves out of
$g$ carry unit vectors, every wire at $g$ stays constant on the support of $x$ while some wire at $g$ keeps
coordinate $s_0$; at the end that wire holds $e_{s_0}$, which is not constant on a support of two items.
\end{outline}

Together these say that, within the XOR-frame template, every gadget costs at least $W$ (so linear time is
out of reach), gadgets with integer frames cost substantially more than $W$, and gadgets that attain the bound (tight runs) are rigid. Whether some gadget has $s/W$ substantially below $m$, which would give an improvement
over $O(L\log n)$ much larger than the claimed $\kappa$, is open (\cref{sec:open}).

\section{The construction side}\label{sec:construction}

Our aim at the outset was a linear-time transposer. We did not find one. This section records what the
results above require of any construction, describes the one architecture we know that passes every
test, and summarizes the mechanisms that fail. It contains no new theorems: every claim is either one of
the results above, a citation, or explicitly labelled as a heuristic.

\subsection{Necessary conditions}

A machine witnessing~(T) must have all of the following properties.
\begin{enumerate}[label=(N\arabic*),leftmargin=3em]
\item \emph{Coding.} It cannot only move data: machines that move indivisible items need time
$\Omega(L\log n)$ (Stoß~\cite{Stoss1973} for Turing machines, Aggarwal and Vitter~\cite{AggarwalVitter1988}
in external memory). Cells must hold functions of several entries.
\item \emph{Many passes, at every scale.} Its run has $\Omega(\log n)$ passes (Theorem~P$'$,
\cref{thm:Pprime}), even on a $2^{-O(nb)}$ fraction of the inputs (\cref{thm:Pprime-on}); the grouped form
(\cref{thm:Pgamma}) gives the same at coarser column scales.
\item \emph{Folding motion.} Along the popular trace, the macroscopic work-head trajectories satisfy $(\star)$
(\cref{cor:mstar-T}): some head makes a macroscopic U-turn, and the heads together waste at least one unit of
range (\cref{thm:excess}).
\item \emph{A maximal coding gain.} The block network of its run carries a causal code for complete exchange,
on a set of co-entropy $O(n)$, at the profile capacities (\cref{thm:reduction,thm:coent-linear}); by
\cref{thm:code-gap} this code beats routing by a factor $\Omega(\log n)$, the largest gain that cut bounds
allow. For an oblivious machine the code works on all inputs (\cref{thm:oblivious}); the same then holds on
the bit-level network~\cite[Thm.~6.2]{AHJKL2006}.
\item \emph{Not an XOR-frame recursion.} It cannot be built from the gadgets of \cref{sec:frames}, which cost at
least the trivial gadget.
\end{enumerate}

Condition~(N4) is the heart of the matter. Coding gains over routing in undirected networks do exist: the
multiple-unicast (network coding) conjecture is false~\cite{BravermanHe2026}. But the known growing gains are
of order $(\log n)^{\varepsilon}$, obtained by amplifying a constant gap by tensoring~\cite{BGS2017}, and they
concern sparse demands; tensoring a constant gap cannot produce the maximal gain $\Theta(\log n)$ for dense
demands (heuristically, and as observed by Zhang, Li and Li~\cite{ZhangLiLi2026}). A linear-time transposer
needs a \emph{native} maximal gain for the densest demand, on a network shaped by tapes.

\subsection{An architecture that survives}

We describe, as a design and not as a result, the one architecture that we could not rule out.

\paragraph{Geometric recursion.} The transposer runs $\Theta(\log n)$ levels. Level $\ell$ works on data of
volume $\theta^\ell L$ for a constant $\theta<1$, with $O(1)$ passes of an encoder $E_\ell$ and a decoder
$D_\ell$: it computes its output as $D_\ell(X_\ell,f_{\ell+1}(E_\ell(X_\ell)))$, where $f_{\ell+1}$ is the next
level's problem. The total time is $\sum_\ell O(\theta^\ell L)=O(L)$, and the pass count is $\Theta(\log n)$, as
Theorem~P requires.

\paragraph{Expanding choreography.} The heads of each level are driven along prepared \emph{rulers} that realize
two non-commuting monotone maps of the interval, so that the relation graph between the encoder's input and
output rows is a bounded-degree expander. Bourgain and Yehudayoff's monotone expanders~\cite{BourgainYehudayoff2013}
provide such maps; the rulers can be written in linear time, since their marks are computed with
$O(\log L)$-bit arithmetic. Near-identity generators are needed: integer or dyadic M\"obius pairs generate
discrete groups whose relation graphs are not expanders (a heuristic supported by numerical spectral gaps).
Repeating each choreography a constant number of times should make the macroscopic motion satisfy the
quantitative inequality $(\star)$; this is a cost constraint, not an obstruction (heuristic).

\paragraph{Two coding cores.} The geometry is the easy part. Each level needs a constant-fraction coding saving at
the encoder, so that $\abs{E_\ell(X_\ell)}\le\theta\abs{X_\ell}$, and, because the decoder writes the level's whole
output from $O(1)$-pass views of the next level's output, also at the decoder. In the simplest linear form the
encoder core is the following problem.

\begin{quote}
\emph{Locally encoded index coding with own-row side information.} Let $\Gamma$ be a bounded-degree graph on
$[n]$. A broadcast consists of linear functionals of an $n\times n$ matrix, each supported on the two rows of an
edge of $\Gamma$. Receiver $j$ sees the broadcast and knows row $j$, and must recover column $j$. Can the broadcast
have size $(1-\varepsilon)n^2$ for a fixed $\varepsilon>0$?
\end{quote}

Without the degree bound, $\varepsilon=\tfrac12$ is achievable: broadcast $A+A^{\T}$ above the diagonal. With the
degree bound, routing-type codes save only $O(n^2/\log n)$ on expanders, and hyperfinite (amenable) graphs save
$o(n^2)$ (heuristic arguments; not included here). A positive answer on expanders would be the first
constant-fraction coded saving in a tape-relevant model. A negative answer, together with its decoder analogue
(restricted-view index coding), would rule out every geometric recursion with $O(1)$-pass levels.

\subsection{Mechanisms that fail}

The following classes of constructions fail for reasons proved in this paper.
\begin{itemize}[leftmargin=2em]
\item \emph{Few passes.} Any schedule with $o(\log n)$ passes, with any coding, including constant-pass
choreographies along expanders (Theorems~P and~P$'$).
\item \emph{Regular or one-way motion.} Machines whose macroscopic head motion is stationary or monotone on every
tape (\cref{thm:excess,thm:mstar-adaptive}).
\item \emph{XOR-frame recursions,} including the template of~\cite{OpenAI2026mult}, which cannot reach linear time
(\cref{sec:frames}).
\item \emph{Codes whose linear structure covers the input set.} For a construction aimed at the plain conjecture's
counterexamples, linear codes, codes correct off a small set, and codes that are linear outside $O(nw)$ bits of
capacity reduce to plain codes (\cref{sec:plainrobust}).
\end{itemize}
Many further attempts failed for reasons we could only argue informally: flat schemes that code across a constant
number of whole streams (they save at most a constant factor), window-local recursions whose levels read only
nearby data, static coding certificates on Ramanujan-type graphs that resist conversion into causal codes, and
re-indexings that turn transposition into a cheap map only through a transposition-hard change of basis. We
mention them as warnings, not as results.

\section{Open problems}\label{sec:open}

\begin{enumerate}[label=\arabic*.,leftmargin=2em]
\item \emph{Decide (T).} Either construct a machine transposing $n\times n$ matrices with $\lceil\log_2n\rceil$-bit
entries in time $O(n^2\log n)$, or prove that none exists. A proof is accepted by the frozen checker in the form
described in \cref{sec:model}.

\item \emph{The Tape Coding Conjecture} (\cref{conj:TCC}). A proof refutes~(T) (\cref{thm:reduction}). It suffices to
rule out codes at co-entropy $Cn$ for every $C$, for infinitely many $n$, on diagrams that write their output at the
end (\cref{thm:coent-linear,thm:io}). By \cref{sec:barriers}, a proof must use the shape of tape networks, and its
techniques would give lower bounds for oblivious computation beyond the state of the art. A refutation would be a
family of codes realizing the maximal coding gain for complete exchange on tape networks; it would not by itself give
a machine.

\item \emph{Plain versus robust.} Does the plain conjecture (codes correct on all inputs) imply the robust one? Already
the first rung is open: does ``no code on all inputs'' imply ``no code on half of the inputs''? A positive answer
for codes of trace classes, which are self-detecting, would suffice for~(T) (\cref{sec:plainrobust}).

\item \emph{Oblivious machines.} Is there an oblivious linear-time transposer? The plain conjecture says no
(\cref{thm:oblivious}). We expect that Harvey and van der Hoeven's reduction~\cite{HvdH2025} preserves obliviousness,
so that a negative answer would also rule out oblivious linear-time multiplication; this requires checking that all
control flow and data layout in their construction depend only on the sizes. Does (T) imply oblivious~(T)?

\item \emph{Locally encoded index coding} (\cref{sec:construction}). On bounded-degree expanders, can a broadcast of
two-row functionals of size $(1-\varepsilon)n^2$ let every receiver, knowing its own row, recover its column?

\item \emph{Sharper passes.} Remove the hypothesis $T\le\poly(L)$ from Theorem~P$'$ (for example through the run-length
lemma of Grohe and Schweikardt~\cite{GroheSchweikardt2005}), and extend Theorems~P and~P$'$ to passes that may move
back by a bounded number of cells.

\item \emph{Beyond levels.} Theorem~M* concerns partitions of macroscopic time into intervals. Cuts of other shapes, or
information inequalities that are not cut bounds, are needed to go further; is there a tape-specific non-cut inequality
for oblivious machines?

\item \emph{XOR-frame gadgets.} Is there a gadget (\cref{def:gadget}) with $s/W\le m^{\tau}$ for some $\tau<1$
bounded away from $1$? This would improve the exponent of the sub-$n\log n$ transposition claim substantially, though
not to linear time (\cref{thm:frame-change}). For integer frames the lower bounds of \cref{thm:mod2m} stay below
$4W$, while the trivial gadget costs $mW$; close this gap.
\end{enumerate}


\appendix
\section{Index of Lean declarations}\label{app:index}

Every declaration cited in this paper, grouped by the Lean module that contains it. All of them
depend only on the axioms \texttt{propext}, \texttt{Classical.choice} and \texttt{Quot.sound}
(\cref{app:verify}).

\paragraph{\texttt{Transpose.Frames.Gadget}}
{\small\raggedright \texttt{Transpose.ShearFrame.Gadget}, \texttt{Transpose.ShearFrame.Gadget.cost}, \texttt{Transpose.ShearFrame.Gadget.cost\_ge\_m}, \texttt{Transpose.ShearFrame.Gadget.cost\_ge\_of\_invariant\_family}, \texttt{Transpose.ShearFrame.Gadget.invariant\_cost\_ge}.\par}

\paragraph{\texttt{Transpose.Frames.GadgetCargo}}
{\small\raggedright \texttt{Transpose.ShearFrame.cargo\_rule}, \texttt{Transpose.ShearFrame.cargo\_rule\_bool}, \texttt{Transpose.ShearFrame.tight\_potentials}, \texttt{Transpose.ShearFrame.tight\_sigma\_split}.\par}

\paragraph{\texttt{Transpose.Frames.GadgetCoarse}}
{\small\raggedright \texttt{Transpose.ShearFrame.det\_ne\_zero\_of\_moves\_le}, \texttt{Transpose.ShearFrame.rank\_sum\_of\_moves\_le}, \texttt{Transpose.ShearFrame.visible\_ge\_W\_of\_coarse}, \texttt{Transpose.ShearFrame.vmoves\_ge\_W}.\par}

\paragraph{\texttt{Transpose.Frames.GadgetItinerary}}
{\small\raggedright \texttt{Transpose.ShearFrame.chain\_itinerary}, \texttt{Transpose.ShearFrame.chain\_itinerary\_matrix}, \texttt{Transpose.ShearFrame.itinerary\_condition}, \texttt{Transpose.ShearFrame.itinerary\_condition\_matrix}.\par}

\paragraph{\texttt{Transpose.Frames.GadgetItineraryC}}
{\small\raggedright \texttt{Transpose.ShearFrame.carrying\_move\_on\_admissible\_chain}, \texttt{Transpose.ShearFrame.carrying\_move\_on\_chain}, \texttt{Transpose.ShearFrame.visible\_move\_on\_admissible\_chain}.\par}

\paragraph{\texttt{Transpose.Frames.GadgetMixing}}
{\small\raggedright \texttt{Transpose.ShearFrame.mixed\_move\_forces\_extra\_hop}, \texttt{Transpose.ShearFrame.mixing\_lemma}, \texttt{Transpose.ShearFrame.mixing\_lemma\_matrix}.\par}

\paragraph{\texttt{Transpose.Frames.GadgetMod2}}
{\small\raggedright \texttt{Transpose.ShearFrame.two\_mul\_le\_rank\_sum}.\par}

\paragraph{\texttt{Transpose.Frames.GadgetOneSided}}
{\small\raggedright \texttt{Transpose.ShearFrame.one\_sided\_bound}.\par}

\paragraph{\texttt{Transpose.Frames.GadgetProjection}}
{\small\raggedright \texttt{Transpose.ShearFrame.PGadget.two\_mul\_cost\_ge}.\par}

\paragraph{\texttt{Transpose.Frames.GadgetRuns}}
{\small\raggedright \texttt{Transpose.ShearFrame.Cfg}, \texttt{Transpose.ShearFrame.Event}, \texttt{Transpose.ShearFrame.cost\_ge\_W}, \texttt{Transpose.ShearFrame.moves}, \texttt{Transpose.ShearFrame.moves\_ge\_W}, \texttt{Transpose.ShearFrame.moves\_ge\_W\_matrix}.\par}

\paragraph{\texttt{Transpose.Frames.GadgetTwoSided}}
{\small\raggedright \texttt{Transpose.ShearFrame.two\_sided\_bound}.\par}

\paragraph{\texttt{Transpose.Frames.GroupAlgebraRuns}}
{\small\raggedright \texttt{Transpose.ShearFrame.GroupAlgebra.changes\_ge\_W}, \texttt{Transpose.ShearFrame.GroupAlgebra.cost\_ge\_W}.\par}

\paragraph{\texttt{Transpose.Statement}}
{\small\raggedright \texttt{Transpose.ComputesWithin}, \texttt{Transpose.LinearTimeTransposable}, \texttt{Transpose.TransposesWithin}, \texttt{Transpose.encode}, \texttt{Transpose.entryBits}.\par}

\paragraph{\texttt{Transpose.Tape.BufTM}}
{\small\raggedright \texttt{Impl.bufTM}.\par}

\paragraph{\texttt{Transpose.Tape.Defs}}
{\small\raggedright \texttt{AdmissibleRes}, \texttt{HaltsAt}, \texttt{IsChord}, \texttt{IsLastVisit}, \texttt{ObliviousOn}, \texttt{Transposes}, \texttt{blockOf}, \texttt{chordsInto}, \texttt{curBlock}, \texttt{cutPairs}, \texttt{cutSet}, \texttt{entryCells}, \texttt{entryPoints}, \texttt{matrixInputs}, \texttt{numSlots}, \texttt{pairRes}, \texttt{residentPair}, \texttt{run}, \texttt{slotOf}, \texttt{snkCols}, \texttt{srcRows}, \texttt{touchRes}, \texttt{tpos}, \texttt{transAt}.\par}

\paragraph{\texttt{Transpose.Tape.DefsCE}}
{\small\raggedright \texttt{IsWalk}, \texttt{NetCode}, \texttt{NetCode.edgeVal}, \texttt{NetCode.inVals}, \texttt{NetCode.outVal}, \texttt{NoCode\_TM}, \texttt{NoCode\_TM\_coent}, \texttt{NoCode\_TM\_hintBits}, \texttt{ObliviousLinearTimeTransposable}, \texttt{RobustCE\_TM}, \texttt{RobustCE\_TM\_dist}, \texttt{RobustNoCode\_TM}, \texttt{Routing}, \texttt{TSDiagram}, \texttt{TSDiagram.CapProfile}, \texttt{TSDiagram.RoutesAt}, \texttt{TSDiagram.Solves}, \texttt{TSDiagram.Valid}, \texttt{TSDiagram.blockCap}, \texttt{TSDiagram.blockSize}, \texttt{TSDiagram.chordEdges}, \texttt{TSDiagram.edges}, \texttt{TSDiagram.memEdges}, \texttt{TSDiagram.nodes}, \texttt{TSDiagram.outNode}, \texttt{TSDiagram.profileCap}, \texttt{TSDiagram.rowDist}, \texttt{TSDiagram.rowsAt}, \texttt{TSDiagram.visible}, \texttt{ball}, \texttt{netDist}, \texttt{runDiagram}.\par}

\paragraph{\texttt{Transpose.Tape.DefsM}}
{\small\raggedright \texttt{ConsecIn}, \texttt{CutIneq}, \texttt{FinUnionIcc}, \texttt{PiecewiseConstOn}, \texttt{blockTraj}, \texttt{cellBits}, \texttt{crossPairs}, \texttt{horizon}, \texttt{levelCross}, \texttt{levelExcess}, \texttt{levelSet}, \texttt{rescaledTraj}, \texttt{slotLen}.\par}

\paragraph{\texttt{Transpose.Tape.DefsP}}
{\small\raggedright \texttt{AllMono}, \texttt{MonoOn}, \texttt{PassPartition}, \texttt{addrCount}, \texttt{addresses}, \texttt{addrsBelow}, \texttt{childNode}, \texttt{combSet}, \texttt{forest}, \texttt{forestNode}, \texttt{greedyBound}, \texttt{greedyPasses}, \texttt{readSet}, \texttt{rootNode}, \texttt{sSet}.\par}

\paragraph{\texttt{Transpose.Tape.DefsPP}}
{\small\raggedright \texttt{nodeCode}, \texttt{sideCode}, \texttt{structCount}.\par}

\paragraph{\texttt{Transpose.Tape.DefsTrace}}
{\small\raggedright \texttt{SameTrace}, \texttt{TraceClass}, \texttt{blockTrace}, \texttt{rescaledTrajOf}, \texttt{traceConst}, \texttt{traceCount}.\par}

\paragraph{\texttt{Transpose.Tape.Statements.Adaptive}}
{\small\raggedright \texttt{adaptive\_hyps\_of\_linearTime}, \texttt{block\_locality}, \texttt{block\_locality\_offsets\_counterexample}, \texttt{cut\_bound\_trace}, \texttt{cut\_bound\_trace\_bits}, \texttt{cut\_bound\_trace\_bits\_on}, \texttt{cut\_bound\_trace\_on}, \texttt{exists\_popular\_traceClass}, \texttt{forest\_agree}, \texttt{forest\_comb\_agree}, \texttt{forest\_comb\_agree\_on}, \texttt{isoperimetry\_trace}, \texttt{isoperimetry\_trace\_on}, \texttt{kraft\_group}, \texttt{kraft\_group\_on}, \texttt{linearTimeTransposable\_popular\_limits}, \texttt{locality\_anchored}, \texttt{locality\_entry}, \texttt{mstar\_discrete\_class\_on}, \texttt{no\_stationary\_limits\_class\_on}, \texttt{rescaledTrajOf\_close}, \texttt{tendstoUniformlyOn\_of\_mem}, \texttt{theoremMStar\_class\_nonuniform}, \texttt{theoremMStar\_excess\_nonuniform}, \texttt{theoremMStar\_excess\_robust}, \texttt{theoremMStar\_level\_adaptive}, \texttt{theoremMStar\_level\_class}, \texttt{theoremMStar\_level\_class\_on}, \texttt{theoremMStar\_level\_nonuniform}, \texttt{theoremMStar\_level\_robust}, \texttt{theoremPprime}, \texttt{theoremPprime\_bits}, \texttt{theoremPprime\_div}, \texttt{theoremPprime\_gamma}, \texttt{theoremPprime\_greedy}, \texttt{theoremPprime\_grouped}, \texttt{theoremPprime\_grouped\_bits}, \texttt{theoremPprime\_grouped\_on}, \texttt{theoremPprime\_grouped\_on\_bits}, \texttt{theoremPprime\_hyps\_of\_oblivious}, \texttt{theoremPprime\_hyps\_tr2dd}, \texttt{theoremPprime\_linearTime}, \texttt{theoremPprime\_no\_five\_pass}, \texttt{theoremPprime\_on}, \texttt{theoremPprime\_on\_bits}, \texttt{theoremPprime\_on\_div}, \texttt{theoremPprime\_on\_greedy}, \texttt{theoremPprime\_on\_pow}, \texttt{theoremPprime\_on\_robust}, \texttt{theoremPprime\_on\_robust\_gamma}, \texttt{theoremPprime\_pow}, \texttt{theoremPprime\_ratio}, \texttt{traceBits\_le}, \texttt{traceClass\_of\_mem}, \texttt{traceClass\_of\_oblivious}, \texttt{traceClass\_of\_obliviousOn}.\par}

\paragraph{\texttt{Transpose.Tape.Statements.Coding}}
{\small\raggedright \texttt{blockTraj\_limit\_lipschitz}, \texttt{bufTM\_run}, \texttt{card\_chordEdges\_le}, \texttt{card\_le\_of\_solves}, \texttt{ce\_TM\_dist\_iff\_noCode}, \texttt{ce\_TM\_dist\_of\_robustCE\_TM\_dist}, \texttt{ce\_TM\_iff\_noCode}, \texttt{ce\_TM\_of\_robustCE\_TM}, \texttt{code\_gap}, \texttt{code\_of\_hint}, \texttt{codes\_coentLinear\_of\_linearTimeTransposable}, \texttt{codes\_of\_fastOn}, \texttt{codes\_of\_fastOn\_bits}, \texttt{codes\_of\_linearTimeTransposable}, \texttt{codes\_of\_obliviousLinearTimeTransposable}, \texttt{coentLinear\_io\_iff\_exists\_coent\_io}, \texttt{coentLinear\_io\_of\_robustNoCode\_io}, \texttt{coentLinear\_of\_robustCE\_TM\_lin}, \texttt{coentLinear\_of\_robustNoCode}, \texttt{coent\_io\_nw\_of\_robustNoCode\_io}, \texttt{coent\_io\_of\_coent}, \texttt{coent\_mono}, \texttt{coent\_of\_hintBits}, \texttt{cutIneq\_of\_solves}, \texttt{cutset\_bits\_bitwise}, \texttt{cutset\_bound}, \texttt{cutset\_bound\_bitwise}, \texttt{exists\_popular\_traceClass\_on}, \texttt{exists\_shift\_cover}, \texttt{exists\_shift\_iff\_unbounded}, \texttt{exists\_subseq\_blockTraj}, \texttt{flow\_rate\_lt}, \texttt{frozen\_V2}, \texttt{hintBits\_of\_coent}, \texttt{hintCodes\_linear\_of\_linearTimeTransposable}, \texttt{hintCodes\_of\_linearTimeTransposable}, \texttt{hintLinear\_iff\_coentLinear}, \texttt{hint\_of\_code}, \texttt{isoperimetry\_code}, \texttt{linearTime\_fails\_of\_coentLinear}, \texttt{linearTime\_fails\_of\_robustNoCode}, \texttt{linearTime\_fails\_of\_robustNoCode\_logBits}, \texttt{mstar\_of\_not\_robustNoCode}, \texttt{mstar\_of\_robustNoCode\_counterexamples}, \texttt{noCode\_bits\_iff}, \texttt{noCode\_iff\_coent\_zero}, \texttt{noCode\_iff\_forall\_cap}, \texttt{noCode\_io\_of\_robustNoCode\_io}, \texttt{noCode\_of\_coentLinear}, \texttt{noCode\_of\_noCode\_hint}, \texttt{noCode\_of\_robustNoCode}, \texttt{no\_fast\_transposer\_of\_coent}, \texttt{no\_fast\_transposer\_of\_coentLinear}, \texttt{no\_fast\_transposer\_of\_robustNoCode}, \texttt{no\_fast\_transposer\_of\_robustNoCode\_bits}, \texttt{not\_linearTimeTransposable\_of\_coent}, \texttt{not\_linearTimeTransposable\_of\_coentLinear}, \texttt{not\_linearTimeTransposable\_of\_coentLinear\_end}, \texttt{not\_linearTimeTransposable\_of\_coentLinear\_end\_io}, \texttt{not\_linearTimeTransposable\_of\_coentLinear\_io}, \texttt{not\_linearTimeTransposable\_of\_coent\_io}, \texttt{not\_linearTimeTransposable\_of\_coent\_nw}, \texttt{not\_linearTimeTransposable\_of\_noCode\_hint}, \texttt{not\_linearTimeTransposable\_of\_noCode\_hintLinear}, \texttt{not\_linearTimeTransposable\_of\_robustCE\_TM}, \texttt{not\_linearTimeTransposable\_of\_robustCE\_TM\_dist}, \texttt{not\_linearTimeTransposable\_of\_robustCE\_TM\_io}, \texttt{not\_linearTimeTransposable\_of\_robustCE\_TM\_lin}, \texttt{not\_linearTimeTransposable\_of\_robustNoCode}, \texttt{not\_linearTimeTransposable\_of\_robustNoCode\_io}, \texttt{not\_linearTimeTransposable\_of\_robustNoCode\_logBits}, \texttt{not\_obliviousLinearTimeTransposable\_of\_ce\_TM}, \texttt{not\_obliviousLinearTimeTransposable\_of\_ce\_TM\_dist}, \texttt{not\_obliviousLinearTimeTransposable\_of\_noCode}, \texttt{not\_obliviousLinearTimeTransposable\_of\_noCode\_io}, \texttt{not\_robustCE\_TM\_of\_linearTimeTransposable}, \texttt{numSlots\_le\_of}, \texttt{pairRes\_congr}, \texttt{plain\_tfae}, \texttt{robustCE\_TM\_att\_iff}, \texttt{robustCE\_TM\_att\_iff\_robustCE\_TM}, \texttt{robustCE\_TM\_dist\_iff\_robustNoCode}, \texttt{robustCE\_TM\_gamma\_iff}, \texttt{robustCE\_TM\_gamma\_per\_iff}, \texttt{robustCE\_TM\_iff\_robustNoCode}, \texttt{robustCE\_TM\_lin\_iff}, \texttt{robustNoCode\_bits\_iff}, \texttt{robustNoCode\_iff\_attached}, \texttt{robustNoCode\_iff\_coent}, \texttt{robustNoCode\_iff\_forall\_cap}, \texttt{robustNoCode\_iff\_hint}, \texttt{robustNoCode\_io\_of\_robustNoCode}, \texttt{robustNoCode\_logBits\_iff}, \texttt{robustNoCode\_of\_gainBound}, \texttt{robustNoCode\_with\_iff}, \texttt{robustNoCode\_with\_k\_iff}, \texttt{robustNoCode\_with\_of\_robustNoCode}, \texttt{robust\_io\_tfae}, \texttt{robust\_tfae}, \texttt{rowDist\_sum\_ge}, \texttt{solves\_mono}, \texttt{solves\_profileCap\_addTapes}, \texttt{solves\_shiftCode}, \texttt{sum\_cap\_le}, \texttt{sum\_cross\_blockCap\_le}, \texttt{theoremMStar\_diagram}, \texttt{theoremMStar\_diagram\_profile}, \texttt{theoremMStar\_level\_diagram}, \texttt{theoremMStar\_level\_walk}, \texttt{theoremMStar\_walk}, \texttt{traceClass\_code}, \texttt{traceClass\_coent\_lt}, \texttt{volume\_gap}, \texttt{weak\_duality}.\par}

\paragraph{\texttt{Transpose.Tape.Statements.Core}}
{\small\raggedright \texttt{addrCount\_mul}, \texttt{ae\_finite\_levelSet}, \texttt{card\_addresses\_le}, \texttt{card\_le\_entryData}, \texttt{charging\_bound}, \texttt{computesWithin\_iff}, \texttt{copyTM\_obliviousOn}, \texttt{copyTM\_transposes}, \texttt{countable\_localExtr\_values}, \texttt{coverage\_grouped}, \texttt{coverage\_per\_address}, \texttt{crossChords\_split}, \texttt{cutPairs\_card\_mul\_le}, \texttt{cut\_bound}, \texttt{cut\_bound\_bits}, \texttt{cut\_bound\_on}, \texttt{cut\_bound\_on\_coentropy}, \texttt{cut\_bound\_pairs}, \texttt{cut\_bound\_sum}, \texttt{cut\_bound\_sum\_bits}, \texttt{cut\_bound\_sum\_on}, \texttt{entryCells\_slotSteps\_card}, \texttt{entryPoints\_slotSteps\_card}, \texttt{entryTimes\_card\_le}, \texttt{entryTimes\_charged}, \texttt{excess\_add\_range\_le}, \texttt{exists\_slot\_block\_eq}, \texttt{exists\_subseq\_tendstoUniformlyOn}, \texttt{forestNode\_disjoint}, \texttt{forest\_entryCells\_card\_le}, \texttt{forest\_entryCells\_subset}, \texttt{forest\_entryPoints\_card\_le}, \texttt{forest\_entryPoints\_subset}, \texttt{forest\_info}, \texttt{forest\_subset}, \texttt{greedyPasses\_indep}, \texttt{greedyPasses\_le}, \texttt{greedy\_passPartition}, \texttt{haltsAt\_unique}, \texttt{isoperimetry}, \texttt{isoperimetry\_res}, \texttt{isoperimetry\_res\_on}, \texttt{isoperimetry\_touch}, \texttt{levelChords\_limsup\_le}, \texttt{levelCross\_le\_levelExcess}, \texttt{levelSet\_encard\_majorant}, \texttt{lintegral\_encard\_levelSet\_le}, \texttt{lintegral\_levelExcess\_eq\_zero\_of\_monotone}, \texttt{locality}, \texttt{locality\_core}, \texttt{locality\_function}, \texttt{locality\_interval}, \texttt{locality\_oblivious}, \texttt{matrixInputs\_eq}, \texttt{no\_stationary\_limits}, \texttt{not\_all\_monotone\_of\_star}, \texttt{numSlots\_le\_horizon}, \texttt{one\_le\_excess\_of\_star}, \texttt{outRev\_reverse\_eq}, \texttt{pairChange\_spacing}, \texttt{pairRes\_V2}, \texttt{pairRes\_admissible}, \texttt{pairRes\_begin\_unique}, \texttt{pairRes\_chord\_source\_unique}, \texttt{pairRes\_chord\_target\_unique}, \texttt{pairRes\_end\_unique}, \texttt{pairRes\_subset}, \texttt{passPartition\_iff}, \texttt{passPartition\_indep}, \texttt{pass\_fresh}, \texttt{range\_add\_one\_le\_of\_star}, \texttt{read\_chain}, \texttt{rescaledTraj\_limit\_lipschitz}, \texttt{slot\_block\_step}, \texttt{slot\_cut\_bound}, \texttt{slot\_cut\_bound\_partition}, \texttt{srcRows\_slotSteps\_card\_le}, \texttt{stepN\_init\_eq\_some\_iff}, \texttt{sum\_cutPairs\_ge}, \texttt{theoremMStar\_conclusion\_can\_fail}, \texttt{theoremMStar\_excess}, \texttt{theoremMStar\_level}, \texttt{theoremMStar\_level\_sub}, \texttt{theoremMStar\_level\_subseq}, \texttt{theoremP}, \texttt{theoremP\_bits}, \texttt{theoremP\_bits\_add\_one}, \texttt{theoremP\_div}, \texttt{theoremP\_greedy}, \texttt{theoremP\_grouped}, \texttt{theoremP\_grouped\_bits}, \texttt{theoremP\_grouped\_on}, \texttt{theoremP\_grouped\_on\_bits}, \texttt{theoremP\_hyps\_tr2}, \texttt{theoremP\_hyps\_tr2k}, \texttt{theoremP\_on}, \texttt{theoremP\_on\_bits}, \texttt{theoremP\_on\_div}, \texttt{theoremP\_on\_pow}, \texttt{theoremP\_on\_robust}, \texttt{theoremP\_on\_robust\_gamma}, \texttt{theoremP\_pow}, \texttt{touchRes\_admissible}, \texttt{touchRes\_begin\_unique}, \texttt{touchRes\_chord\_source\_unique}, \texttt{touchRes\_chord\_target\_unique}, \texttt{touchRes\_end\_unique}, \texttt{tpos\_spec}, \texttt{tr2\_cutPairs}, \texttt{tr2\_obliviousOn}, \texttt{tr2\_transposes}, \texttt{tr2k\_obliviousOn}, \texttt{tr2k\_transposes}, \texttt{transposes\_of\_transposesWithin}, \texttt{uncovered\_cols\_le}, \texttt{uncovered\_cols\_le\_len}, \texttt{uncovered\_cols\_le\_slots}, \texttt{uncovered\_cols\_le\_slots\_len}, \texttt{unitSteps\_inTraj}, \texttt{unitSteps\_workTraj}, \texttt{work\_eq\_blank}, \texttt{work\_eq\_written}.\par}

\paragraph{\texttt{Transpose.Tape.Statements.PlainRobust}}
{\small\raggedright \texttt{PlainRobust.noCode\_additiveCover\_of\_noCode}, \texttt{PlainRobust.noCode\_additive\_of\_noCode}, \texttt{PlainRobust.noCode\_affine\_exc\_of\_noCode}, \texttt{PlainRobust.noCode\_affine\_of\_noCode}, \texttt{PlainRobust.noCode\_colAggregate\_of\_noCode}, \texttt{PlainRobust.noCode\_colCover\_of\_noCode}, \texttt{PlainRobust.noCode\_cosmall\_of\_noCode}, \texttt{PlainRobust.noCode\_crossQuad\_of\_noCode}, \texttt{PlainRobust.noCode\_decomp\_of\_noCode}, \texttt{PlainRobust.noCode\_fewErrorsCover\_of\_noCode}, \texttt{PlainRobust.noCode\_fewErrors\_of\_noCode}, \texttt{PlainRobust.noCode\_hintConst\_of\_noCode}, \texttt{PlainRobust.noCode\_iff\_hintConst}, \texttt{PlainRobust.noCode\_linearExcept\_of\_noCode}, \texttt{PlainRobust.noCode\_linearPlusImage\_of\_noCode}, \texttt{PlainRobust.noCode\_linearPlus\_of\_noCode}, \texttt{PlainRobust.noCode\_linear\_of\_noCode}, \texttt{PlainRobust.noCode\_lowDegree\_of\_noCode}, \texttt{PlainRobust.noCode\_netCover\_of\_noCode}, \texttt{PlainRobust.noCode\_netHint\_of\_noCode}, \texttt{PlainRobust.noCode\_of\_affine\_subset}, \texttt{PlainRobust.noCode\_quadForm\_of\_noCode}, \texttt{PlainRobust.noCode\_rowMapCover\_of\_noCode}, \texttt{PlainRobust.noCode\_sparseEntryErrors\_of\_noCode}, \texttt{PlainRobust.noCode\_sparseErrors\_of\_noCode}.\par}

\paragraph{\texttt{Transpose.Tape.Witness}}
{\small\raggedright \texttt{Witness.tr2}, \texttt{Witness.tr2k}.\par}

\paragraph{\texttt{Transpose.Tape.WitnessDD}}
{\small\raggedright \texttt{Witness.tr2dd}.\par}

\section{Building and verifying the formalization}\label{app:verify}

The Lean project is the directory \leanfile{lean/} next to this paper. It uses Lean and Mathlib, both at version
v4.34.1 (\leanfile{lean-toolchain}, \leanfile{lakefile.toml}, \leanfile{lake-manifest.json}).

\paragraph{Build.} With \texttt{elan} installed, run in \leanfile{lean/}:
\begin{center}
\texttt{lake exe cache get \&\& lake build}
\end{center}
The first command downloads the compiled Mathlib; the second compiles the library (86 modules).

\paragraph{The frozen files.} \leanfile{Transpose/Statement.lean} (the statement (T)) and \leanfile{Verify.lean} (the
checker of a future \leanfile{Transpose/Answer.lean}) are byte-identical to the files fixed before this work began;
\texttt{shasum -a 256 -c FROZEN.sha256} checks them. \leanfile{Verify.lean} imports
\leanfile{Transpose/Answer.lean}, which does not exist because (T) is not decided; running it therefore fails with a
missing-module error.

\paragraph{Axioms.} The file \leanfile{PaperAxioms.lean}, generated from the \LaTeX{} sources of this paper, runs
\texttt{\#print axioms} on every declaration cited here (\cref{app:index}):
\begin{center}
\texttt{lake env lean PaperAxioms.lean}
\end{center}
Every declaration depends at most on \texttt{propext}, \texttt{Classical.choice} and \texttt{Quot.sound}. No Lean file uses
\texttt{sorry}, \texttt{admit} or \texttt{native\_decide}, or declares an \texttt{axiom} (the words occur only in
comments of the frozen files).

\paragraph{Where things are.} The public statements are in four files. They are proved from theorems in \texttt{Impl}
namespaces, usually of the same name, which the other files of the library prove:
\begin{itemize}[leftmargin=2em]
\item \leanfile{Tape/Statements/Core.lean}: \cref{sec:model,sec:locality,sec:blocks}, Theorem~P (\cref{sec:passes}) and
Theorem~M* for oblivious machines (\cref{sec:mstar});
\item \leanfile{Tape/Statements/Adaptive.lean}: Theorem~P$'$, trace classes, and Theorem~M* for adaptive, robust and
non-uniform machines (\cref{sec:passes,sec:traces,sec:mstar});
\item \leanfile{Tape/Statements/Coding.lean}: codes, Theorem~M* for walks and for code-carrying diagrams, the Tape
Coding Conjecture and its consequences
(\cref{sec:codes,sec:conjecture,sec:plainrobust});
\item \leanfile{Tape/Statements/PlainRobust.lean}: the partial results of \cref{sec:plainrobust}.
\end{itemize}
The results on XOR-frame recursions (\cref{sec:frames}) are in \leanfile{Frames/}. Definitions are in
\leanfile{Tape/Defs.lean}, \leanfile{Tape/DefsP.lean}, \leanfile{Tape/DefsPP.lean}, \leanfile{Tape/DefsTrace.lean},
\leanfile{Tape/DefsM.lean} and \leanfile{Tape/DefsCE.lean}.


\bibliographystyle{plainurl}
\bibliography{references}

\end{document}
